% Lesson 55 — Reading: Texts on using ICTs as recreational resources at home or in the community — Anglais Tle A — Cours signé OnBuch+
% Programme officiel MINESEC. Compiler avec : tectonic EN55-reading-texts-on-using-icts-as-recreational-resources-a.tex
\newcommand{\DOCMATIERE}{Anglais}
\newcommand{\DOCNIVEAU}{Tle A}
\newcommand{\DOCMODULE}{Chapter 5 : Media and Communication}
\newcommand{\DOCLECON}{EN55}
\newcommand{\DOCTITRE}{Lesson 55 — Reading: Texts on using ICTs as recreational resources at home or in the community}
% OnBuch+ — préambule « cours signés » ENRICHI (compilé avec tectonic / XeLaTeX)
% Système visuel « Pop » d'OnBuch+ : fond crème (jamais blanc), bordures encre
% épaisses, ombres « dures » décalées, titres Archivo Black en capitales.
% Version MATHÉMATIQUES : graphes (pgfplots/tikz), boîtes pédagogiques
% (définition, propriété, méthode, attention, le savais-tu, exercice résolu,
% exercice, TP, bilan), page de garde et signature OnBuch+ ancrée en bas.
%
% Macros attendues dans main.tex AVANT \input{preamble} :
%   \DOCMATIERE  \DOCNIVEAU  \DOCMODULE  \DOCLECON (numéro)  \DOCTITRE
\documentclass[11pt]{article}

\usepackage{fontspec}
\usepackage[english,french]{babel} % français = langue principale ; l'anglais via \eng{..} / textebox
\usepackage[a4paper,margin=2cm,top=3.4cm,bottom=2.8cm,headheight=1.4cm,footskip=1.3cm]{geometry}
\usepackage{amsmath,amssymb,amsfonts}
\usepackage{mathtools} % \xmapsto, \coloneqq... souvent utilisés par les modèles
\usepackage{cancel} % simplifications barrées (bilans de forces)
% \abs{z} (module, valeur absolue) et \norm{u} (norme), utilisés par les modèles
\providecommand{\abs}{}\let\abs\relax\DeclarePairedDelimiter{\abs}{\lvert}{\rvert}
\providecommand{\norm}{}\let\norm\relax\DeclarePairedDelimiter{\norm}{\lVert}{\rVert}
\usepackage[table]{xcolor} % [table] : fournit aussi \rowcolors (alternance de lignes)
\usepackage{pagecolor}
\usepackage{tikz}
\usetikzlibrary{arrows.meta,positioning,calc,shapes.geometric,shapes.misc,shapes.arrows,decorations.pathmorphing,decorations.pathreplacing,decorations.markings,patterns,shadows,mindmap,trees,fit,backgrounds,matrix,intersections,angles,quotes}
\usepackage{pgfplots}
\usepackage[version=4]{mhchem} % équations nucléaires \ce{^{235}_{92}U}
\usepackage[european,siunitx]{circuitikz} % schémas électriques
\pgfplotsset{compat=1.18}
\usepgfplotslibrary{fillbetween,groupplots} % aires entre courbes (calcul intégral)
\usepackage{enumitem}
\usepackage{fancyhdr}
% [most] charge toutes les bibliothèques tcolorbox (listings, minted, raster,
% documentation...) et gonfle inutilement la mémoire TeX sur un document long
% (~300 boîtes) : on ne charge que ce qui est réellement utilisé.
\usepackage[skins,breakable]{tcolorbox}
\usepackage{graphicx}
\usepackage{colortbl}
\usepackage{booktabs}
\usepackage{tabularx}
\usepackage{diagbox} % case d'en-tête diagonale des tableaux à double entrée
% Colonnes extensibles centrée / gauche / droite (C, L, R), souvent utilisées par les modèles
\newcolumntype{C}{>{\centering\arraybackslash}X}
\newcolumntype{L}{>{\raggedright\arraybackslash}X}
\newcolumntype{R}{>{\raggedleft\arraybackslash}X}
\usepackage{array}
\usepackage{multicol}
\usepackage{multirow}
\usepackage{siunitx}
% parse-numbers=false : accepte aussi \SI{2,5 \times 10^{-3}}{...}
\sisetup{locale=FR,output-decimal-marker={,},group-minimum-digits=4,parse-numbers=false}
\DeclareSIUnit\ohm{\ensuremath{\symup{\Omega}}} % U+2126 absent de la police
\DeclareSIUnit\division{div} % réglages d'oscilloscope (ms/div, V/div)
\DeclareSIUnit\tr{tr} % tours (vitesse de rotation, ex. \SI{1500}{\tr\per\minute})
\DeclareSIUnit\molar{M} % concentration molaire (ex. \SI{3}{\molar}, \SI{1}{\micro\molar})
\DeclareSIUnit\an{an} % année (le modèle confond parfois avec \year, un compteur TeX) — ex. \SI{5}{\kilo\meter\per\an}
\DeclareSIUnit\FCFA{FCFA} % franc CFA (ex. \SI{500}{\FCFA\per\kilo\gram})
\DeclareSIUnit\degreeBrix{\textdegree Brix} % teneur en sucre d'un sirop/jus (réfractométrie)
\DeclareSIUnit\month{mois} % le modèle utilise parfois \month, un compteur TeX, comme unité de durée
\usepackage{qrcode}
% \degree : commande fréquemment utilisée par les modèles pour un angle ou une
% température en dehors de \SI{}{} (non fournie par les paquets chargés).
\providecommand{\degree}{^{\circ}}
% \qed (fin de démonstration), utilisé par les modèles sans amsthm
\providecommand{\qed}{\hfill$\square$}
% \barre : double barre des valeurs interdites dans un tableau de variations (tkz-tab)
\providecommand{\barre}{\ensuremath{\|}}
% environnement proof (amsthm non chargé) utilisé par les modèles
\ifdefined\proof\else\newenvironment{proof}[1][Démonstration]{\par\noindent\textit{#1.}\ }{\qed\par}\fi
\usepackage{lastpage}
\usepackage{hyperref}
\hypersetup{hidelinks}

% ============================================================
% Palette « Pop » OnBuch+ — mêmes hex que la classe P (plus_ui.dart)
% ============================================================
\definecolor{poporange}{HTML}{F59321}
\definecolor{popdark}{HTML}{E07A0C}
\definecolor{popcream}{HTML}{FFF6E8}
\definecolor{popink}{HTML}{1C1714}
\definecolor{poppurple}{HTML}{7A5AE0}
\definecolor{popgreen}{HTML}{2E9E6A}
\definecolor{poppink}{HTML}{E0457B}
\definecolor{popblue}{HTML}{2F7DE1}
\definecolor{popgold}{HTML}{C98A12}
\definecolor{popmuted}{HTML}{8A7F76}
\definecolor{popline}{HTML}{EADFD2}
\definecolor{popgray}{HTML}{8A7F76}
\colorlet{popgrey}{popgray}
% Alias tolérants (le modèle nomme parfois les couleurs différemment)
\colorlet{popwhite}{white}
\colorlet{popblack}{popink}
\colorlet{popred}{poppink}
\colorlet{popyellow}{popgold}
% Teintes claires (fonds de tableaux/encadrés) — le modèle nomme parfois
% les couleurs avec un suffixe L (ex. popblueL) sans les avoir définies.
\colorlet{poporangeL}{poporange!20}
\colorlet{popdarkL}{popdark!20}
\colorlet{poppurpleL}{poppurple!20}
\colorlet{popgreenL}{popgreen!20}
\colorlet{poppinkL}{poppink!20}
\colorlet{popblueL}{popblue!20}
\colorlet{popgoldL}{popgold!20}
\colorlet{popredL}{popred!20}
\colorlet{popgrayL}{popgray!20}
% Teintes claires pour fonds de tableaux / zones
\colorlet{popblueL}{popblue!10!white}
\colorlet{poporangeL}{poporange!14!white}
\colorlet{popgreenL}{popgreen!12!white}
\colorlet{poppurpleL}{poppurple!12!white}
\colorlet{poppinkL}{poppink!12!white}
\colorlet{popgoldL}{popgold!14!white}

\pagecolor{popcream}

% ============================================================
% Typographie
% ============================================================
\defaultfontfeatures{Ligatures=TeX}
\setmainfont{PlusJakartaSans-Medium.ttf}[Path=../../../fonts/,BoldFont=PlusJakartaSans-Bold.ttf]
% Maths en sans-serif (Fira Math) avec chiffres et lettres droites de Plus
% Jakarta, pour que les formules se fondent dans le texte.
\usepackage[math-style=ISO,bold-style=ISO]{unicode-math}
\setmathfont{FiraMath-Regular.otf}[Path=../../../fonts/]
\setmathfont{PlusJakartaSans-Medium.ttf}[Path=../../../fonts/,range={up/num,up/{Latin,latin},\mathup}]
\usepackage{icomma} % virgule décimale sans espace en mode math
% Caractères mathématiques Unicode absents des polices de texte (Plus Jakarta,
% Archivo Black) : rendus par leur équivalent mathématique, en texte comme en
% formule — sinon ils s'affichent en carré vide (ex. « Arithmétique dans ℤ »).
\usepackage{newunicodechar}
\newunicodechar{ℝ}{\ensuremath{\mathbb{R}}}
\newunicodechar{ℤ}{\ensuremath{\mathbb{Z}}}
\newunicodechar{ℂ}{\ensuremath{\mathbb{C}}}
\newunicodechar{ℕ}{\ensuremath{\mathbb{N}}}
\newunicodechar{ℚ}{\ensuremath{\mathbb{Q}}}
\newunicodechar{𝔻}{\ensuremath{\mathbb{D}}}
\newunicodechar{α}{\ensuremath{\alpha}}
\newunicodechar{β}{\ensuremath{\beta}}
\newunicodechar{γ}{\ensuremath{\gamma}}
\newunicodechar{δ}{\ensuremath{\delta}}
\newunicodechar{ε}{\ensuremath{\varepsilon}}
\newunicodechar{θ}{\ensuremath{\theta}}
\newunicodechar{λ}{\ensuremath{\lambda}}
\newunicodechar{μ}{\ensuremath{\mu}}
\newunicodechar{σ}{\ensuremath{\sigma}}
\newunicodechar{τ}{\ensuremath{\tau}}
\newunicodechar{φ}{\ensuremath{\varphi}}
\newunicodechar{ψ}{\ensuremath{\psi}}
\newunicodechar{ω}{\ensuremath{\omega}}
\newunicodechar{Σ}{\ensuremath{\Sigma}}
% majuscules produites par \MakeUppercase dans les titres (α-aminés -> Α)
\newunicodechar{Α}{\ensuremath{\alpha}}
\newunicodechar{Β}{\ensuremath{\beta}}
\newunicodechar{Γ}{\ensuremath{\Gamma}}
\newunicodechar{Δ}{\ensuremath{\Delta}}
\newunicodechar{Ε}{\ensuremath{\varepsilon}}
\newunicodechar{Θ}{\ensuremath{\Theta}}
\newunicodechar{Λ}{\ensuremath{\Lambda}}
\newunicodechar{Μ}{\ensuremath{\mu}}
\newunicodechar{Τ}{\ensuremath{\tau}}
\newunicodechar{Φ}{\ensuremath{\Phi}}
\newunicodechar{Ψ}{\ensuremath{\Psi}}
\newunicodechar{Ω}{\ensuremath{\Omega}}
\newunicodechar{↦}{\ensuremath{\mapsto}}
\newunicodechar{∈}{\ensuremath{\in}}
\newunicodechar{∉}{\ensuremath{\notin}}
\newunicodechar{∧}{\ensuremath{\wedge}}
\newunicodechar{⇔}{\ensuremath{\Leftrightarrow}}
\newunicodechar{⟺}{\ensuremath{\Longleftrightarrow}}
\newunicodechar{⇒}{\ensuremath{\Rightarrow}}
\newunicodechar{⟹}{\ensuremath{\Longrightarrow}}
\newunicodechar{∘}{\ensuremath{\circ}}
\newunicodechar{∪}{\ensuremath{\cup}}
\newunicodechar{∩}{\ensuremath{\cap}}
\newunicodechar{≡}{\ensuremath{\equiv}}
\newunicodechar{⊂}{\ensuremath{\subset}}
\newunicodechar{⊥}{\ensuremath{\perp}}
\newunicodechar{∃}{\ensuremath{\exists}}
\newunicodechar{∀}{\ensuremath{\forall}}
\newunicodechar{∛}{\ensuremath{\sqrt[3]{\,}}}
\newunicodechar{∜}{\ensuremath{\sqrt[4]{\,}}}
\newunicodechar{⃗}{\ensuremath{{}^{\to}}}
\newunicodechar{ⁿ}{\ifmmode^{n}\else\textsuperscript{\ensuremath{n}}\fi}
\newunicodechar{ˣ}{\ifmmode^{x}\else\textsuperscript{\ensuremath{x}}\fi}
\newunicodechar{ᵇ}{\ifmmode^{b}\else\textsuperscript{\ensuremath{b}}\fi}
\newunicodechar{ʳ}{\ifmmode^{r}\else\textsuperscript{\ensuremath{r}}\fi}
\newunicodechar{ᵘ}{\ifmmode^{u}\else\textsuperscript{\ensuremath{u}}\fi}
\newunicodechar{ʸ}{\ifmmode^{y}\else\textsuperscript{\ensuremath{y}}\fi}
\newunicodechar{ᵃ}{\ifmmode^{a}\else\textsuperscript{\ensuremath{a}}\fi}
\newunicodechar{ᵉ}{\ifmmode^{e}\else\textsuperscript{\ensuremath{e}}\fi}
\newunicodechar{ᵏ}{\ifmmode^{k}\else\textsuperscript{\ensuremath{k}}\fi}
\newunicodechar{ᵖ}{\ifmmode^{p}\else\textsuperscript{\ensuremath{p}}\fi}
\newunicodechar{ᴺ}{\ifmmode^{N}\else\textsuperscript{\ensuremath{N}}\fi}
\newunicodechar{ᶜ}{\ifmmode^{c}\else\textsuperscript{\ensuremath{c}}\fi}
\newunicodechar{ʰ}{\ifmmode^{h}\else\textsuperscript{\ensuremath{h}}\fi}
\newunicodechar{ᵠ}{\ifmmode^{q}\else\textsuperscript{\ensuremath{q}}\fi}
\newunicodechar{ₙ}{\ifmmode_{n}\else\textsubscript{\ensuremath{n}}\fi}
\newunicodechar{ₐ}{\ifmmode_{a}\else\textsubscript{\ensuremath{a}}\fi}
\newunicodechar{ᵢ}{\ifmmode_{i}\else\textsubscript{\ensuremath{i}}\fi}
\newunicodechar{ᵣ}{\ifmmode_{r}\else\textsubscript{\ensuremath{r}}\fi}
\newunicodechar{ₖ}{\ifmmode_{k}\else\textsubscript{\ensuremath{k}}\fi}
\newunicodechar{ᵦ}{\ifmmode_{\beta}\else\textsubscript{\ensuremath{\beta}}\fi}
\newunicodechar{ₓ}{\ifmmode_{x}\else\textsubscript{\ensuremath{x}}\fi}
\newfontfamily\popdisplay{ArchivoBlack-Regular.ttf}[Path=../../../fonts/]
\newfontfamily\poplogo{SpaceGrotesk-Bold.ttf}[Path=../../../fonts/]
\newfontfamily\popsemi{PlusJakartaSans-SemiBold.ttf}[Path=../../../fonts/]
\newfontfamily\popxbold{PlusJakartaSans-ExtraBold.ttf}[Path=../../../fonts/]
\newfontfamily\popxboldls{PlusJakartaSans-ExtraBold.ttf}[Path=../../../fonts/,LetterSpace=3.0]

\setlist[enumerate]{leftmargin=1.7em,itemsep=5pt,topsep=4pt}
\setlist[itemize]{leftmargin=1.7em,itemsep=4pt,topsep=4pt,label={\color{poporange}\textbullet}}
\setlength{\parindent}{0pt}
\setlength{\parskip}{5pt}
\renewcommand{\arraystretch}{1.25}

% pgfplots : style Pop par défaut
\pgfplotsset{
  every axis/.append style={axis line style={popink,line width=1pt},tick style={popink},
    grid style={popline,line width=0.5pt},label style={font=\small},tick label style={font=\footnotesize}},
  popcurve/.style={poporange,line width=1.8pt,no markers,smooth},
}
% Même style disponible dans un \draw TikZ simple (hors pgfplots)
\tikzset{popcurve/.style={poporange,line width=1.8pt}}
\tikzset{popgrid/.style={popline,very thin}}
\colorlet{popcurve}{poporange} % le nom du style est parfois utilisé comme couleur

% ============================================================
% Pill / squiggle
% ============================================================
\newcommand{\poppill}[3]{%
  \tikz[baseline=-0.55ex]\node[draw=popink,line width=1.2pt,rounded corners=2.6pt,%
    fill=#2,inner xsep=6.5pt,inner ysep=3pt]{%
    \popxboldls\fontsize{7}{7}\selectfont\color{#3}\MakeUppercase{#1}};%
}
\newcommand{\popsquiggle}[1]{%
  \tikz[baseline=-0.4ex]{
    \draw[#1,line width=2.6pt,line cap=round]
      plot [smooth] coordinates {(0,0.09) (0.34,0.01) (0.68,0.21) (1.02,0.01) (1.36,0.10)};
  }%
}
% Mot-clé surligné (marqueur orange sous le texte)
\newcommand{\cle}[1]{{\popxbold\color{popdark}#1}}

% ============================================================
% En-tête / pied
% ============================================================
\pagestyle{fancy}
\fancyhf{}
\renewcommand{\headrulewidth}{0pt}
\renewcommand{\footrulewidth}{0pt}
\lhead{\raisebox{-0.15ex}{\poplogo\fontsize{13}{13}\selectfont\color{popink}OnBuch\color{poporange}+}}
\rhead{\poppill{\DOCMATIERE\ \textbullet\ \DOCNIVEAU\ \textbullet\ Leçon \DOCLECON}{white}{popink}}
\lfoot{\popxbold\fontsize{7.6}{7.6}\selectfont\color{popmuted}\MakeUppercase{Cours signé OnBuch+}\ \textbullet\ {\popsemi\DOCTITRE}}
\rfoot{\tikz[baseline=-0.6ex]\node[rounded corners=8.5pt,fill=popink,minimum height=17pt,minimum width=17pt,inner xsep=5pt,inner sep=0]{\popxbold\fontsize{8}{8}\selectfont\color{popcream}\thepage\,/\,\pageref*{LastPage}};}

% ============================================================
% Titres
% ============================================================
\newcommand{\chaptitre}[2]{%
  \begin{tcolorbox}[enhanced,colback=poporange,colframe=popink,boxrule=2.6pt,arc=11pt,
      drop shadow=popink,left=15pt,right=15pt,top=12pt,bottom=11pt,before skip=3pt,after skip=18pt]
    \poppill{#2}{popink}{popcream}\par\vspace{9pt}
    {\raggedright\hyphenpenalty=10000\exhyphenpenalty=10000\popdisplay\fontsize{22}{24}\selectfont\color{popcream}\MakeUppercase{#1}\par}
  \end{tcolorbox}
}
% Section numérotée : pastille orange avec le numéro + titre Archivo
\newcounter{popsec}
\newcommand{\coursec}[1]{%
  \stepcounter{popsec}\setcounter{popsub}{0}%
  \par\vspace{16pt}\noindent
  \begin{minipage}{\linewidth}
  \tikz[baseline=-0.7ex]\node[draw=popink,line width=1.6pt,fill=poporange,rounded corners=4pt,minimum size=22pt,inner sep=0]{\popdisplay\fontsize{12}{12}\selectfont\color{popcream}\thepopsec};%
  \hspace{8pt}{\popdisplay\fontsize{15}{17}\selectfont\color{popink}\MakeUppercase{#1}}\par
  \vspace{4pt}\noindent{\color{poporange}\rule{36pt}{3.6pt}}
  \end{minipage}\par\nopagebreak\vspace{8pt}
}
\newcounter{popsub}
\newcommand{\courssub}[1]{%
  \stepcounter{popsub}%
  \par\vspace{9pt}\noindent{\popxbold\fontsize{11.5}{13}\selectfont\color{popink}{\color{poporange}\thepopsec.\thepopsub}\hspace{6pt}#1}\par\nopagebreak\vspace{4pt}
}

% ============================================================
% Boîtes pédagogiques — même recette : fond blanc, bordure épaisse colorée,
% ombre dure encre, badge plein. Titre optionnel [#1].
% ============================================================
\tcbset{popbase/.style={breakable,enhanced,colback=white,boxrule=2.3pt,arc=9pt,
  shadow={3pt}{-3pt}{0pt}{popink},left=14pt,right=14pt,top=11pt,bottom=11pt,before skip=11pt,after skip=15pt,
  fontupper=\popsemi\fontsize{10.6}{14.5}\selectfont\color{popink},
  fontlower=\popsemi\fontsize{10.6}{14.5}\selectfont\color{popink}}}
% Coche dessinée (le glyphe ✓ n'existe pas dans les polices du document)
\newcommand{\popcheck}[1]{\tikz[baseline=-0.5ex]\draw[#1,line width=1.6pt,line cap=round,line join=round](-0.13,0.0)--(-0.04,-0.09)--(0.13,0.1);}
\providecommand{\coche}{\popcheck{popgreen}}
\providecommand{\resultat}[1]{\par\smallskip\noindent\fcolorbox{popgreen}{popgreenL}{\parbox{\dimexpr\linewidth-2\fboxsep-2\fboxrule\relax}{\textbf{Résultat.} #1}}\par\smallskip}
\newcommand{\popboxhead}[3]{\poppill{#1}{#2}{white}\ifx\relax#3\relax\else\hspace{6pt}{\popxbold\fontsize{10.6}{12}\selectfont\color{popink}#3}\fi\par\vspace{7pt}}

\newtcolorbox{aretenir}[1][]{popbase,colframe=popgreen,before upper={\popboxhead{À retenir}{popgreen}{#1}}}
% Texte en anglais (document de lecture, dialogue, extrait) : cadre
% d'étude. Un titre optionnel (source, auteur) s'affiche dans l'en-tête.
\newtcolorbox{textebox}[1][]{popbase,colframe=popblue,before upper={\popboxhead{Text}{popblue}{#1}\begin{otherlanguage}{english}},after upper={\end{otherlanguage}}}
% Marques ✓ / ✗ (forme correcte / fautive) souvent utilisées par les modèles
\providecommand{\cross}{\textcolor{poppink}{\ensuremath{\times}}}
\providecommand{\xmark}{\cross}\providecommand{\crossmark}{\cross}\providecommand{\cmark}{\ensuremath{\checkmark}}
% Ligne de réponse / justification à compléter (inventée par les modèles)
\providecommand{\justification}[1]{\par\noindent\textit{Justification~:} \dotfill\par}
% \textipa (package tipa non chargé) : la police ne couvre pas l'API -> simple passage
\providecommand{\textipa}[1]{#1}
% Soulignement (ulem non chargé) : \uline, \ul
\providecommand{\uline}[1]{\underline{#1}}\providecommand{\ul}[1]{\underline{#1}}
% \glossaire{\item ...} inventée par les modèles : liste de glossaire
\providecommand{\glossaire}[1]{\begin{itemize}#1\end{itemize}}
% \alert (beamer) inventée par les modèles : texte mis en évidence
\providecommand{\alert}[1]{\textbf{\textcolor{poppink}{#1}}}
% variantes ASCII d'ordinaux écrites en macro
\providecommand{\iere}{\textsuperscript{re}}\providecommand{\ieme}{\textsuperscript{e}}\providecommand{\ere}{\textsuperscript{re}}\providecommand{\eme}{\textsuperscript{e}}\providecommand{\er}{\textsuperscript{er}}\providecommand{\ie}{\textsuperscript{e}}
% Ordinaux français écrits en macro par les modèles : 1\ière, 3\ième
\newcommand{\ière}{\textsuperscript{re}}\newcommand{\ième}{\textsuperscript{e}}
% \exemple (inventée par les modèles) : étiquette « Exemple : »
\providecommand{\exemple}{\textbf{Exemple~:}\ }
% \square utilisable aussi hors mode math (cases à cocher)
\AtBeginDocument{\let\squareMath\square\renewcommand{\square}{\ensuremath{\squareMath}}}
% Mot ou phrase en anglais à l'intérieur d'un paragraphe français
\newcommand{\eng}[1]{\ifmmode\text{\textit{#1}}\else\begingroup\selectlanguage{english}\itshape #1\endgroup\fi}
\let\esp\eng % alias toléré
% flèches d'intonation inventées par les modèles
\providecommand{\textsearrow}{}\renewcommand{\textsearrow}{\ensuremath{\searrow}}\providecommand{\textnearrow}{}\renewcommand{\textnearrow}{\ensuremath{\nearrow}}
% \fr{..} : mot français (inventé par les modèles)
\providecommand{\fr}[1]{\textit{#1}}
\newtcolorbox{exemplebox}[1][]{popbase,colframe=poporange,before upper={\popboxhead{Exemple}{poporange}{#1}}}
\newtcolorbox{definition}[1][]{popbase,colframe=popblue,before upper={\popboxhead{Définition}{popblue}{#1}}}
\newtcolorbox{propriete}[1][]{popbase,colframe=poppurple,before upper={\popboxhead{Propriété}{poppurple}{#1}}}
\newtcolorbox{methode}[1][]{popbase,colframe=popgold,before upper={\popboxhead{Méthode}{popgold}{#1}}}
\newtcolorbox{attention}[1][]{popbase,colframe=poppink,before upper={\popboxhead{Attention}{poppink}{#1}}}
\newtcolorbox{experience}[1][]{popbase,colframe=popblue,colback=popblueL,before upper={\popboxhead{Expérience}{popblue}{#1}}}
% « Le savais-tu ? » : carte violette pleine (contexte, culture, Cameroun)
\newtcolorbox{savaistu}[1][]{popbase,colback=poppurple,colframe=popink,
  fontupper=\popsemi\fontsize{10.4}{14.2}\selectfont\color{white},
  before upper={\poppill{Le savais-tu\,?}{white}{poppurple}\ifx\relax#1\relax\else\hspace{6pt}{\popxbold\color{white}#1}\fi\par\vspace{7pt}}}
% Exercice résolu : énoncé en haut, \tcblower puis la solution sur fond crème
\newcounter{popexo}\newcounter{popexr}
\newtcolorbox{exoresolu}[1][]{popbase,colframe=poporange,colbacklower=popcream,
  segmentation style={popink,line width=1.2pt,dashed},
  before upper={\stepcounter{popexr}\popboxhead{Exercice résolu \thepopexr}{poporange}{#1}},
  before lower={\poppill{Solution}{popink}{popcream}\par\vspace{6pt}}}
% Exercice (non résolu en place) — la correction est dans la partie Corrigés
\newtcolorbox{exercice}[1][]{popbase,colframe=popink,
  before upper={\stepcounter{popexo}\popboxhead{Exercice \thepopexo}{popink}{#1}}}
% Corrigé
\newtcolorbox{corrige}[1][]{popbase,colframe=popgreen,colback=popgreenL,
  before upper={\popboxhead{Corrigé}{popgreen}{#1}}}
% Objectifs / prérequis / bilan
\newtcolorbox{objectifs}{popbase,colframe=popink,colback=poporangeL,
  before upper={\popboxhead{Objectifs de la leçon}{popink}{}}}
\newtcolorbox{prerequis}{popbase,colframe=popink,colback=popblueL,
  before upper={\popboxhead{Prérequis}{popblue}{}}}
\newtcolorbox{bilan}{popbase,colframe=popink,colback=white,boxrule=2.8pt,
  before upper={\popboxhead{Fiche bilan}{poporange}{L'essentiel à connaître pour l'examen}}}
% Figure encadrée avec légende
\newtcolorbox{popfigure}[1][]{enhanced,breakable=false,colback=white,colframe=popline,boxrule=1.4pt,arc=8pt,
  left=10pt,right=10pt,top=10pt,bottom=8pt,before skip=10pt,after skip=12pt,halign=center,
  lower separated=false,halign lower=center,
  fontlower=\popsemi\fontsize{9}{11.5}\selectfont\color{popmuted},
  #1}
\newcommand{\legende}[1]{\par\vspace{6pt}{\popsemi\fontsize{9}{11.5}\selectfont\color{popmuted}#1}}

% ============================================================
% Signature OnBuch+
% ============================================================
\newcommand{\poplogoinline}[1]{{\poplogo\fontsize{#1}{#1}\selectfont\color{popink}OnBuch\color{poporange}+}}
\newcommand{\signatureonbuch}{%
  \begin{tcolorbox}[enhanced,colback=popink,colframe=popink,boxrule=2.2pt,arc=10pt,
      drop shadow=poporange,left=14pt,right=14pt,top=10pt,bottom=10pt,before skip=0pt,after skip=0pt]
    \tikz[baseline=-0.7ex]\node[circle,fill=popgreen,draw=popcream,line width=1.3pt,minimum size=22pt,inner sep=0]{\popcheck{white}};%
    \hspace{9pt}\begin{minipage}[c]{0.62\linewidth}
      {\popxbold\fontsize{12}{13}\selectfont\color{popcream}Signé\ \textbullet\ {\poplogo OnBuch\color{poporange}+}}\par\vspace{2pt}
      {\raggedright\popsemi\fontsize{8}{10.5}\selectfont\color{popline}\MakeUppercase{Équipe pédagogique OnBuch+}\\\MakeUppercase{Conforme au programme officiel MINESEC}\par}
    \end{minipage}\hfill
    {\poplogo\fontsize{22}{22}\selectfont\color{popcream}OnBuch\color{poporange}+}
  \end{tcolorbox}%
}

% ============================================================
% Page de garde
% #1 titre · #2 matière · #3 niveau · #4 module · #5 n° leçon · #6 durée ·
% #7 description / objectifs (texte ou itemize)
% ============================================================
\newcommand{\pagegarde}[7]{%
  \thispagestyle{empty}
  \newgeometry{a4paper,margin=1.7cm,top=1.6cm,bottom=1.4cm}%
  \begin{tikzpicture}[remember picture,overlay]
    \fill[poporange!18] ([xshift=-3.2cm,yshift=-3.6cm]current page.north east) circle (4.6cm);
    \fill[poppurple!14] ([xshift=2.4cm,yshift=7.6cm]current page.south west) circle (3.8cm);
  \end{tikzpicture}%
  {\poplogo\fontsize{30}{30}\selectfont\color{popink}OnBuch\color{poporange}+}\hfill
  \raisebox{0.6ex}{\poppill{Cours signé \textbullet\ Premium}{popink}{popcream}}\par
  \vspace{1pt}\popsquiggle{poppurple}\par
  \vspace{8pt}
  \begin{tcolorbox}[enhanced,colback=poporange,colframe=popink,boxrule=2.8pt,arc=13pt,
      drop shadow=popink,left=18pt,right=18pt,top=12pt,bottom=13pt,after skip=10pt]
    \poppill{#2 \textbullet\ #3}{popink}{popcream}\hspace{5pt}\poppill{#4}{white}{popink}\par\vspace{8pt}
    {\popdisplay\fontsize{14}{14}\selectfont\color{popink}LEÇON #5}\par\vspace{4pt}
    {\raggedright\hyphenpenalty=10000\exhyphenpenalty=10000\popdisplay\fontsize{25}{27}\selectfont\color{popcream}\MakeUppercase{#1}\par}
    \vspace{8pt}
    {\popxbold\fontsize{9.5}{9.5}\selectfont\color{popink}\MakeUppercase{Durée conseillée : #6}}
  \end{tcolorbox}
  \begin{tcolorbox}[enhanced,colback=white,colframe=popink,boxrule=2.2pt,arc=10pt,
      drop shadow=popink,left=16pt,right=16pt,top=10pt,bottom=10pt,after skip=8pt]
    \poppill{Dans cette leçon}{poppurple}{white}\par\vspace{5pt}
    \popsemi\fontsize{9.3}{12.6}\selectfont\color{popink}#7
  \end{tcolorbox}
  \noindent\begin{minipage}[t]{0.487\linewidth}
    \begin{tcolorbox}[enhanced,colback=popgreen,colframe=popink,boxrule=2.2pt,arc=10pt,
        drop shadow=popink,left=11pt,right=11pt,top=10pt,bottom=10pt,height=3.1cm,valign=center]
      \tcbox[colback=white,colframe=popink,boxrule=1.3pt,arc=5pt,boxsep=0pt,nobeforeafter,
        left=3pt,right=3pt,top=3pt,bottom=3pt,valign=center]{\qrcode[height=1.9cm]{https://play.google.com/store/apps/details?id=com.luvvix.onbuch&hl=fr}}%
      \hspace{8pt}\begin{minipage}[c]{0.5\linewidth}
        {\popdisplay\fontsize{11}{12.5}\selectfont\color{white}\MakeUppercase{Télécharge OnBuch}}\par\vspace{4pt}
        {\raggedright\popsemi\fontsize{8.4}{11}\selectfont\color{white}Gratuit \textbullet\ Google Play\\Tuteur IA, annales, concours, résultats\par}
      \end{minipage}
    \end{tcolorbox}
  \end{minipage}\hfill
  \begin{minipage}[t]{0.487\linewidth}
    \begin{tcolorbox}[enhanced,colback=poppurple,colframe=popink,boxrule=2.2pt,arc=10pt,
        drop shadow=popink,left=11pt,right=11pt,top=10pt,bottom=10pt,height=3.1cm,valign=center]
      \tikz[baseline=-0.7ex]\node[circle,fill=white,draw=popink,line width=1.3pt,minimum size=1.6cm,inner sep=0]{\popdisplay\fontsize{24}{24}\selectfont\color{poppurple}+};%
      \hspace{8pt}\begin{minipage}[c]{0.6\linewidth}
        {\popdisplay\fontsize{11}{12.5}\selectfont\color{white}\MakeUppercase{Passe à OnBuch+}}\par\vspace{4pt}
        {\raggedright\popsemi\fontsize{8.4}{11}\selectfont\color{white}Tous les cours signés, Tuteur IA illimité, fiches PDF — onglet OnBuch+ de l'app\par}
      \end{minipage}
    \end{tcolorbox}
  \end{minipage}
  \vspace{8pt}
  \popcontenu
  \vfill
  \signatureonbuch
  \restoregeometry
  \newpage
}

% Rangée « Ce cours contient » (page de garde)
\newcommand{\poptile}[3]{% #1 couleur #2 gros texte #3 légende
  \begin{tcolorbox}[enhanced,colback=#1,colframe=popink,boxrule=2pt,arc=9pt,shadow={2.5pt}{-2.5pt}{0pt}{popink},
      left=6pt,right=6pt,top=9pt,bottom=9pt,halign=center,valign=center,height=1.75cm,width=\linewidth,nobeforeafter]
    {\popdisplay\fontsize{12.5}{13}\selectfont\color{white}\MakeUppercase{#2}}\par\vspace{3pt}
    {\centering\popsemi\fontsize{7.8}{9.5}\selectfont\color{white}#3\par}
  \end{tcolorbox}}
\newcommand{\popcontenu}{%
  \par\noindent{\popxboldls\fontsize{7.5}{7.5}\selectfont\color{popmuted}CE COURS SIGNÉ CONTIENT}\par\vspace{7pt}
  \noindent\begin{minipage}[t]{0.235\linewidth}\poptile{popblue}{Cours}{complet, illustré, pas à pas}\end{minipage}\hfill
  \begin{minipage}[t]{0.235\linewidth}\poptile{poporange}{Méthodes}{et exercices résolus}\end{minipage}\hfill
  \begin{minipage}[t]{0.235\linewidth}\poptile{poppink}{Exercices}{type examen + corrigés}\end{minipage}\hfill
  \begin{minipage}[t]{0.235\linewidth}\poptile{popgreen}{Fiche bilan}{l'essentiel à retenir}\end{minipage}\par}

% Sommaire
\newtcolorbox{sommaire}{enhanced,colback=white,colframe=popink,boxrule=2.2pt,arc=10pt,
  drop shadow=popink,left=18pt,right=18pt,top=13pt,bottom=13pt,before skip=6pt,after skip=20pt,
  before upper={\poppill{Sommaire}{poppurple}{white}\par\vspace{9pt}\popsemi\fontsize{10.6}{14.5}\selectfont\color{popink}}}

% Signature de fin de cours — poussée tout en bas de la dernière page
\newcommand{\findecours}{%
  \par\vfill\nopagebreak
  \begin{center}\popsquiggle{poporange}\end{center}\vspace{4pt}
  \signatureonbuch
}
\AtBeginDocument{\def\llbracket{\mathopen{[\mkern-3.5mu[}}\def\rrbracket{\mathclose{]\mkern-3.5mu]}}} % intervalles d'entiers (glyphe ⟦ absent de la police)
% Glyphes absents de Fira Math : redéfinis à partir de glyphes disponibles
\AtBeginDocument{%
  \def\setminus{\mathbin{\backslash}}\let\smallsetminus\setminus
  \def\vdots{\mathord{\vbox{\baselineskip3.2pt\lineskiplimit0pt\kern4pt\hbox{.}\hbox{.}\hbox{.}}}}%
  \def\ddots{\mathinner{\mkern1mu\raise6.5pt\hbox{.}\mkern2mu\raise3.6pt\hbox{.}\mkern2mu\raise0.7pt\hbox{.}\mkern1mu}}%
  \def\triangle{\mathord{\scalebox{0.9}{$\symup{\Delta}$}}}%
  \def\nabla{\mathord{\raisebox{\depth}{\scalebox{1}[-1]{$\symup{\Delta}$}}}}%
  \def\checkmark{\popcheck{popdark}}%
  \def\star{\mathbin{\ast}}\def\bigstar{\mathord{\ast}}%
  \def\diamond{\mathbin{\scalebox{0.6}{$\Diamond$}}}%
  \def\bigcup{\mathop{\mathchoice{\raisebox{-0.3em}{\scalebox{1.7}{$\cup$}}}{\scalebox{1.25}{$\cup$}}{\cup}{\cup}}\displaylimits}%
  \def\bigcap{\mathop{\mathchoice{\raisebox{-0.3em}{\scalebox{1.7}{$\cap$}}}{\scalebox{1.25}{$\cap$}}{\cap}{\cap}}\displaylimits}%
  \def\lesseqqgtr{\mathrel{\lessgtr}}\def\bigoplus{\mathop{\oplus}\displaylimits}%
}
\providecommand{\ssi}{si et seulement si} % abréviation fréquente
\AtBeginDocument{\providecommand{\wideparen}{\overparen}} % arc de cercle

%% FIGFIX-BEGIN v3 — correctifs globaux des figures (docs/onbuch-generation/tools/figfix)
\usepackage{adjustbox}\usepackage{etoolbox}
% toute figure TikZ/pgfplots plus large que la ligne est réduite à la largeur disponible
\BeforeBeginEnvironment{tikzpicture}{\begin{adjustbox}{max width=\linewidth}}
\AfterEndEnvironment{tikzpicture}{\end{adjustbox}}
% légendes pgfplots lisibles par-dessus les courbes
\pgfplotsset{every axis/.append style={legend style={fill=white,fill opacity=0.92,text opacity=1,draw=black!15,inner sep=3pt}}}
% étiquettes d'arêtes (syntaxe edge["..."]) avec fond blanc
\tikzset{every edge quotes/.append style={fill=white,inner sep=1.2pt}}
% bulles de cartes mentales : texte à la ligne dans la bulle, bulle un peu plus grande
\tikzset{every concept/.append style={text width=2.0cm,align=center,minimum size=2.3cm,inner sep=2pt}}
%% FIGFIX-END
\begin{document}

\pagegarde{Lesson 55 — Reading: Texts on using ICTs as recreational resources at home or in the community}{Anglais}{Tle A}{Chapter 5 : Media and Communication}{EN55}{2 h}{Cette leçon de lecture propose des textes authentiques sur l'usage des TIC comme ressources de loisir à la maison et dans la communauté. Elle entraîne les élèves aux stratégies de skimming, scanning et déduction du sens, tout en enrichissant leur lexique thématique et en les préparant aux épreuves de compréhension écrite du Bac.
\vspace{4pt}
\begin{multicols}{2}\small
\begin{itemize}
\item À la fin de cette leçon, je sais identifier l'idée générale d'un texte grâce au skimming.
\item Je sais repérer rapidement une information précise grâce au scanning.
\item Je sais déduire le sens d'un mot inconnu à partir du contexte.
\item Je connais le lexique essentiel des TIC et des loisirs numériques en anglais.
\item Je sais répondre à des questions de compréhension variées (vrai/faux, QCM, réponses courtes).
\item Je sais distinguer faits et opinions dans un texte.
\end{itemize}\end{multicols}}

\begin{sommaire}
\begin{multicols}{2}
\begin{enumerate}
\item Warm-up and pre-reading
\item Reading text with glossary
\item Reading strategies : skimming, scanning, guessing
\item Comprehension practice
\item Thematic vocabulary : ICTs and digital leisure
\item Guided production : summary and opinion
\item Activité d'intégration
\item Méthodes et exercices résolus
\item Exercices
\item Corrigés des exercices
\item Fiche bilan
\end{enumerate}
\end{multicols}
\end{sommaire}

\begin{objectifs}
\begin{itemize}
\item À la fin de cette leçon, je sais identifier l'idée générale d'un texte grâce au skimming.
\item Je sais repérer rapidement une information précise grâce au scanning.
\item Je sais déduire le sens d'un mot inconnu à partir du contexte.
\item Je connais le lexique essentiel des TIC et des loisirs numériques en anglais.
\item Je sais répondre à des questions de compréhension variées (vrai/faux, QCM, réponses courtes).
\item Je sais distinguer faits et opinions dans un texte.
\item Je sais résumer un texte en quelques phrases avec mes propres mots.
\item Je sais exprimer mon opinion sur les avantages et dangers des loisirs numériques.
\end{itemize}
\end{objectifs}

\begin{prerequis}
\begin{itemize}
\item Lexique de base des médias et de la communication (Première : radio, TV, internet, social media).
\item Présent simple et présent continu pour décrire des habitudes.
\item Expressions d'opinion simples : I think, I believe, in my opinion.
\item Connecteurs de base : and, but, because, however.
\item Stratégies de lecture vues en Première : lire le titre et les images avant le texte.
\end{itemize}
\end{prerequis}

\newpage

\coursec{Warm-up and pre-reading}

It is 8 p.m. in Bamenda. Junior has just finished his homework and picks up his smartphone: thirty minutes of TikTok, a quick game of FIFA Mobile with his cousin in Douala, then a Nollywood film on YouTube. His grandmother complains that “this phone has eaten the boy”, but Junior argues that he is learning English and relaxing at the same time. Are ICTs at home and in our communities a blessing or a trap? In this lesson, we read texts about how people use ICTs for recreation and learn to read them efficiently.

Mais avant de lire le moindre texte, un bon lecteur se prépare. Cette première section est un échauffement : nous allons activer ce que vous savez déjà, faire des prédictions, et découvrir les trois stratégies de lecture qui vous serviront toute l'année — et à l'examen.

\courssub{Brainstorming : what digital leisure do you practise?}

Regardez la scène de Bamenda. Junior pratique, en une seule soirée, trois loisirs numériques différents : \eng{social media} (TikTok), \eng{mobile gaming} (FIFA Mobile) et \eng{streaming} (watching a film on YouTube). Et vous ?

\begin{experience}[Class brainstorming — Our digital free time]
En groupes de quatre, répondez en anglais aux questions suivantes, puis un rapporteur présente les résultats au tableau.
\begin{enumerate}
\item \eng{How many hours a day do you spend on your phone or computer for fun?}
\item \eng{Which apps do you use most: WhatsApp, TikTok, YouTube, games?}
\item \eng{Do you ever go to a cybercafé or a viewing centre? What for?}
\item \eng{Do your parents think ICTs are good or bad for you? Why?}
\end{enumerate}
Phrases utiles pour s'exprimer : \eng{I spend about two hours a day on…} — \eng{I use my phone to…} — \eng{In my opinion, …} — \eng{Most of us… but only a few of us…}
\end{experience}

\begin{definition}[ICTs]
\cle{ICTs} means \emph{Information and Communication Technologies}: all the digital tools we use to communicate and to access information — phones, computers, tablets, the internet, apps, radio and television. « TIC : tous les outils numériques utilisés pour communiquer et accéder à l'information. »
\par
Exemple : \eng{ICTs have changed the way young people relax.} « Les TIC ont changé la façon dont les jeunes se détendent. »
\end{definition}

\begin{attention}[Faux ami : “recreation”]
\eng{Recreation} (prononcez « ré-kri-é-cheune ») signifie « loisirs, détente », et non « récréation » au sens de la pause des écoliers. La « récréation » de la cour d'école se dit \eng{break} ou \eng{playtime}. De même, \eng{actually} signifie « en fait » et non « actuellement » (\eng{currently}).
\end{attention}

\begin{popfigure}
\begin{tikzpicture}[mindmap, grow cyclic, every node/.style={concept, align=center, font=\small},
root concept/.append style={concept color=popblue, font=\bfseries},
level 1/.append style={level distance=3.4cm, sibling angle=60},
level 1 concept/.append style={concept color=poporangeL}]
\node[root concept]{ICT recreation}
child { node[align=center]{gaming\\ \emph{FIFA Mobile, Candy Crush}} }
child { node[align=center]{social media\\ \emph{TikTok, WhatsApp}} }
child { node[align=center]{streaming\\ \emph{films, series}} }
child { node[align=center]{music\\ \emph{Afrobeats, makossa}} }
child { node[align=center]{e-learning fun\\ \emph{quizzes, videos}} }
child { node[align=center]{community centres\\ \emph{cybercafés, viewing centres}} };
\end{tikzpicture}
\legende{Carte mentale : les grandes familles de loisirs numériques (ICT recreation).}
\end{popfigure}

\courssub{Recreation at home vs recreation in the community}

Observez ces deux phrases :

\begin{exemplebox}[Two places, two experiences]
\eng{At home, Junior watches YouTube on his own phone with mobile data.} « À la maison, Junior regarde YouTube sur son propre téléphone avec ses données mobiles. »
\par
\eng{On Saturday, the whole neighbourhood meets at the viewing centre to watch the Premier League.} « Samedi, tout le quartier se retrouve au viewing centre pour regarder la Premier League. »
\end{exemplebox}

La distinction est essentielle pour comprendre les textes de cette leçon. Le loisir numérique peut être \cle{private} (à la maison, seul ou en famille) ou \cle{community-based} (dans un lieu public partagé : cybercafé, \emph{community multimedia centre}, viewing centre).

\begin{center}
\begin{tabularx}{\linewidth}{|X|X|X|}
\hline
\rowcolor{popblueL} At home (private) & In the community (shared) & Remarque \\
\hline
smartphone, tablet & cybercafé & \eng{cybercafé} = même mot qu'en français \\
\hline
home Wi-Fi & community multimedia centre & lieu public avec ordinateurs et internet \\
\hline
personal data bundle & viewing centre & on paie pour regarder un match \\
\hline
headphones, alone & big screen, crowd & \eng{crowd} = foule \\
\hline
\end{tabularx}
\end{center}

\begin{savaistu}[Viewing centres in Cameroon]
Au Cameroun, les \emph{viewing centres} — petites salles où l'on paie une somme modeste pour regarder les matchs de la Premier League ou de la Ligue des champions sur grand écran — sont un exemple typique de loisir numérique \textbf{communautaire}. On y va autant pour le football que pour l'ambiance : commenter, débattre, rire ensemble. Le loisir numérique n'est donc pas toujours solitaire !
\end{savaistu}

\courssub{Activating known vocabulary}

Vous connaissez déjà beaucoup de mots du numérique, souvent identiques ou proches du français. Attention toutefois à l'orthographe et à la prononciation.

\begin{center}
\begin{tabularx}{\linewidth}{|X|X|X|}
\hline
\rowcolor{popblueL} English & Français & Prononciation / Remarque \\
\hline
phone / smartphone & téléphone / smartphone & « faune / smart-fone » \\
\hline
data & données (forfait internet) & « dé-ta » ; \eng{mobile data} = données mobiles \\
\hline
Wi-Fi & Wi-Fi & « ouaï-faï » \\
\hline
app (application) & application & « ap » ; toujours avec l'article : \eng{an app} \\
\hline
to download & télécharger (depuis internet) & « daune-lode » ; opposé : \eng{to upload} \\
\hline
to stream & regarder/écouter en streaming & \eng{streaming} = diffusion en continu \\
\hline
screen & écran & « skrine » \\
\hline
password & mot de passe & « pass-weurd » \\
\hline
\end{tabularx}
\end{center}

\begin{attention}[“Télécharger” : download ou upload ?]
Les francophones confondent souvent les deux. \eng{To download} = télécharger \emph{vers} son appareil (recevoir) ; \eng{to upload} = téléverser, envoyer \emph{depuis} son appareil (publier une vidéo, par exemple). \eng{I downloaded a film} « j'ai téléchargé un film » ; \eng{She uploaded a video to TikTok} « elle a publié une vidéo sur TikTok ».
\end{attention}

\courssub{Predicting from the title and the first paragraph}

Un bon lecteur ne commence jamais un texte « les yeux fermés ». Avant de lire, il observe le titre, les images éventuelles et le premier paragraphe, puis il fait des \cle{predictions} (prédictions).

\begin{methode}[Before you read: ask yourself three questions]
\begin{enumerate}
\item \eng{What is the text about?} — De quoi parle le texte ? (le titre donne le sujet)
\item \eng{Who wrote it and why?} — Qui l'a écrit et dans quel but ? (informer, convaincre, divertir)
\item \eng{What do I already know about the topic?} — Que sais-je déjà sur ce sujet ? (activez vos connaissances et votre vocabulaire)
\end{enumerate}
Ensuite, formulez une prédiction : \eng{I think the text will talk about…} / \eng{The text will probably explain…}
\end{methode}

Appliquons immédiatement cette méthode. Voici le premier paragraphe du texte que nous étudierons dans la section suivante.

\begin{textebox}[“Fun with a Screen” — opening paragraph]
All over Cameroon, from Yaoundé to Buea, free time has gone digital. At home, teenagers chat on WhatsApp, play mobile games and watch films on their phones. In the community, cybercafés and viewing centres welcome young people who do not have internet at home. For many families, ICTs are now the cheapest and most popular form of recreation. But is all this screen time good for us? This text looks at how Cameroonians use ICTs to relax, at home and in their communities, and asks what we gain — and what we may lose.
\end{textebox}

Glossaire : \eng{free time} (n.) temps libre ; \eng{to chat} (v.) bavarder ; \eng{screen time} (n.) temps passé devant un écran ; \eng{to relax} (v.) se détendre ; \eng{to gain} (v.) gagner ; \eng{to lose} (v.) perdre.

D'après ce seul paragraphe, on peut prédire que le texte : (1) parlera des loisirs numériques au Cameroun ; (2) comparera la maison et la communauté ; (3) présentera des avantages \emph{et} des inconvénients. Notez la question rhétorique \eng{is all this screen time good for us?} : elle annonce un texte équilibré, pas un simple éloge.

\courssub{The three reading strategies of the day}

Trois stratégies vont transformer votre façon de lire en anglais. Retenez leurs noms anglais : ils apparaissent souvent dans les consignes d'examen.

\begin{definition}[Skimming, scanning, guessing from context]
\begin{itemize}
\item \cle{Skimming} (« survol ») : lire rapidement le titre, la première phrase de chaque paragraphe et la conclusion pour saisir l'\textbf{idée générale} (\eng{the gist}). On ne lit pas chaque mot.
\item \cle{Scanning} (« repérage ») : parcourir le texte des yeux pour trouver une \textbf{information précise} (un nom, une date, un chiffre), comme on cherche un numéro dans un annuaire.
\item \cle{Guessing from context} (« déduction ») : deviner le sens d'un mot inconnu grâce aux mots qui l'entourent, sans dictionnaire.
\end{itemize}
\end{definition}

\begin{exemplebox}[Guessing from context in action]
Phrase : \eng{Junior is addicted to his phone: he cannot stop checking it, even during meals.}
\par
Vous ne connaissez pas \eng{addicted} ? Le contexte vous aide : « il ne peut pas s'arrêter de le consulter, même pendant les repas » → \eng{addicted} signifie « accro, dépendant ». Indices possibles : la définition après les deux-points, les synonymes, les exemples, le contraire (\eng{but}, \eng{however}).
\end{exemplebox}

\begin{exoresolu}[Prédiction et survol]
Énoncé. Voici le titre d'un article : \eng{“Why Grandmother Now Loves WhatsApp”}. (a) Prédisez le contenu en une phrase. (b) La première phrase est : \eng{Sixty-five-year-old Mama Rose from Bafoussam used to hate mobile phones.} Que va probablement raconter la suite ?
\tcblower
Solution. (a) \eng{I think the text will explain how an old woman discovered WhatsApp and now enjoys using it.} Le titre oppose \eng{grandmother} (on s'attend à une personne âgée éloignée de la technologie) et \eng{loves WhatsApp} (surprise) : l'article racontera un changement d'attitude. (b) \eng{Used to hate} (« elle détestait autrefois ») annonce un contraste passé/présent : la suite expliquera comment et pourquoi elle a changé d'avis — probablement grâce à un petit-enfant qui lui a montré l'application. Stratégie utilisée : \emph{skimming} du titre et de la première phrase + prédiction.
\end{exoresolu}

\begin{exercice}[Warm-up practice]
1. Classez ces activités en deux colonnes \eng{at home} / \eng{in the community} : \eng{watching a match at a viewing centre ; chatting on WhatsApp in bed ; using a computer at the cybercafé ; playing FIFA on your phone ; printing documents at a multimedia centre}.
\par
2. Devinez le sens de \eng{cheap} dans : \eng{For many families, ICTs are the cheapest form of recreation: a data bundle costs less than a cinema ticket.}
\par
3. En une phrase anglaise, prédisez le contenu d'un texte intitulé \eng{“The Cybercafé is Dead: Long Live the Smartphone!”}.
\end{exercice}

\begin{corrige}[Exercice « Warm-up practice »]
1. \eng{At home}: chatting on WhatsApp in bed ; playing FIFA on your phone. \eng{In the community}: watching a match at a viewing centre ; using a computer at the cybercafé ; printing documents at a multimedia centre.
\par
2. \eng{Cheap} signifie « bon marché, pas cher » : le contexte le dit explicitement — un forfait coûte \emph{moins cher} qu'une place de cinéma (comparaison \eng{costs less than}).
\par
3. Réponse possible : \eng{I think the text will explain that smartphones have replaced cybercafés because people now have internet at home.} Le titre joue sur la formule “The King is dead, long live the King!” : il annonce le déclin des cybercafés face au téléphone personnel.
\end{corrige}

\begin{aretenir}[L'essentiel de l'échauffement]
\begin{itemize}
\item Avant de lire : observez le titre et le premier paragraphe, posez-vous trois questions (\eng{What? Who? What do I know?}) et faites une prédiction.
\item Distinguez le loisir numérique \textbf{privé} (at home) du loisir \textbf{communautaire} (cybercafé, viewing centre).
\item Trois stratégies : \eng{skimming} (idée générale), \eng{scanning} (information précise), \eng{guessing from context} (mots inconnus).
\item Attention aux faux amis : \eng{recreation} = loisirs ; \eng{actually} = en fait ; \eng{download} ≠ \eng{upload}.
\end{itemize}
\end{aretenir}

Vous êtes maintenant prêts : dans la section suivante, nous lirons le texte complet \eng{“Fun with a Screen”} et nous mettrons ces stratégies à l'épreuve.

\coursec{Reading text with glossary}

Après notre échauffement (\eng{warm-up}) et nos prédictions sur le contenu du texte — rappelez-vous : nous avions deviné, à partir du titre et de l'image, que le texte parlerait des loisirs numériques des jeunes Camerounais —, il est temps de passer à la lecture proprement dite. Dans cette section, nous allons lire le texte \emph{deux fois} : une première fois rapidement, pour vérifier nos prédictions (c'est le \cle{skimming}), puis une seconde fois lentement, avec le glossaire, pour comprendre les détails (c'est la \cle{detailed reading}). C'est exactement la démarche exigée à l'épreuve d'anglais du baccalauréat A.

\courssub{First reading : silent reading and skimming}

\begin{methode}[Comment faire une première lecture globale (skimming)]
\begin{enumerate}
\item Lis le texte \textbf{en silence}, sans dictionnaire, sans t'arrêter sur les mots inconnus.
\item Lis le \textbf{titre}, le \textbf{premier paragraphe} et la \textbf{première phrase de chaque paragraphe} avec plus d'attention : ils contiennent l'essentiel.
\item Pose-toi une seule question : \eng{What is this text about?} (De quoi parle ce texte ?)
\item Vérifie tes prédictions de la section précédente : avais-tu deviné juste ?
\item Formule le thème en une phrase, par exemple : \eng{The text is about how Cameroonian families use ICTs for entertainment.}
\end{enumerate}
\end{methode}

\begin{attention}[Ne traduis pas mot à mot !]
L'erreur n° 1 des élèves francophones est de traduire le texte \textbf{mot à mot} et de paniquer dès qu'un mot est inconnu. Lis d'abord le \textbf{paragraphe entier} pour saisir le sens global. Très souvent, le contexte suffit à deviner le sens d'un mot. Exemple : dans la phrase \eng{whole families gather around a laptop to watch Nigerian films}, même si tu ne connais pas \eng{gather}, le contexte (\eng{whole families}, \eng{around a laptop}) te permet de deviner « se rassembler ».
\end{attention}

\courssub{The reading text}

\begin{textebox}[Fun at Our Fingertips — a magazine article]
In many Cameroonian homes today, the smartphone has become the new television. After school, teenagers stream music, watch comedy skits on YouTube and chat with friends on WhatsApp. At weekends, whole families gather around a laptop to watch Nigerian films. In towns, young people who have no internet at home go to cybercafés or community multimedia centres, where they play online games, learn video editing or follow football matches live.

Mrs Ngum, a mother of four in Buea, says: “My children relax with their phones, but they also learn English songs and typing skills.” However, some parents worry. “My son spends five hours a day gaming and forgets his books,” complains Mr Ekane.

Experts advise balance: ICTs are wonderful recreational resources when they are used wisely — for fun, creativity and learning — but harmful when they replace sleep, sport and family conversation. The key, they say, is not to ban technology, but to manage it.
\end{textebox}

\begin{definition}[Glossary — Glossaire]
\textbf{fingertips} (n. pl.) = bout des doigts ; \textbf{to stream} (v.) = regarder/écouter en ligne (streaming) ; \textbf{skit} (n.) = sketch comique, saynète ; \textbf{to gather} (v.) = se rassembler, se réunir ; \textbf{typing skills} (n. pl.) = compétences en saisie (au clavier) ; \textbf{wisely} (adv.) = sagement, avec discernement ; \textbf{harmful} (adj.) = nuisible, préjudiciable ; \textbf{to ban} (v.) = interdire ; \textbf{to manage} (v.) = gérer, administrer.
\end{definition}

\courssub{Second reading : detailed reading with the glossary}

\begin{methode}[Comment faire une lecture détaillée]
\begin{enumerate}
\item Relis le texte \textbf{paragraphe par paragraphe}, lentement.
\item Souligne les mots inconnus. Avant de consulter le glossaire, essaie de \textbf{deviner} leur sens grâce au contexte, à la famille de mots ou à la ressemblance avec le français.
\item Consulte ensuite le glossaire pour vérifier tes hypothèses.
\item Repère les \textbf{mots de liaison} (\eng{however}, \eng{but}, \eng{when}) : ils indiquent la structure argumentative du texte.
\item Résume chaque paragraphe en une phrase anglaise simple.
\end{enumerate}
\end{methode}

\begin{propriete}[Déduire le sens d'un mot inconnu : trois stratégies]
\begin{enumerate}
\item \textbf{Le contexte} : \eng{My son spends five hours a day gaming} → \eng{gaming} est une activité qui dure des heures sur écran → « jouer à des jeux vidéo ».
\item \textbf{La famille de mots} : \eng{harmful} = \eng{harm} (mal, tort) + suffixe \eng{-ful} (plein de) → « qui fait du mal, nuisible ». De même : \eng{wonderful} = \eng{wonder} + \eng{-ful}.
\item \textbf{Les faux amis à surveiller} : certains mots ressemblent au français mais ont un sens différent (voir le tableau ci-dessous).
\end{enumerate}
\end{propriete}

\begin{exemplebox}[Guessing words from context — exemples tirés du texte]
\begin{itemize}
\item \eng{teenagers stream music} → après \eng{music}, verbe d'usage numérique → « écouter en streaming ».
\item \eng{comedy skits on YouTube} → associé à \eng{comedy} → « courts sketchs comiques ».
\item \eng{cybercafés or community multimedia centres} → lieux (\eng{go to}) où l'on utilise internet.
\item \eng{they replace sleep, sport and family conversation} → \eng{replace} = « remplacent » (transparent avec le français).
\end{itemize}
\end{exemplebox}

\courssub{Reading strategy : Scanning (reading for specific information)}

\begin{methode}[Comment faire du scanning (recherche d'informations précises)]
\begin{enumerate}
\item Identifie le \textbf{mot-clé} ou le \textbf{chiffre} cherché dans la question (ex. : \eng{How many hours?} → cherche un nombre ; \eng{Who?} → cherche un nom propre).
\item Fais glisser ton regard \textbf{rapidement} sur le texte \textbf{sans lire chaque mot}.
\item Repère le mot-clé (ou un synonyme) dans le texte.
\item Lis \textbf{uniquement la phrase concernée} pour extraire la réponse.
\item Écris la réponse courte et précise.
\end{enumerate}
\end{methode}

\begin{exercice}[Scanning practice — Find the information quickly]
Ne relis pas tout le texte. Trouve uniquement les réponses aux questions ci-dessous :
\begin{enumerate}
\item What device has become the \eng{new television} in Cameroonian homes?
\item Where do young people go if they have no internet at home?
\item What does Mrs Ngum say her children learn?
\item How many hours a day does Mr Ekane's son spend gaming?
\item What do experts advise instead of banning technology?
\end{enumerate}
\end{exercice}

\begin{corrige}[Exercice 2 — Scanning]
\begin{enumerate}
\item The smartphone. (Paragraphe 1, phrase 1)
\item Cybercafés or community multimedia centres. (Paragraphe 1, phrase 3)
\item English songs and typing skills. (Paragraphe 2, citation de Mrs Ngum)
\item Five hours. (Paragraphe 2, citation de Mr Ekane)
\item To manage it (wisely) / Balance. (Paragraphe 3, dernière phrase)
\end{enumerate}
\end{corrige}

\courssub{Identifying the text type}

\begin{propriete}[Quel type de texte ?]
Ce texte est un \textbf{article de magazine} (ou un \eng{blog post}) de type \textbf{reportage social}. Indices :
\begin{itemize}
\item un titre accrocheur et imagé : \eng{Fun at Our Fingertips} (jeu de mots sur \eng{at our fingertips} = « à portée de main ») ;
\item des \textbf{témoignages cités entre guillemets} (\eng{Mrs Ngum says...}, \eng{complains Mr Ekane}) ;
\item une structure \textbf{équilibrée} : avantages (\eng{relax, learn}) puis inquiétudes (\eng{However, some parents worry}) puis conseil d'experts (\eng{Experts advise balance}) ;
\item un ton informatif mais accessible, sans statistiques savantes.
\end{itemize}
\end{propriete}

\begin{center}
\begin{tabularx}{\linewidth}{|X|X|X|}
\hline
\rowcolor{popblueL} English & Français & Remarque \\
\hline
recreational resources & ressources de loisirs & \eng{recreation} = loisir, détente \\
\hline
to relax & se détendre & faux ami : ≠ « relâcher » \\
\hline
typing skills & compétences en saisie & \eng{to type} = taper au clavier \\
\hline
harmful & nuisible & \eng{harm} + \eng{-ful} ; contraire : \eng{harmless} \\
\hline
wisely & sagement & adverbe en \eng{-ly} formé sur \eng{wise} \\
\hline
to ban / to manage & interdire / gérer & \eng{manage} ≠ « manier » ; = « gérer, s'en sortir » \\
\hline
actually (piège hors texte) & en fait & faux ami : ≠ « actuellement » (\eng{currently}) \\
\hline
\end{tabularx}
\end{center}

\courssub{Thematic vocabulary : ICTs and digital leisure}

\begin{definition}[Vocabulaire thématique : les TIC et les loisirs numériques]
\begin{center}
\begin{tabular}{|l|l|l|}
\hline
\rowcolor{popblueL} \textbf{Category} & \textbf{English} & \textbf{Français} \\
\hline
Devices & smartphone, laptop, tablet, desktop computer & téléphone intelligent, portable, tablette, ordinateur fixe \\
\hline
Places & cybercafé, multimedia centre, community centre & cybercafé, centre multimédia, centre communautaire \\
\hline
Actions (Verbs) & stream, chat, browse, download, upload, edit, game & regarder en ligne, discuter, naviguer, télécharger, mettre en ligne, monter (vidéo), jouer \\
\hline
Content & skit, playlist, podcast, livestream, tutorial, vlog & sketch, liste de lecture, podcast, direct, tutoriel, vlog \\
\hline
Adjectives & addictive, entertaining, educational, harmful, harmless, virtual & addictif, divertissant, éducatif, nuisible, inoffensif, virtuel \\
\hline
Nouns (Concepts) & screen time, bandwidth, connectivity, digital literacy, balance & temps d'écran, bande passante, connectivité, culture numérique, équilibre \\
\hline
\end{tabular}
\end{center}
\end{definition}

\begin{popfigure}
\begin{tikzpicture}[mindmap, grow cyclic, every node/.style={concept, align=center}, root concept/.append style={concept color=popblue, font=\small\bfseries}, level 1/.append style={level distance=3.2cm, sibling angle=90}, concept connection/.append style={color=popmuted}]
\node[root concept]{ICTs as recreational resources}
  child[concept color=popgreen]{ node[align=center]{Fun\\ \footnotesize stream music\\ \footnotesize watch skits\\ \footnotesize online games} }
  child[concept color=poporange]{ node[align=center]{Learning\\ \footnotesize English songs\\ \footnotesize typing skills\\ \footnotesize video editing} }
  child[concept color=poppink]{ node[align=center]{Dangers\\ \footnotesize less sleep\\ \footnotesize less sport\\ \footnotesize less conversation} }
  child[concept color=popgold]{ node[align=center]{Advice\\ \footnotesize balance\\ \footnotesize manage, not ban} };
\end{tikzpicture}
\legende{Carte mentale des idées du texte : les quatre grands axes de l'article.}
\end{popfigure}

\begin{exoresolu}[Check your global understanding]
Énoncé — Réponds en anglais par une phrase complète : \eng{What is the text about? What do experts advise?}
\tcblower
Solution — \eng{The text is about how Cameroonian teenagers and families use ICTs — smartphones, laptops and the internet — as recreational resources at home and in the community. Experts advise balance: we should not ban technology, but manage it wisely, so that it does not replace sleep, sport and family conversation.} (« Le texte explique comment les adolescents et les familles camerounais utilisent les TIC comme ressources de loisirs. Les experts conseillent l'équilibre : ne pas interdire la technologie, mais la gérer sagement. »)
\end{exoresolu}

\begin{exercice}[Before the glossary — guess first!]
Sans consulter le glossaire, devine le sens des mots soulignés grâce au contexte, puis vérifie : \eng{whole families \textbf{gather} around a laptop} ; \eng{they also learn \textbf{typing skills}} ; \eng{ICTs are wonderful but \textbf{harmful} when they replace sleep}.
\end{exercice}

\begin{corrige}[Exercice 1]
\eng{gather} = se rassembler (contexte : \eng{whole families around a laptop}) ; \eng{typing skills} = compétences en saisie au clavier (contexte : \eng{learn}, activité sur téléphone/ordinateur) ; \eng{harmful} = nuisible (contexte : opposition avec \eng{wonderful} introduite par \eng{but}). Bravo si tu as deviné avant de regarder le glossaire : c'est la compétence clé du \eng{reading} au baccalauréat.
\end{corrige}

\courssub{Guided production : Summary and Opinion}

\begin{methode}[Comment rédiger un résumé et donner son avis (épreuve écrite)]
\begin{enumerate}
\item \textbf{Résumé (50--60 mots)} : Identifie l'idée principale de chaque paragraphe (1 phrase par paragraphe). Relie-les avec des connecteurs (\eng{Firstly, Secondly, Finally}). N'inclus pas de détails (noms, chiffres) ni ton opinion.
\item \textbf{Opinion (30--40 mots)} : Utilise \eng{In my opinion / I think that / As far as I am concerned}. Donne un argument personnel + un exemple concret (ton expérience, ton entourage). Utilise le vocabulaire du thème (\eng{screen time, balance, addictive, educational}).
\item \textbf{Relecture} : Vérifie l'orthographe, les temps (présent simple pour faits généraux), la ponctuation et le nombre de mots.
\end{enumerate}
\end{methode}

\begin{experience}[Production guidée : Écris ton paragraphe]
\textbf{Consigne :} Écris un paragraphe de 80--100 mots au total (résumé + opinion) sur le sujet : \eng{``ICTs as recreational resources in Cameroon: benefits and risks.''}
\begin{enumerate}
\item Rédige d'abord ton brouillon en suivant la méthode ci-dessus.
\item Échange avec ton voisin pour une relecture croisée (vérification du vocabulaire, de la structure, de l'anglais).
\item Recopie ta version finale proprement.
\end{enumerate}
\textbf{Amorce possible :} \eng{The text highlights that ICTs have become essential for entertainment in Cameroon. Teenagers stream music and play games, while families watch movies together. However, addiction is a real risk. In my opinion, ICTs are fantastic tools for learning and fun, provided parents set clear limits on screen time. For instance, my younger brother improved his English by watching tutorials, but he now negotiates his gaming hours with my parents. Balance is the key.}
\end{experience}

\begin{savaistu}[Le saviez-vous ? — Le Cameroun et le numérique]
Le Cameroun compte plus de \textbf{10 millions d'internautes} (données 2023, ARTP). Le mobile money (Orange Money, MTN MoMo) a révolutionné l'inclusion financière. Les \eng{cybercafés} restent des lieux majeurs d'accès à internet pour les jeunes hors grands centres urbains. Le terme \eng{Silicon Mountain} (Buea) désigne l'écosystème technologique anglophone en plein essor (start-ups, incubateurs comme ActivSpaces). Le gouvernement promeut le \eng{Cameroun Numérique 2020} pour développer l'économie numérique.
\end{savaistu}

Maintenant que le texte est compris globalement et dans ses détails, et que tu as pratiqué le \eng{skimming}, le \eng{scanning} et l'inférence lexicale, tu es prêt à aborder les exercices de compréhension de type bac et à produire ton propre texte argumentatif sur ce thème.

\coursec{Reading strategies : skimming, scanning, guessing}

Dans la section précédente, nous avons lu le texte \eng{“Screen Time and Free Time”} avec son glossaire et nous avons découvert comment Mr Ekane et sa famille utilisent les TIC comme ressources de loisir à la maison. Mais au Bac, personne ne vous donnera un glossaire, et le temps sera compté. Il faut donc apprendre à lire \cle{stratégiquement}. Trois stratégies professionnelles vont transformer votre façon de lire : \cle{skimming}, \cle{scanning} et \cle{guessing from context}. Ce sont exactement les stratégies attendues dans les épreuves de compréhension de l'écrit.

\courssub{Observation : deux lecteurs, deux méthodes}

\begin{textebox}[Two readers, one exam]
Read this short scene before studying the rules.

\textbf{Student A} opens the paper and starts translating every single word of the text, line by line, with his dictionary. After twenty minutes, he has read only one paragraph and he has not answered any question.

\textbf{Student B} first reads the title, the first sentence of each paragraph and the conclusion. In two minutes, she knows the text is about ICTs used for leisure at home. Then she reads the questions. Question 1 asks: \eng{“How many hours does Mr Ekane's son spend gaming every day?”} She does not re-read everything: her eyes slide over the text until she finds the name \eng{“Ekane”} and a number. In thirty seconds, she has the answer: \eng{“three hours”}. When she meets the unknown word \eng{“wisely”}, she does not panic: the sentence contrasts it with \eng{“harmful”}, so she guesses it means \eng{“in a good, sensible way”}.

Student B finishes early and gets excellent marks. Student A runs out of time.
\end{textebox}

\textbf{Glossaire :} \emph{to run out of time} (v.) : manquer de temps ; \emph{to slide} (v.) : glisser ; \emph{marks} (n.) : notes ; \emph{to guess} (v.) : deviner ; \emph{sensible} (adj.) : raisonnable, sensé.

\textbf{Question d'observation :} quelle différence fondamentale voyez-vous entre les deux élèves ? Student A lit \emph{tout} de la même façon ; Student B \emph{adapte sa lecture à son objectif}. C'est exactement cela, une stratégie de lecture.

\courssub{Skimming : lire vite pour l'idée générale}

\begin{definition}[Skimming]
Le \cle{skimming} (lecture rapide, « survol ») consiste à lire un texte très vite pour en dégager \textbf{l'idée générale} (the gist), sans chercher à comprendre chaque mot. On lit sélectivement : le titre, la première phrase de chaque paragraphe (la \emph{topic sentence}) et la dernière phrase ou la conclusion.
\end{definition}

\begin{methode}[SKIMMING — la méthode en 4 étapes]
\begin{enumerate}
\item Lis le \textbf{titre} et regarde les éventuelles illustrations : ils annoncent le sujet.
\item Lis la \textbf{première phrase de chaque paragraphe} : en anglais journalistique, elle contient presque toujours l'idée principale du paragraphe.
\item Lis la \textbf{dernière phrase} ou la conclusion : elle donne souvent l'opinion ou le message de l'auteur.
\item \textbf{Ne t'arrête jamais sur un mot inconnu.} Saute-le et continue : l'idée générale n'a pas besoin de tous les mots.
\end{enumerate}
\end{methode}

\begin{exemplebox}[Skimming appliqué à notre texte]
Appliquons le skimming au texte de la section 2 :
\begin{itemize}
\item \textbf{Titre :} \eng{“Screen Time and Free Time”} → le texte parle des écrans et des loisirs.
\item \textbf{Première phrase du §1 :} \eng{“In many Cameroonian homes, ICTs have become the favourite recreational resource.”} → sujet : les TIC comme loisirs.
\item \textbf{Première phrase du §2 :} \eng{“Mr Ekane's family in Yaoundé is a good example.”} → un exemple concret va suivre.
\item \textbf{Conclusion :} \eng{“ICTs are wonderful recreational resources when they are used wisely.”} → message : les TIC sont positifs \emph{si} on les utilise avec sagesse.
\end{itemize}
En moins de deux minutes, sans traduire, nous connaissons le sujet, la structure et le message du texte. \textbf{Traduction de la conclusion :} « Les TIC sont de merveilleuses ressources de loisir quand elles sont utilisées sagement. »
\end{exemplebox}

\begin{attention}[Comparaison avec le français]
En français scolaire, on vous apprend souvent à lire « attentivement, du début à la fin ». En compréhension de l'écrit en anglais, c'est l'inverse : lire \emph{tout} lentement est une \textbf{perte de temps}. Le skimming n'est pas de la paresse : c'est une compétence évaluée. Les questions du type \eng{“What is the text about?”} ou \eng{“Choose the best title”} se traitent \emph{uniquement} par skimming.
\end{attention}

\courssub{Scanning : chercher une information précise}

\begin{definition}[Scanning]
Le \cle{scanning} (lecture ciblée, « balayage ») consiste à chercher \textbf{une information précise} (une date, un nom, un chiffre, un lieu) sans lire le reste du texte. Les yeux glissent rapidement sur les lignes jusqu'à repérer un \textbf{mot-repère} (a keyword), puis on lit attentivement seulement la phrase qui l'entoure.
\end{definition}

\begin{methode}[SCANNING — la méthode en 4 étapes]
\begin{enumerate}
\item Lis d'abord la \textbf{question} (pas le texte !) et \textbf{souligne ses mots-clés} : nom propre, chiffre, verbe d'action.
\item Prévois la \textbf{forme} du mot-repère dans le texte : un nom propre aura une majuscule, une durée sera un chiffre suivi de \eng{hours}, \eng{minutes}, \eng{days}.
\item Fais \textbf{glisser tes yeux} en diagonale sur le texte, sans lire les phrases, jusqu'à trouver le mot-repère ou un mot de la même famille (\eng{gaming} ↔ \eng{games}).
\item Lis attentivement \textbf{uniquement la phrase} qui contient le repère, et réponds.
\end{enumerate}
\end{methode}

\begin{exemplebox}[Exemple traité pas à pas : \eng{“How many hours does Mr Ekane's son spend gaming?”}]
\textbf{Étape 1 — Mots-clés de la question :} \eng{how many hours} (on cherche un \textbf{chiffre} + \eng{hours}) ; \eng{Mr Ekane's son} (on cherche le \textbf{nom propre} “Ekane” ou le mot \eng{son}) ; \eng{gaming} (mot de la famille de \eng{games}).

\textbf{Étape 2 — Balayage :} les yeux glissent sur le texte… on repère “Ekane” au paragraphe 2, puis juste après : \eng{“His sixteen-year-old son, Junior, spends three hours every day playing online games with his friends.”}

\textbf{Étape 3 — Réponse :} \eng{“Mr Ekane's son spends three hours gaming every day.”} (« Le fils de M. Ekane passe trois heures par jour à jouer. »)

Temps total : environ 30 secondes, sans avoir relu tout le texte.
\end{exemplebox}

\begin{attention}[Au Bac, ne perds pas de temps à tout traduire]
L'erreur classique du candidat francophone : traduire mentalement tout le texte avant de regarder les questions. Résultat : il ne reste plus de temps pour rédiger. \textbf{Réponds d'abord aux questions faciles par scanning} (celles qui demandent un chiffre, un nom, une date), puis traite les questions d'interprétation. Et n'écris jamais une réponse en français : la réponse doit être en anglais, souvent en reprenant les mots du texte.
\end{attention}

\courssub{Guessing from context : déduire le sens d'un mot inconnu}

\begin{definition}[Guessing meaning from context]
Déduire le sens d'un mot inconnu (\cle{guessing}) consiste à exploiter les \textbf{indices de contexte} au lieu de chercher dans le dictionnaire. Les indices les plus fréquents sont :
\begin{itemize}
\item la \textbf{définition} ou la paraphrase : \eng{“a cybercafé, that is, a place where people pay to use the Internet”} (\eng{that is} = c'est-à-dire) ;
\item le \textbf{synonyme} : \eng{“the film was amusing and entertaining”} (\eng{amusing} ≈ \eng{entertaining}) ;
\item l'\textbf{exemple} : \eng{“outdoor games such as football and hide-and-seek”} ;
\item le \textbf{contraste} : \eng{“gaming is not harmful; it is fun when used wisely”} (\eng{but}, \eng{while}, \eng{unlike}, \eng{instead of} signalent un contraire) ;
\item la \textbf{liste} : une énumération révèle la catégorie du mot (\eng{“fun, creativity and learning”} = des choses positives).
\end{itemize}
\end{definition}

\begin{exemplebox}[Exemple traité : déduire le sens de “wisely”]
Phrase du texte : \eng{“ICTs are wonderful recreational resources when they are used wisely, but they can be harmful when they are misused.”}

\textbf{Étape 1 — Repérer les connecteurs :} \eng{but} annonce un \textbf{contraste}.

\textbf{Étape 2 — Identifier les contraires :} d'un côté \eng{wonderful} (merveilleux) et \eng{used wisely} ; de l'autre \eng{harmful} (nuisible) et \eng{misused} (mal utilisés). Donc \eng{wisely} est le \textbf{contraire} de \eng{misused}.

\textbf{Étape 3 — Vérifier avec la liste :} le texte ajoute \eng{“they bring fun, creativity and learning”} (« ils apportent plaisir, créativité et apprentissage ») : trois choses \textbf{positives}. Donc \eng{wisely} = « sagement, de façon raisonnable ».

\textbf{Traduction complète :} « Les TIC sont de merveilleuses ressources de loisir quand elles sont utilisées sagement, mais elles peuvent être nuisibles quand on en abuse. »
\end{exemplebox}

\begin{propriete}[Les connecteurs : des panneaux indicateurs pour deviner]
Les connecteurs logiques sont vos meilleurs alliés pour déduire le sens :
\begin{center}
\begin{tabularx}{\linewidth}{|l|X|X|}\hline
\rowcolor{popblueL} \textbf{Connecteur} & \textbf{Fonction} & \textbf{Exemple} \\ \hline
\eng{but, however, while} & contraste & \eng{It is not harmful but helpful.} \\ \hline
\eng{that is, in other words} & reformulation & \eng{a gadget, that is, a small device} \\ \hline
\eng{such as, for example} & exemple & \eng{games such as chess} \\ \hline
\eng{because, so} & cause / conséquence & \eng{He is addicted, so he plays all night.} \\ \hline
\end{tabularx}
\end{center}
\end{propriete}

\courssub{Tableau récapitulatif des trois stratégies}

\begin{center}
\begin{tabularx}{\linewidth}{|l|X|X|X|}\hline
\rowcolor{popblueL} \textbf{Stratégie} & \textbf{Objectif} & \textbf{Comment ?} & \textbf{Type de question} \\ \hline
\cle{Skimming} & Idée générale (gist) & Titre + 1\textsuperscript{re} phrase de chaque § + conclusion, très vite & \eng{What is the text about?} \\ \hline
\cle{Scanning} & Information précise & Yeux qui glissent jusqu'au mot-repère (nom, chiffre, date) & \eng{How many hours…? Where…? When…?} \\ \hline
\cle{Guessing} & Sens d'un mot inconnu & Indices de contexte : contraste, synonyme, exemple, définition & \eng{What does “wisely” mean?} \\ \hline
\end{tabularx}
\end{center}

\begin{popfigure}
\begin{tikzpicture}[
  font=\small,
  quest/.style={rectangle, rounded corners, draw=poppurple, fill=poppurpleL, align=center, text width=3.4cm, minimum height=1cm},
  test/.style={diamond, draw=popblue, fill=popblueL, align=center, text width=2.6cm, aspect=2},
  strat/.style={rectangle, rounded corners, draw=popgreen, fill=popgreenL, align=center, text width=3cm, minimum height=1.1cm},
  fleche/.style={-{Stealth[length=3mm]}, thick, popdark}
]
\node[quest, align=center] (q) at (0,0){\textbf{QUESTION}\\ à traiter};
\node[test] (t1) at (5,1.8) {Idée générale ?};
\node[test] (t2) at (5,0) {Information précise ?};
\node[test] (t3) at (5,-1.8) {Mot inconnu ?};
\node[strat, align=center] (s1) at (10,1.8){\textbf{SKIMMING}\\ titre + 1\textsuperscript{re} phrase + conclusion};
\node[strat, align=center] (s2) at (10,0){\textbf{SCANNING}\\ mot-repère : nom, chiffre, date};
\node[strat, align=center] (s3) at (10,-1.8){\textbf{CONTEXTE}\\ contraste, synonyme, exemple};
\draw[fleche] (q) -- (t1);
\draw[fleche] (q) -- (t2);
\draw[fleche] (q) -- (t3);
\draw[fleche] (t1) -- (s1);
\draw[fleche] (t2) -- (s2);
\draw[fleche] (t3) -- (s3);
\end{tikzpicture}
\legende{Organigramme de décision : quelle stratégie pour quelle question ?}
\end{popfigure}

\courssub{Entraînement guidé}

\begin{textebox}[Digital leisure in Buea]
Read this short text and apply the three strategies.

In Buea, at the foot of Mount Cameroon, young people have changed the way they spend their free time. Ten years ago, most teenagers played outdoor games after school. Today, many of them prefer digital leisure.

Take the example of Mrs Ndi's family. Her daughter, Claris, aged seventeen, spends two hours every evening watching films on her tablet. Her younger brother, Godlove, prefers educational apps: he learns English vocabulary with a game on his mother's phone. On Saturdays, the whole family gathers in the living room to watch football matches on television.

However, digital leisure has a dark side. Some teenagers in Buea spend so much time online that they neglect their homework and their sleep. Doctors say that too much screen time can damage the eyes and cause headaches.

Mrs Ndi has found a solution: during the week, the children may use screens for only one hour a day, and never after nine o'clock at night. “ICTs are wonderful recreational resources when they are used wisely,” she says, “but they can be harmful when they are misused. Screens should bring fun, creativity and learning, not problems.”
\end{textebox}

\textbf{Glossaire :} \emph{free time} (n.) : temps libre ; \emph{to gather} (v.) : se réunir ; \emph{a dark side} (n.) : un aspect négatif ; \emph{to neglect} (v.) : négliger ; \emph{to damage} (v.) : abîmer ; \emph{harmful} (adj.) : nuisible ; \emph{misused} (adj.) : mal utilisé.

\begin{exoresolu}[Applying the three strategies]
\textbf{1. (Skimming)} What is the text about? Answer in one sentence. \textbf{2. (Scanning)} How many hours does Claris spend watching films every evening? \textbf{3. (Guessing)} What does “neglect” mean in paragraph 3? Use the context.
\tcblower
\textbf{1. Skimming :} titre + premières phrases + conclusion → \eng{“The text is about how a family in Buea uses ICTs for leisure, and the need to use screens wisely.”} (« Le texte parle de l'usage des TIC comme loisirs dans une famille de Buea et de la nécessité d'utiliser les écrans sagement. »)

\textbf{2. Scanning :} mots-clés : \eng{Claris} + \eng{how many hours} → on repère le nom propre “Claris” au paragraphe 2, puis le chiffre : \eng{“spends two hours every evening”}. Réponse : \eng{“Claris spends two hours watching films every evening.”} (« Claris passe deux heures chaque soir à regarder des films. »)

\textbf{3. Guessing :} indice de conséquence : \eng{“spend so much time online \textbf{that} they neglect their homework and their sleep”}. Passer trop de temps en ligne a une conséquence négative sur les devoirs et le sommeil ; le paragraphe annonce \eng{“a dark side”}. Donc \eng{neglect} = « négliger, délaisser ».
\end{exoresolu}

\begin{exercice}[Your turn]
Toujours avec le texte \eng{“Digital leisure in Buea”} :
\begin{enumerate}
\item (Scanning) At what time must the children stop using screens at night? \eng{(indice : nine o'clock)}
\item (Guessing) What does \eng{“gathers”} mean in paragraph 2? (Indice : \eng{the whole family… in the living room}.)
\item (Skimming) Which is the best title: (a) \eng{The dangers of television} (b) \eng{Digital leisure in a Buea family} (c) \eng{Football in Cameroon}?
\end{enumerate}
\end{exercice}

\begin{corrige}[Exercice : Your turn]
\textbf{1.} Scanning avec le repère horaire \eng{nine o'clock} : \eng{“They must stop using screens at nine o'clock at night.”} (« Ils doivent arrêter les écrans à 21 h. »)
\textbf{2.} La famille entière se retrouve dans le salon pour regarder un match : \eng{gathers} = « se réunit ».
\textbf{3.} Le meilleur titre est (b) : le texte parle des loisirs numériques d'une famille de Buea, avec leurs avantages et leurs limites ; (a) est trop étroit et (c) est hors sujet.
\end{corrige}

\begin{aretenir}[L'essentiel de la section]
\begin{itemize}
\item \textbf{Skimming} = lecture rapide pour l'idée générale : titre, premières phrases, conclusion. Jamais de dictionnaire à cette étape.
\item \textbf{Scanning} = recherche ciblée d'une information précise grâce à un mot-repère (nom propre, chiffre, date).
\item \textbf{Guessing} = déduction du sens d'un mot inconnu grâce aux indices de contexte : contraste (\eng{but}), reformulation (\eng{that is}), exemples (\eng{such as}), listes.
\item Au Bac : questions faciles d'abord (scanning), jamais de traduction intégrale, réponses en anglais.
\end{itemize}
\end{aretenir}

\coursec{Comprehension practice}

Dans la section précédente, tu as appris à \cle{lire stratégiquement} : survoler le texte avec le \eng{skimming}, repérer une information précise avec le \eng{scanning}, et deviner le sens des mots inconnus grâce au contexte. Ces stratégies n'ont de valeur que si elles te permettent de \textbf{répondre correctement aux questions}. C'est exactement l'objet de cette section : transformer ta lecture en réponses précises, justifiées et rédigées en anglais correct — comme à l'examen du Baccalauréat.

\courssub{The reading text (reminder)}

Pour travailler, relisons le texte support de la leçon. Lis-le une première fois rapidement (\eng{skimming}), puis une deuxième fois lentement en soulignant les informations clés.

\begin{textebox}[Weekend screens: how one Bamenda neighbourhood relaxes with ICTs]
In the small neighbourhood of Nkwen, in Bamenda, the weekend no longer means only football and church choirs. Information and Communication Technologies have quietly entered homes and changed the way people spend their free time.

On Saturday evenings, whole families gather around a laptop to watch Nigerian films downloaded during the week. “My children prefer comedies, but my wife and I love the old classics,” says Mr. Tanyi, a retired teacher. “It is cheaper than going to a cinema, and we can pause the film to discuss it together.”

Teenagers, however, often prefer their own screens. Sixteen-year-old Blessing spends about two hours every evening on her smartphone, chatting with friends, watching music videos and playing online games. Her mother worries that she spends too much time alone, but Blessing disagrees: “I am not alone! I talk to my cousins in Douala every day. Without my phone, we would never speak.”

The community has also found collective uses for technology. At the Nkwen cybercafé, young people who cannot afford a computer at home pay a small fee to play video games, learn graphic design or follow football matches online. Every Sunday afternoon, the owner organises a free “digital hour” where he teaches elderly people how to use WhatsApp to see photos of their grandchildren abroad.

Not everyone is satisfied, though. Some parents complain that their children neglect their homework, and health workers warn against sitting in front of screens for too long. Experts advise balance: enjoy technology, they say, but keep time for sport, reading and face-to-face conversation.

For Mr. Tanyi, the conclusion is simple: “A screen is like fire. It can warm the house or burn it down. It depends on how you use it.”
\end{textebox}

\textbf{Glossaire}

\begin{center}
\begin{tabularx}{\linewidth}{|l|l|X|}
\hline
\rowcolor{popblueL} \textbf{English} & \textbf{Nature} & \textbf{Français} \\
\hline
to gather & verb & se rassembler, se réunir \\
\hline
retired & adjective & retraité, à la retraite \\
\hline
fee & noun & droit, tarif (somme payée) \\
\hline
elderly people & noun phrase & les personnes âgées \\
\hline
abroad & adverb & à l'étranger \\
\hline
to neglect & verb & négliger \\
\hline
to warn against & verb & mettre en garde contre \\
\hline
balance & noun & équilibre, juste milieu \\
\hline
\end{tabularx}
\end{center}

\courssub{True or False questions: justify every answer}

Le premier type de question que tu rencontreras presque toujours est le \cle{vrai/faux} (\eng{True or False}). La règle d'or : une réponse \eng{True} ou \eng{False} sans justification ne vaut presque rien. Tu dois citer la phrase du texte qui prouve ta réponse.

\begin{exemplebox}[Modèle de réponse justifiée]
\textbf{Question :} True or False? \eng{Families in the text watch Nigerian films together.}

\textbf{Réponse attendue :} \eng{True. The text says: “whole families gather around a laptop to watch Nigerian films.”}

Traduction : Vrai. Le texte dit : « des familles entières se rassemblent autour d'un ordinateur portable pour regarder des films nigérians ».
\end{exemplebox}

\begin{methode}[Réussir un True/False en 4 étapes]
\begin{enumerate}
\item Lis la phrase proposée et souligne ses mots-clés (personnes, actions, lieux, durées).
\item Retrouve dans le texte le passage qui parle de la même chose (\eng{scanning}).
\item Compare mot à mot : le texte \textbf{confirme}-t-il la phrase (\eng{True}) ou la \textbf{contredit}-il (\eng{False}) ?
\item Rédige : \eng{True/False. The text says: “...”} en recopiant la ligne exacte qui justifie ta réponse.
\end{enumerate}
\end{methode}

\begin{attention}[False ne veut pas dire “not mentioned”]
\eng{False} signifie que le texte \textbf{dit le contraire} de la phrase proposée. Si le texte ne dit \emph{rien} sur le sujet, la réponse n'est pas \eng{False} mais \eng{Not stated} (« non mentionné ») — quand cette option existe. Exemple : \eng{Mr. Tanyi has three children.} Le texte ne donne pas le nombre de ses enfants : ce n'est ni vrai ni faux, c'est \eng{Not stated}. Confondre \eng{False} et \eng{Not stated} est une erreur classique des candidats francophones.
\end{attention}

\begin{center}
\begin{tabularx}{\linewidth}{|l|X|X|}
\hline
\rowcolor{popblueL} \textbf{Option} & \textbf{Signification} & \textbf{Quand la choisir} \\
\hline
True & Vrai & Le texte confirme la phrase. \\
\hline
False & Faux & Le texte contredit la phrase. \\
\hline
Not stated & Non mentionné & Le texte ne dit rien à ce sujet. \\
\hline
\end{tabularx}
\end{center}

\courssub{Multiple choice questions (MCQ): main idea and attitudes}

Le \cle{QCM} (\eng{multiple choice question}) teste deux choses : ta compréhension de l'\textbf{idée principale} (\eng{main idea}) et ta capacité à repérer l'\textbf{attitude} ou le sentiment d'un personnage.

\begin{propriete}[Stratégie pour le QCM]
\begin{itemize}
\item Pour l'idée principale : demande-toi « de quoi parle le texte \emph{en entier} ? ». Élimine les options qui ne concernent qu'un seul paragraphe (ce sont des détails, pas l'idée principale).
\item Pour l'attitude d'un personnage : cherche les verbes de sentiment (\eng{worries, disagrees, complain, prefer}) et les adjectifs (\eng{satisfied, cheaper}), puis les citations directes du personnage.
\item Lis les quatre options \textbf{avant} de retourner dans le texte : tu sauras ainsi quoi chercher.
\end{itemize}
\end{propriete}

\begin{exemplebox}[Exemple de QCM sur l'attitude]
\textbf{Question :} \eng{How does Blessing's mother feel about her daughter's phone use?}

\eng{A) She is proud of her. \quad B) She is worried about her. \quad C) She is angry with her. \quad D) She is indifferent.}

\textbf{Réponse :} B. Justification : \eng{“Her mother worries that she spends too much time alone.”} — le verbe \eng{worries} (s'inquiète) indique l'inquiétude, pas la colère ni la fierté.
\end{exemplebox}

\courssub{Short answer questions: who, where, how long, why}

Les questions à réponse courte commencent par un mot interrogatif. Chaque mot interrogatif appelle un type précis de réponse :

\begin{center}
\begin{tabularx}{\linewidth}{|l|l|X|}
\hline
\rowcolor{popblueL} \textbf{Mot interrogatif} & \textbf{Type de réponse} & \textbf{Exemple tiré du texte} \\
\hline
Who...? & une personne & \eng{Mr. Tanyi, a retired teacher.} \\
\hline
Where...? & un lieu & \eng{In Nkwen, in Bamenda.} \\
\hline
When...? / How long...? & un moment / une durée & \eng{About two hours every evening.} \\
\hline
Why...? & une raison (\eng{because}) & \eng{Because it is cheaper than going to a cinema.} \\
\hline
What...? & une chose, une action & \eng{A free “digital hour” for elderly people.} \\
\hline
How...? & une manière, un moyen & \eng{By paying a small fee at the cybercafé.} \\
\hline
\end{tabularx}
\end{center}

\begin{attention}[Réponds par une phrase complète]
À l'examen, une réponse d'un seul mot (\eng{Bamenda.}) est souvent acceptée, mais la phrase complète est toujours plus sûre : \eng{The story takes place in Nkwen, a neighbourhood of Bamenda.} Attention aussi à l'ordre des mots : en anglais, on ne dit jamais \eng{Why she uses her phone?} mais \eng{Why does she use her phone?} — l'auxiliaire \eng{does} est obligatoire dans la question.
\end{attention}

\courssub{Fact or opinion? Interpreting the text}

Dernière compétence : distinguer un \cle{fait} (\eng{fact}) — quelque chose de vérifiable — d'une \cle{opinion} (\eng{opinion}) — un jugement, un conseil ou un sentiment personnel.

\begin{definition}[Fact vs opinion]
Un \eng{fact} peut être prouvé ou vérifié : \eng{Blessing spends about two hours every evening on her smartphone.} (durée mesurable).

Une \eng{opinion} exprime un point de vue, souvent introduite par des verbes comme \eng{advise, think, believe, should} ou des comparaisons subjectives : \eng{Experts advise balance.} (c'est un conseil, pas un fait mesurable).
\end{definition}

\begin{exemplebox}[Application au texte]
\begin{itemize}
\item \eng{“The owner organises a free digital hour every Sunday afternoon.”} → \textbf{Fact} (événement vérifiable).
\item \eng{“Experts advise balance.”} → \textbf{Opinion} (c'est un conseil, une recommandation).
\item \eng{“A screen is like fire. It can warm the house or burn it down.”} → \textbf{Opinion} (comparaison personnelle de Mr. Tanyi).
\end{itemize}
\end{exemplebox}

\begin{popfigure}
\begin{center}
\begin{tikzpicture}[mindmap, grow cyclic,
  every node/.style={concept, font=\small},
  root concept/.append style={concept color=popblue, font=\bfseries},
  level 1 concept/.append style={level distance=3.2cm, sibling angle=120},
  level 2 concept/.append style={level distance=2.6cm, sibling angle=60}]
\node[root concept]{Comprehension questions}
  child[concept color=poporange] { node {True / False} 
    child[concept color=poporangeL] { node[align=center] {justify with\\ a quotation} } }
  child[concept color=popgreen] { node {MCQ}
    child[concept color=popgreenL] { node[align=center] {main idea\\ or attitude} } }
  child[concept color=poppurple] { node[align=center] {Short answers\\ Fact or opinion}
    child[concept color=poppurpleL] { node[align=center] {who, where,\\ why, how long} } };
\end{tikzpicture}
\end{center}
\legende{Carte mentale : les quatre types de questions de compréhension et leur exigence principale.}
\end{popfigure}

\courssub{Guided practice: solved exercise}

\begin{exoresolu}[Compréhension complète du texte]
\textbf{Énoncé.} Read the text again and answer the following questions.

\textbf{A. True or False? Justify with a line from the text.}
\begin{enumerate}
\item \eng{Mr. Tanyi's family goes to the cinema every Saturday.}
\item \eng{Blessing talks to her cousins in Douala thanks to her phone.}
\item \eng{The “digital hour” at the cybercafé is expensive.}
\end{enumerate}

\textbf{B. Choose the correct answer.} The main idea of the text is:
\eng{A) Nigerian films are the best entertainment. B) ICTs have changed leisure in a Bamenda neighbourhood, with advantages and risks. C) Teenagers should not use smartphones. D) Cybercafés are too expensive.}

\textbf{C. Short answers.}
\begin{enumerate}
\item \eng{Who teaches elderly people how to use WhatsApp?}
\item \eng{How long does Blessing spend on her phone every evening?}
\item \eng{Why does Mr. Tanyi prefer watching films at home?}
\end{enumerate}

\textbf{D. Fact or opinion?} \eng{“Experts advise balance.”}

\tcblower

\textbf{Solution détaillée.}

\textbf{A. True or False.}
\begin{enumerate}
\item \eng{False. The text says: “whole families gather around a laptop to watch Nigerian films” and “It is cheaper than going to a cinema.”} (La famille regarde les films à la maison, pas au cinéma.)
\item \eng{True. Blessing says: “I talk to my cousins in Douala every day. Without my phone, we would never speak.”}
\item \eng{False. The text says the owner organises “a free digital hour”.} (\eng{Free} = gratuit ; c'est le contraire de « cher ».)
\end{enumerate}

\textbf{B. MCQ.} Réponse : \textbf{B}. Le texte entier décrit les nouveaux usages des TIC pour les loisirs (films, téléphone, cybercafé) et mentionne à la fois les avantages (communication, apprentissage) et les risques (devoirs négligés, santé). Les options A, C et D ne concernent qu'un détail du texte.

\textbf{C. Short answers.}
\begin{enumerate}
\item \eng{The owner of the Nkwen cybercafé teaches elderly people how to use WhatsApp.}
\item \eng{She spends about two hours every evening on her smartphone.}
\item \eng{Because it is cheaper than going to a cinema, and the family can pause the film to discuss it together.}
\end{enumerate}

\textbf{D. Fact or opinion?} C'est une \textbf{opinion} : \eng{advise} exprime un conseil, une recommandation, et non un fait vérifiable.
\end{exoresolu}

\courssub{Collective correction: the golden sentence}

Lors de la correction collective en classe, chaque élève doit justifier sa réponse à voix haute. Une seule formule à retenir, valable pour tous les types de questions :

\begin{aretenir}[La phrase magique de la justification]
\eng{I know this because the text says: “...”} — « Je le sais parce que le texte dit : ... »

Variantes utiles :
\begin{itemize}
\item \eng{It is written in paragraph two that...} (il est écrit au paragraphe 2 que...)
\item \eng{According to the text, ...} (selon le texte, ...)
\item \eng{Mr. Tanyi says that...} (Mr. Tanyi dit que...)
\end{itemize}
\end{aretenir}

\begin{experience}[Correction collective en classe]
Le professeur interroge un élève par question. L'élève répond en deux temps obligatoires :
\begin{enumerate}
\item Sa réponse : \eng{Number one is false.}
\item Sa justification : \eng{I know this because the text says: “It is cheaper than going to a cinema.”}
\end{enumerate}
Les autres élèves écoutent et peuvent contester : \eng{I disagree, because the text says...} (« je ne suis pas d'accord, parce que le texte dit... »). Celui qui cite la bonne ligne du texte gagne le débat.
\end{experience}

\begin{savaistu}[La compréhension écrite au Bac anglais Tle A]
Au Baccalauréat anglais de la série A, la section \eng{Reading comprehension} vaut environ 10 points de l'épreuve. Les correcteurs évaluent \textbf{systématiquement la justification} : une réponse \eng{True} ou \eng{False} sans la phrase du texte qui la prouve perd la moitié des points, voire la totalité. Entraîne-toi donc dès maintenant à ne jamais répondre sans citer le texte — c'est l'habitude qui fait la différence entre une note moyenne et une excellente note.
\end{savaistu}

\begin{exercice}[À toi de jouer]
Réponds aux questions suivantes en rédigeant des phrases complètes, avec justification pour les vrai/faux.
\begin{enumerate}
\item True or False? \eng{Blessing's mother thinks her daughter spends too much time alone.}
\item True or False? \eng{Mr. Tanyi is still working as a teacher.}
\item Where do young people go if they cannot afford a computer at home?
\item Why do some parents complain?
\item Fact or opinion? \eng{“Health workers warn against sitting in front of screens for too long.”}
\end{enumerate}
\end{exercice}

\begin{corrige}[Exercice « À toi de jouer »]
\begin{enumerate}
\item \eng{True. The text says: “Her mother worries that she spends too much time alone.”}
\item \eng{False. The text says Mr. Tanyi is “a retired teacher”.} (retraité, donc il ne travaille plus)
\item \eng{They go to the Nkwen cybercafé, where they pay a small fee to use a computer.}
\item \eng{Some parents complain that their children neglect their homework.}
\item Opinion (ou conseil) : \eng{warn against} exprime une mise en garde, une recommandation, pas un fait mesurable.
\end{enumerate}
\end{corrige}

\coursec{Thematic vocabulary : ICTs and digital leisure}

Après avoir lu le texte et répondu aux questions de compréhension, nous allons maintenant \cle{organiser le vocabulaire} qui permet de parler des technologies de l'information et de la communication (ICTs) comme ressources de loisir. Un bon élève de Terminale ne mémorise pas des listes isolées : il classe les mots par \cle{champs lexicaux}, il apprend les \cle{collocations} (les mots qui vont ensemble) et il se méfie des \cle{faux amis}. C'est exactement ce que nous allons faire ici, en quatre champs : les appareils, les activités, les lieux d'accès communautaire, et enfin les avantages et les dangers.

\courssub{Field 1 : Devices and tools (les appareils et les outils)}

\begin{definition}[ICTs]
\eng{ICTs} est l'abréviation de \eng{Information and Communication Technologies} (les technologies de l'information et de la communication). Le mot \eng{technology} se prononce « tek-no-lo-dji » et \eng{device} (appareil) se prononce « di-vaïss ».
\end{definition}

Voici les mots essentiels pour nommer les appareils et outils numériques. Remarquez que l'anglais utilise souvent des mots composés (\eng{smartphone}, \eng{keyboard}) ou des noms suivis d'un complément (\eng{memory card}).

\begin{center}
\begin{tabularx}{\linewidth}{|X|X|X|}\hline
\rowcolor{popblueL} English & Français & Exemple en contexte \\ \hline
a smartphone & un smartphone & \eng{Most students in Yaoundé own a smartphone.} \\ \hline
a laptop / a computer & un ordinateur portable / un ordinateur & \eng{She does her homework on her laptop.} \\ \hline
a tablet & une tablette & \eng{The children watch cartoons on a tablet.} \\ \hline
a television set & un poste de télévision & \eng{The family gathers around the television set.} \\ \hline
a game console & une console de jeux & \eng{He received a game console for Christmas.} \\ \hline
a headset / headphones & un casque (audio) & \eng{Put on your headphones to listen quietly.} \\ \hline
a charger & un chargeur & \eng{My charger is broken, so my phone is off.} \\ \hline
a memory card & une carte mémoire & \eng{The memory card is full of photos.} \\ \hline
an app / an application & une application & \eng{WhatsApp is a very popular app in Cameroon.} \\ \hline
a data bundle & un forfait internet & \eng{I bought a data bundle of 2 GB yesterday.} \\ \hline
\end{tabularx}
\end{center}

\begin{attention}[Faux ami : « data »]
\eng{Data} se prononce « déï-ta » et signifie « les données », pas « la date ». \eng{Mobile data} = les données mobiles (la connexion internet du téléphone). Pour « la date », on dit \eng{the date}. Exemple : \eng{I have no data left.} « Je n'ai plus de forfait internet. »
\end{attention}

\courssub{Field 2 : Online leisure activities (les activités de loisir numérique)}

\begin{exemplebox}[Observer avant de mémoriser]
Observez ces phrases tirées de situations réelles. Quel verbe accompagne chaque activité ?

\eng{After school, Divine goes online and streams music while chatting with friends. Her brother prefers to download games and play online. On Saturdays, they watch films together.}

« Après l'école, Divine se connecte et écoute de la musique en streaming tout en discutant avec des amis. Son frère préfère télécharger des jeux et jouer en ligne. Le samedi, ils regardent des films ensemble. »

On observe que chaque activité a \emph{son} verbe : \eng{to stream music}, \eng{to chat with friends}, \eng{to download games}, \eng{to play online}, \eng{to watch films}. Ce sont des \cle{collocations} : il faut les apprendre comme des blocs.
\end{exemplebox}

\begin{center}
\begin{tabularx}{\linewidth}{|X|X|}\hline
\rowcolor{popblueL} Collocation anglaise & Traduction française \\ \hline
to go online & se connecter à internet, aller en ligne \\ \hline
to spend time online & passer du temps en ligne \\ \hline
to stream music / videos & écouter de la musique / regarder des vidéos en streaming \\ \hline
to download an app / a game & télécharger une application / un jeu \\ \hline
to upload photos & mettre des photos en ligne \\ \hline
to chat with friends & discuter (par messages) avec des amis \\ \hline
to play video games / to game & jouer à des jeux vidéo \\ \hline
to browse the internet & naviguer sur internet \\ \hline
to watch a film / a match & regarder un film / un match \\ \hline
to listen to the radio & écouter la radio \\ \hline
to post on social media & publier sur les réseaux sociaux \\ \hline
to share a video & partager une vidéo \\ \hline
\end{tabularx}
\end{center}

\begin{attention}[Piège de traduction : « écouter » et « regarder »]
En français, on « écoute la radio » et on « regarde la télé ». En anglais, attention aux verbes :
\begin{itemize}
\item \eng{to listen TO the radio} (la préposition \eng{to} est obligatoire) : \eng{I listen to music every evening.} « J'écoute de la musique chaque soir. »
\item \eng{to watch TV / a film} (sans préposition) : \eng{We watched a film last night.} « Nous avons regardé un film hier soir. »
\end{itemize}
Ne dites jamais \eng{to listen music} ni \eng{to watch to TV}.
\end{attention}

\begin{definition}[Screen time]
\eng{Screen time} désigne \cle{le temps passé devant un écran} (téléphone, ordinateur, tablette, télévision). C'est un nom indénombrable : on ne dit pas \eng{a screen time}. Exemple : \eng{Doctors advise teenagers to reduce their screen time.} « Les médecins conseillent aux adolescents de réduire leur temps d'écran. »
\end{definition}

\courssub{Field 3 : Community access (les lieux d'accès communautaire)}

Tout le monde ne possède pas d'appareil personnel. Dans beaucoup de quartiers et de villages camerounais, l'accès aux TIC se fait dans des \cle{lieux communautaires}. Voici le vocabulaire pour en parler.

\begin{center}
\begin{tabularx}{\linewidth}{|X|X|X|}\hline
\rowcolor{popblueL} English & Français & Remarque \\ \hline
a cybercafé & un cybercafé & Lieu payant où l'on utilise des ordinateurs connectés. \\ \hline
a viewing centre & un centre de visionnage & Salle payante où l'on regarde des matchs ou des films, très courant au Cameroun. \\ \hline
a community radio station & une radio communautaire & Radio locale qui informe et divertit le quartier ou le village. \\ \hline
a library & une bibliothèque & Lieu où l'on emprunte ou consulte des livres. \\ \hline
a youth centre & un foyer des jeunes & Centre d'activités pour les jeunes. \\ \hline
a game centre & une salle de jeux & Lieu où l'on joue à des jeux vidéo contre paiement. \\ \hline
free Wi-Fi & le Wi-Fi gratuit & \eng{The town hall offers free Wi-Fi in the square.} \\ \hline
\end{tabularx}
\end{center}

\begin{attention}[Faux ami : « a library » n'est pas « une librairie »]
\eng{A library} = une \cle{bibliothèque} (on y lit et on y emprunte des livres gratuitement). Une \emph{librairie} (où l'on achète des livres) se dit \eng{a bookshop} ou \eng{a bookstore}. Exemple : \eng{I borrowed this novel from the school library.} « J'ai emprunté ce roman à la bibliothèque de l'école. » — \eng{I bought this dictionary at a bookshop in Douala.} « J'ai acheté ce dictionnaire dans une librairie à Douala. »
\end{attention}

\courssub{Field 4 : Benefits and risks (avantages et dangers)}

Pour donner son opinion (section suivante), il faut pouvoir peser le pour et le contre. Voici les mots-clés des deux camps.

\begin{exemplebox}[Exemples traduits]
\begin{itemize}
\item \eng{to stream music} = écouter de la musique en streaming : \eng{Streaming music helps me relax after exams.} « Écouter de la musique en streaming m'aide à me détendre après les examens. »
\item \eng{to chat with friends} = discuter avec des amis : \eng{She chats with her friends in Bamenda every night.} « Elle discute avec ses amis de Bamenda chaque soir. »
\item \eng{to download a game} = télécharger un jeu : \eng{He downloaded a football game on his phone.} « Il a téléchargé un jeu de football sur son téléphone. »
\item \eng{a cybercafé} = un cybercafé : \eng{The cybercafé near the market is always full.} « Le cybercafé près du marché est toujours plein. »
\item \eng{a viewing centre} = un centre de visionnage : \eng{Fans went to the viewing centre to watch the Indomitable Lions.} « Les supporters sont allés au centre de visionnage pour regarder les Lions Indomptables. »
\item \eng{addiction} = la dépendance : \eng{Video game addiction is a real problem.} « La dépendance aux jeux vidéo est un vrai problème. »
\item \eng{a data bundle} = un forfait internet : \eng{This data bundle costs 500 francs.} « Ce forfait internet coûte 500 francs. »
\end{itemize}
\end{exemplebox}

\begin{popfigure}
\begin{center}
\begin{tikzpicture}[node distance=0.4cm]
\node[draw=popgreen, fill=popgreenL, rounded corners, thick, minimum width=6.2cm, align=center, font=\bfseries] (ben) {BENEFITS};
\node[draw=popgreen, fill=white, rounded corners, minimum width=6.2cm, align=left, below=of ben, text width=5.8cm] (benlist) {\textbullet\ relaxation\\ \textbullet\ learning\\ \textbullet\ creativity\\ \textbullet\ communication};
\node[draw=poppink, fill=poppinkL, rounded corners, thick, minimum width=6.2cm, align=center, font=\bfseries, right=1.2cm of ben] (risk) {RISKS};
\node[draw=poppink, fill=white, rounded corners, minimum width=6.2cm, align=left, below=of risk, text width=5.8cm] (risklist) {\textbullet\ addiction\\ \textbullet\ less sleep\\ \textbullet\ isolation\\ \textbullet\ poor grades};
\draw[-{Stealth}, thick, popmuted] (ben.east) -- (risk.west) node[midway, above, font=\itshape\small, align=center] {balance};
\end{tikzpicture}
\end{center}
\legende{Les deux visages des loisirs numériques : avantages et risques à équilibrer.}
\end{popfigure}

\begin{propriete}[Constructions à connaître : addicted to, spend time, good/bad at]
Trois structures reviennent constamment quand on parle de loisirs numériques :
\begin{enumerate}
\item \cle{be addicted to + nom ou -ing} = être accro à : \eng{He is addicted to gaming.} « Il est accro aux jeux. » / \eng{She is addicted to her phone.} « Elle est accro à son téléphone. » Attention : la préposition est \eng{to}, jamais \eng{with}.
\item \cle{spend time + -ing} = passer du temps à faire : \eng{They spend three hours playing online every day.} « Ils passent trois heures à jouer en ligne chaque jour. » Le verbe qui suit \eng{spend time} est toujours au gérondif.
\item \cle{be good at / be bad at + -ing} = être bon / mauvais en : \eng{My sister is good at video editing.} « Ma sœur est douée en montage vidéo. » / \eng{I am bad at online chess.} « Je suis mauvais aux échecs en ligne. »
\end{enumerate}
\end{propriete}

\begin{attention}[Faux amis : « actually » et « sensible »]
\begin{itemize}
\item \eng{Actually} ne signifie pas « actuellement » mais « en fait, en réalité ». \eng{I am actually gaming.} = « Je suis en train de jouer, en fait. » Pour dire « actuellement », utilisez \eng{currently} ou \eng{at the moment} : \eng{I am currently watching a film.} « Je regarde un film actuellement. »
\item \eng{Sensible} ne signifie pas « sensible » mais « raisonnable, sensé ». \eng{It is sensible to limit your screen time.} « Il est raisonnable de limiter son temps d'écran. » Pour « sensible » (émotif), dites \eng{sensitive} : \eng{She is sensitive to criticism.} « Elle est sensible à la critique. »
\end{itemize}
\end{attention}

\courssub{Reading : un texte pour réemployer le vocabulaire}

\begin{textebox}[Weekend screens in Buea]
Every weekend, the Njoh family in Buea faces the same question: how much screen time is too much? Sixteen-year-old Claudia loves her smartphone. She goes online as soon as she wakes up, streams music while doing her chores, and chats with friends until midnight. Her younger brother Tabi is addicted to gaming; he spends hours playing football video games and often downloads new apps without asking his parents.

Their father, Mr Njoh, does not own a computer, but he is not cut off from the digital world. On Sunday afternoons, he walks to the viewing centre near the market to watch Premier League matches with other fans. Their mother prefers the community radio station, which broadcasts music, news and health advice in the local language.

Mr Njoh sees both sides of the debate. On the one hand, ICTs are wonderful recreational resources: Claudia has learned video editing online, and Tabi's English has improved because he plays with gamers from Nigeria and Ghana. On the other hand, he worries about the risks. Tabi sleeps less, his grades are falling, and the family hardly talks at dinner any more.

Last month, the family held a meeting. They agreed on simple rules: no phones during meals, homework before games, and one screen-free afternoon every week. “Technology is a good servant but a bad master,” Mr Njoh concluded. The children nodded — and then Claudia posted a photo of the family meeting on social media.

\textbf{Glossary}: screen time (n.) = temps d'écran ; chores (n.) = tâches ménagères ; addicted to (adj.) = accro à ; viewing centre (n.) = centre de visionnage ; to broadcast (v.) = diffuser ; on the one hand / on the other hand = d'une part / d'autre part ; grades (n.) = notes ; hardly (adv.) = à peine ; screen-free (adj.) = sans écran.

\textbf{Questions}: 1) What does Claudia do online? 2) Why is Mr Njoh worried about Tabi? 3) Where does Mr Njoh watch football matches? 4) Mention two benefits of ICTs given in the text. 5) What rules did the family agree on?
\end{textebox}

\begin{exoresolu}[Complétez avec le mot ou la structure correcte]
Complétez les phrases avec : \eng{addicted to}, \eng{spend}, \eng{actually}, \eng{library}, \eng{screen time}, \eng{streaming}.

1. Teenagers should reduce their \ldots\ before bedtime.
2. My cousin is \ldots\ online games; he plays all night.
3. I \ldots\ two hours chatting with friends every evening.
4. I borrowed this book from the town \ldots .
5. She is \ldots\ a very good dancer! (en fait)
6. \ldots\ music uses a lot of mobile data.

\tcblower
\textbf{Solution détaillée} :
1. \eng{screen time} — indénombrable, désigne le temps passé devant un écran. « Les adolescents devraient réduire leur temps d'écran avant le coucher. »
2. \eng{addicted to} — structure \eng{be addicted to + nom}. « Mon cousin est accro aux jeux en ligne ; il joue toute la nuit. »
3. \eng{spend} — structure \eng{spend time + -ing}, d'où \eng{chatting} ensuite. « Je passe deux heures à discuter avec des amis chaque soir. »
4. \eng{library} — bibliothèque (faux ami : une librairie = \eng{a bookshop}). « J'ai emprunté ce livre à la bibliothèque municipale. »
5. \eng{actually} — ici « en fait », et non « actuellement » (qui serait \eng{currently}). « C'est en fait une très bonne danseuse ! »
6. \eng{Streaming} — gérondif en tête de phrase comme sujet. « Écouter de la musique en streaming consomme beaucoup de données mobiles. »
\end{exoresolu}

\begin{exercice}[À vous de jouer]
Traduisez en anglais : 1) « Il passe trop de temps à jouer en ligne. » 2) « Actuellement, elle télécharge une application. » 3) « Le cybercafé près de l'école offre le Wi-Fi gratuit. » 4) « Ma sœur est douée pour le montage vidéo. » 5) « Regarder des matchs au centre de visionnage est un loisir populaire. »
\end{exercice}

\begin{corrige}[Exercice « À vous de jouer »]
1. \eng{He spends too much time playing online.} (structure \eng{spend time + -ing}). 2. \eng{She is currently downloading an app.} (\eng{currently} = actuellement ; présent en \eng{be + -ing}). 3. \eng{The cybercafé near the school offers free Wi-Fi.} 4. \eng{My sister is good at video editing.} (\eng{be good at + -ing}). 5. \eng{Watching matches at the viewing centre is a popular leisure activity.} (gérondif sujet).
\end{corrige}

\begin{aretenir}[L'essentiel du vocabulaire]
\begin{itemize}
\item Appareils : \eng{smartphone, laptop, tablet, game console, data bundle}.
\item Activités : \eng{to go online, to stream music, to download an app, to chat with friends, to spend time online}.
\item Lieux : \eng{a cybercafé, a viewing centre, a community radio station, a library}.
\item Opinion : \eng{benefits} (relaxation, learning, creativity, communication) vs \eng{risks} (addiction, less sleep, isolation, poor grades).
\item Structures : \eng{be addicted to + -ing}, \eng{spend time + -ing}, \eng{be good/bad at + -ing}.
\item Faux amis : \eng{actually} = en fait ; \eng{currently} = actuellement ; \eng{a library} = une bibliothèque ; \eng{sensible} = raisonnable.
\end{itemize}
\end{aretenir}

\coursec{Guided production : summary and opinion}

Dans la section précédente, \eng{Thematic vocabulary : ICTs and digital leisure}, tu as enrichi ton lexique : \eng{screen time}, \eng{streaming}, \eng{online gaming}, \eng{social media}, \eng{to relax}, \eng{addiction}... Il est temps maintenant de \cle{produire} : utiliser ce vocabulaire et les idées du texte pour écrire un court paragraphe personnel. C'est exactement ce que l'on te demandera à l'examen : comprendre un texte, puis \eng{give your opinion} en quelques phrases claires et correctes.

\courssub{Objectif de la section}

À la fin de cette section, tu seras capable de :
\begin{itemize}
\item exprimer ton opinion avec des expressions correctes (\eng{In my opinion}, \eng{I agree}, \eng{On the one hand... on the other hand}) ;
\item rédiger un \cle{résumé guidé} de 5 à 6 phrases selon un plan précis ;
\item relire et corriger la production d'un camarade à l'aide d'une grille simple.
\end{itemize}

\courssub{Observation : deux opinions modèles}

Avant d'énoncer les règles, lis attentivement ces deux extraits de copies d'élèves de Terminale. Observe les expressions en gras : que remarques-tu ?

\begin{textebox}[Two students' opinions]
\textbf{Student A (Ngo Bassa, Lycée de Nkolfoulou):} “\textbf{In my opinion}, ICTs are wonderful recreational resources. \textbf{On the one hand}, they help us relax after school: we can listen to music, watch films and chat with friends. \textbf{On the other hand}, too much gaming can harm our studies. \textbf{I think that} we should limit our screen time.”

\textbf{Student B (Etoudi, Lycée Bilingue de Yaoundé):} “\textbf{I agree that} smartphones are useful for entertainment in our community. \textbf{However}, \textbf{I disagree with} the idea that they are always good. \textbf{In my view}, young people should balance digital leisure and sport.”
\end{textebox}

Tu l'as remarqué : chaque élève annonce clairement sa position grâce à une \cle{expression d'opinion} placée en tête de phrase, puis nuance son propos avec un connecteur d'opposition. C'est cette mécanique que nous allons maintenant systématiser.

\begin{propriete}[Les expressions d'opinion : forme et emploi]
Pour donner ton avis en anglais, utilise ces structures fixes :
\begin{itemize}
\item \eng{In my opinion / In my view, + phrase} = « À mon avis, ... » (en tête de phrase, suivi d'une virgule).
\item \eng{I think (that) + phrase} = « Je pense que ... ». Le mot \eng{that} est facultatif : \eng{I think (that) ICTs are useful.}
\item \eng{I agree / I disagree (+ with + nom ou + that + phrase)} = « Je suis d'accord / je ne suis pas d'accord ». \textbf{Attention :} \eng{agree} est un \cle{verbe}, jamais un adjectif.
\item \eng{On the one hand, ... On the other hand, ...} = « D'une part, ... d'autre part, ... » : pour présenter les deux faces d'un sujet (avantages et dangers des TIC, par exemple).
\item \eng{However / But} = « Cependant / Mais » : pour introduire une nuance ou une opposition.
\end{itemize}
\end{propriete}

\begin{center}
\begin{tabularx}{\linewidth}{|X|X|X|}
\hline
\rowcolor{popblueL} Français & English & Remarque \\
\hline
À mon avis, les TIC sont utiles. & In my opinion, ICTs are useful. & Virgule obligatoire après l'expression. \\
\hline
Je pense qu'ils nous aident à nous détendre. & I think (that) they help us (to) relax. & \eng{that} facultatif ; \eng{help} + base verbale ou \eng{to}-infinitif. \\
\hline
Je suis d'accord avec cette idée. & I agree with this idea. & Jamais \emph{I am agree} ! \\
\hline
Je ne suis pas d'accord avec toi. & I disagree with you. & Préfixe négatif \eng{dis-}, pas de double négation. \\
\hline
D'une part, ils nous divertissent ; d'autre part, ils peuvent créer une dépendance. & On the one hand, they entertain us; on the other hand, they can be addictive. & Les deux membres vont toujours ensemble. \\
\hline
Cependant, nous devons limiter notre temps d'écran. & However, we must limit our screen time. & \eng{However} + virgule en tête de phrase. \\
\hline
\end{tabularx}
\end{center}

\begin{attention}[L'erreur n° 1 du francophone : « I am agree »]
N'écris \textbf{jamais} \emph{I am agree}. En français, « être d'accord » utilise le verbe \emph{être} ; en anglais, \eng{agree} est un \cle{verbe} à part entière : on dit simplement \eng{I agree}. De même, \emph{I am disagree} est faux — dis \eng{I disagree}. Autres pièges fréquents :
\begin{itemize}
\item \emph{In my opinion I think that...} : redondant ! Choisis \eng{In my opinion, ...} \textbf{ou} \eng{I think that ...}, pas les deux.
\item \emph{according to me} : faux ami de « selon moi ». Dis \eng{In my opinion} ou \eng{In my view}. (\eng{According to} s'emploie seulement avec une source : \eng{According to the text, ...} = « Selon le texte, ... ».)
\item \emph{the TICs} : faux ami ! En anglais, on dit \eng{ICTs} (\eng{Information and Communication Technologies}), et en général sans article : \eng{ICTs are useful}.
\item \emph{to make sport} : faux ami de « faire du sport ». Dis \eng{to play sport} ou \eng{to do sport} : \eng{We should do sport every week.}
\end{itemize}
\end{attention}

\begin{exemplebox}[Exemples traduits]
\eng{In my opinion, ICTs are useful for relaxation, but we must limit our screen time.} = « À mon avis, les TIC sont utiles pour se détendre, mais nous devons limiter notre temps d'écran. »

\eng{On the other hand, too much gaming can harm our studies.} = « D'un autre côté, trop de jeux vidéo peut nuire à nos études. »

\eng{I think that social media help young people stay in touch with their friends.} = « Je pense que les réseaux sociaux aident les jeunes à rester en contact avec leurs amis. »

\eng{I disagree with the idea that television is a waste of time.} = « Je ne suis pas d'accord avec l'idée que la télévision est une perte de temps. »
\end{exemplebox}

\courssub{Le plan du résumé guidé : l'escalier en cinq marches}

Un bon paragraphe d'opinion-résumé suit une progression logique en cinq étapes. Retiens ce plan : c'est ta rampe de sécurité pour ne jamais bloquer devant la feuille blanche.

\begin{popfigure}
\begin{center}
\begin{tikzpicture}[font=\small]
\node[draw=popblue, fill=popblueL, rounded corners, minimum width=4.2cm, minimum height=0.9cm, align=center] (s1) at (0,0) {\textbf{1. General idea}\\ \eng{ICTs are popular recreational resources.}};
\node[draw=popgreen, fill=popgreenL, rounded corners, minimum width=4.2cm, minimum height=0.9cm, align=center] (s2) at (4.8,-1.3) {\textbf{2. Advantages (2 sentences)}\\ \eng{They entertain and educate us.}};
\node[draw=poporange, fill=poporangeL, rounded corners, minimum width=4.2cm, minimum height=0.9cm, align=center] (s3) at (0,-2.6) {\textbf{3. Danger (1 sentence)}\\ \eng{Too much screen time is harmful.}};
\node[draw=poppurple, fill=poppurpleL, rounded corners, minimum width=4.2cm, minimum height=0.9cm, align=center] (s4) at (4.8,-3.9) {\textbf{4. My opinion}\\ \eng{In my opinion, we should...}};
\node[draw=popgold, fill=popgoldL, rounded corners, minimum width=4.2cm, minimum height=0.9cm, align=center] (s5) at (0,-5.2) {\textbf{5. Conclusion}\\ \eng{Let us use ICTs wisely.}};
\draw[-{Stealth}, thick, popblue] (s1.south) -- (s2.north west);
\draw[-{Stealth}, thick, popgreen] (s2.south) -- (s3.north east);
\draw[-{Stealth}, thick, poporange] (s3.south) -- (s4.north west);
\draw[-{Stealth}, thick, poppurple] (s4.south) -- (s5.north east);
\end{tikzpicture}
\end{center}
\legende{Le plan en escalier : General idea $\rightarrow$ Advantages $\rightarrow$ Danger $\rightarrow$ My opinion $\rightarrow$ Conclusion.}
\end{popfigure}

\begin{methode}[Rédiger un bon résumé d'opinion en cinq étapes]
\begin{enumerate}
\item \textbf{Phrase 1 — Idée générale :} reformule le thème du texte avec \emph{tes propres mots}. Exemple : \eng{Nowadays, ICTs are the favourite recreational resources of many young Cameroonians.}
\item \textbf{Phrases 2 et 3 — Deux avantages :} choisis deux bénéfices mentionnés dans le texte (divertissement, information, contact social, détente) et relie-les avec \eng{First}, \eng{Moreover} ou \eng{On the one hand}.
\item \textbf{Phrase 4 — Un danger :} signale un risque (addiction, mauvais résultats scolaires, isolement) avec \eng{However} ou \eng{On the other hand}.
\item \textbf{Phrase 5 — Ta conclusion personnelle :} donne ton avis avec \eng{In my opinion} ou \eng{I think that}, éventuellement avec un conseil (\eng{should}, \eng{must}).
\end{enumerate}
\textbf{Règles d'or :} un bon résumé = \cle{tes propres mots}, \textbf{pas de phrases copiées du texte}, et \textbf{5 à 6 phrases maximum}. Ne traduis pas mot à mot depuis le français : pense directement en phrases anglaises simples.
\end{methode}

\courssub{Un modèle complet : « ICTs as recreational resources : my view »}

Lis ce modèle de production rédigé par une élève de Tle A. Repère les cinq marches de l'escalier et souligne mentalement les expressions d'opinion.

\begin{textebox}[Model paragraph — ICTs as recreational resources: my view]
Nowadays, ICTs are the most popular recreational resources in our homes and communities. First, they entertain us: after a long day at school, we can watch films, listen to music or play online games with friends. Moreover, they educate us, because the Internet gives access to free tutorials, documentaries and e-books. On the other hand, too much screen time can be dangerous: some students become addicted to gaming and neglect their homework. In my opinion, ICTs are useful for relaxation, but we must limit our screen time and practise sport too. To conclude, let us use digital leisure wisely, so that technology remains a friend and not an enemy.

\textbf{Glossary:} nowadays (adv.) = de nos jours ; to entertain (v.) = divertir ; tutorial (n.) = tutoriel ; screen time (n.) = temps d'écran ; addicted (adj.) = dépendant ; to neglect (v.) = négliger ; wisely (adv.) = sagement.
\end{textebox}

Compte les phrases : il y en a exactement six, une par marche de l'escalier (deux pour les avantages). Chaque connecteur a une fonction précise : \eng{First} et \eng{Moreover} ajoutent, \eng{On the other hand} oppose, \eng{In my opinion} engage l'élève, \eng{To conclude} ferme le paragraphe.

\begin{exoresolu}[Transformer des notes en paragraphe d'opinion]
\textbf{Énoncé.} Voici les notes d'un élève après la lecture du texte : \emph{ICTs = fun ; music and films ; danger = addiction ; my view = balance}. Rédige un paragraphe de 5 phrases en suivant le plan en escalier et en utilisant au moins deux expressions d'opinion.
\tcblower
\textbf{Solution détaillée.} Marche 1 (idée générale) : \eng{ICTs are wonderful recreational resources for young people today.} Marche 2 (avantages) : \eng{On the one hand, they allow us to listen to music and watch films at home, which helps us relax.} Marche 3 (danger) : \eng{On the other hand, too much gaming can harm our studies because it creates addiction.} Marche 4 (opinion) : \eng{In my opinion, we should find a balance between digital leisure and other activities like reading or football.} Marche 5 (conclusion) : \eng{I think that a wise use of ICTs makes our free time both enjoyable and useful.} Vérification : 5 phrases, aucune copiée d'un texte, deux expressions d'opinion (\eng{In my opinion}, \eng{I think that}), un connecteur d'opposition (\eng{On the other hand}) et du lexique TIC (\eng{gaming}, \eng{digital leisure}, \eng{addiction}).
\end{exoresolu}

\courssub{Relecture en binôme : la grille simple}

Après ta rédaction individuelle, échange ta copie avec ton voisin. La relecture par les pairs (\eng{peer review}) est une compétence clé : elle t'apprend à repérer les erreurs... y compris les tiennes.

\begin{methode}[Grille de relecture en binôme]
Coche chaque critère sur la copie de ton camarade :
\begin{enumerate}
\item \textbf{Lexique TIC :} au moins trois mots du thème sont utilisés correctement (\eng{screen time}, \eng{online gaming}, \eng{social media}, \eng{to relax}, \eng{addiction}...).
\item \textbf{Connecteurs :} au moins deux connecteurs de logique (\eng{First}, \eng{Moreover}, \eng{However}, \eng{On the one hand / On the other hand}, \eng{To conclude}).
\item \textbf{Opinion claire :} une expression d'opinion correcte est présente (\eng{In my opinion}, \eng{I think that}, \eng{I agree/disagree}) — vérifie qu'il n'y a pas de \emph{I am agree} !
\item \textbf{Plan respecté :} les cinq marches sont dans l'ordre (idée générale $\rightarrow$ avantages $\rightarrow$ danger $\rightarrow$ opinion $\rightarrow$ conclusion).
\item \textbf{Longueur :} 5 à 6 phrases, pas de phrase copiée du texte support.
\end{enumerate}
Termine par un commentaire encourageant : \eng{Well done! Your opinion is very clear.} ou \eng{Good ideas, but check the verb “agree”.}
\end{methode}

\begin{exercice}[À toi de produire]
Rédige individuellement un paragraphe intitulé \eng{ICTs as recreational resources: my view} (5 à 6 phrases) en suivant le plan en escalier. Contraintes : utilise au moins trois mots du lexique TIC, deux connecteurs et une expression d'opinion. Ensuite, échange avec ton binôme et applique la grille de relecture.
\end{exercice}

\begin{corrige}[Exercice — éléments de correction]
Il n'existe pas une seule bonne réponse, mais toute copie réussie contient : (1) une phrase d'idée générale reformulée ; (2) deux avantages reliés par \eng{First} / \eng{Moreover} ; (3) un danger introduit par \eng{However} ou \eng{On the other hand} ; (4) une opinion personnelle correcte (\eng{I agree}, jamais \emph{I am agree}) ; (5) une conclusion brève avec un conseil (\eng{should}, \eng{must}). Exemple de conclusion attendue : \eng{To conclude, ICTs are excellent tools for leisure if we control our screen time.}
\end{corrige}

\begin{experience}[Débat éclair : “Are ICTs good or bad for our free time?”]
En groupes de quatre : deux élèves défendent les avantages des TIC comme ressources de loisir (\eng{for}), deux élèves soulignent les dangers (\eng{against}). Chacun parle 30 secondes en commençant obligatoirement par une expression d'opinion (\eng{In my opinion...}, \eng{I strongly believe that...}). Le groupe vote ensuite et justifie son choix en une phrase écrite au tableau.
\end{experience}

\begin{aretenir}[Complément Bac]
Cette tâche entraîne \textbf{directement} la question d'expression personnelle fréquente en épreuve de compréhension écrite au baccalauréat : \eng{Do you think that ICTs are good recreational resources? Justify your answer.} Stratégie gagnante : réponds en 3 à 5 phrases avec (1) ta position claire (\eng{I think that...}), (2) un argument tiré du texte reformulé, (3) un contre-argument ou une nuance (\eng{However...}), (4) une courte conclusion. C'est exactement l'escalier que tu viens de t'entraîner à gravir.
\end{aretenir}

\begin{savaistu}[Pourquoi “screen time” est entré dans les dictionnaires]
L'expression \eng{screen time} (« temps d'écran ») est apparue dans les années 1990 avec la démocratisation de la télévision et des ordinateurs, puis s'est imposée avec les smartphones. Elle est aujourd'hui utilisée par les médecins et les enseignants du monde entier, y compris au Cameroun, où les débats sur l'usage des téléphones portables en classe sont vifs. Maîtriser ce vocabulaire, c'est pouvoir participer à une conversation véritablement mondiale !
\end{savaistu}

\newpage

\coursec{Activité d'intégration}

\coursec{Lesson 55 — Reading: Texts on using ICTs as recreational resources at home or in the community}

\courssub{Warm-up and pre-reading}

\begin{experience}[Brainstorming : les écrans dans notre quotidien]
L'enseignant projette ou affiche une série de 4 photos : un adolescent sur un smartphone dans un bendskin, une famille devant la télévision au salon, des jeunes dans un cybercafé à Buea, un élève qui révise sur une tablette. En classe entière, on nomme les activités : \eng{chatting on WhatsApp}, \eng{watching streaming videos}, \eng{online gaming}, \eng{studying online}. L'enseignant écrit au tableau le mot-clé \cle{ICTs} (Information and Communication Technologies) et demande : \eng{“What do you do for fun with your phone or computer at home?”}. Les réponses sont notées en \eng{mind map} autour de \eng{DIGITAL LEISURE}.
\tcblower
\begin{tikzpicture}[mindmap, grow cyclic, every node/.style=concept, concept color=popblue, level 1/.append style={level distance=4cm, sibling angle=90}, level 2/.append style={level distance=3cm, sibling angle=45}]
\node {Digital Leisure}
    child { node {Social Media} child { node {WhatsApp} } child { node {TikTok} } child { node {Facebook} } }
    child { node {Gaming} child { node {Mobile games} } child { node {Console/PC} } child { node {e-sport} } }
    child { node {Streaming} child { node {Music (Boomplay)} } child { node {Video (YouTube)} } child { node {Series/Films} } }
    child { node {Creative} child { node {Editing} } child { node {Coding} } child { node {Blogging} } };
\end{tikzpicture}
\legende{Mind map : Digital Leisure activities in Cameroon}
\end{experience}

\begin{aretenir}[Lexique d'entrée : \eng{Digital Leisure}]
\begin{center}
\begin{tabularx}{\linewidth}{|X|X|X|}\hline
\rowcolor{popblueL} \textbf{Français} & \textbf{English} & \textbf{Collocations / Exemples} \\ \hline
Réseaux sociaux & \eng{social media} & \eng{scroll through social media}, \eng{post a story} \\ \hline
Jeux en ligne & \eng{online gaming} & \eng{play online games}, \eng{a gamer} \\ \hline
Temps d'écran & \eng{screen time} & \eng{limit screen time}, \eng{excessive screen time} \\ \hline
Regarder en streaming & \eng{to stream} & \eng{stream a movie}, \eng{a streaming platform} \\ \hline
Faire défiler & \eng{to scroll} & \eng{scroll down}, \eng{doomscrolling} \\ \hline
Rester connecté & \eng{to stay connected} & \eng{stay connected with relatives} \\ \hline
Dépendance & \eng{addiction} & \eng{screen addiction}, \eng{addicted to gaming} \\ \hline
Équilibre numérique & \eng{digital balance} & \eng{find a digital balance}, \eng{digital wellbeing} \\ \hline
Activités hors ligne & \eng{offline activities} & \eng{enjoy offline activities}, \eng{go offline} \\ \hline
\end{tabularx}
\end{center}
\end{aretenir}

\courssub{Reading text with glossary}

\begin{methode}[Première lecture : Skimming (lecture rapide pour l'idée générale)]
\begin{enumerate}
\item Lisez le titre et observez la structure (paragraphes, mots en gras, images éventuelles).
\item Lisez la première et la dernière phrase de chaque paragraphe.
\item Répondez en une phrase : \eng{“What is the main idea of the text?”} (Sujet : avantages et dangers des TIC à la maison ; thèse : l'équilibre est nécessaire).
\end{enumerate}
\end{methode}

\begin{textebox}[“Screens at home: joy or danger?” — Texte support original]
In many Cameroonian homes today, the smartphone has become the family's favourite companion. After school, teenagers in Yaoundé and Douala chat with friends on WhatsApp, watch comedy videos, or play online games with players from Lagos to London. Used wisely, ICTs are wonderful recreational resources: they help young people relax after a long day, discover new music, learn skills through tutorials, and stay connected with relatives who live far away.

However, digital leisure can quickly become a screen trap. Some students spend five or six hours a day scrolling through social media and forget their homework, their sleep and even their meals. Doctors warn that too much screen time causes tired eyes, headaches and poor concentration. Parents also worry about violent games and strangers who contact teenagers online.

The solution is not to throw away our phones, but to use them smartly. Experts advise families to set clear rules: no phones during meals, a maximum of two hours of recreational screen time per day, and at least one hour of outdoor activity. In some communities, youth centres now offer “offline evenings” with board games, football and storytelling — proof that real life can be fun too.
\end{textebox}

\begin{definition}[Glossaire du texte]
\begin{itemize}
\item \eng{wisely} (adv.) : sagement, avec discernement.
\item \eng{recreational resources} (n. pl.) : ressources de loisir / de détente.
\item \eng{to scroll (through)} (v.) : faire défiler (un fil d'actualité).
\item \eng{a screen trap} (n.) : un piège de l'écran (usage excessif et nocif).
\item \eng{to warn} (v.) : avertir, mettre en garde.
\item \eng{outdoor activity} (n.) : activité de plein air.
\item \eng{youth centre} (n.) : centre de jeunes (Maison des jeunes).
\item \eng{board games} (n. pl.) : jeux de société.
\item \eng{storytelling} (n.) : art de conter / raconter des histoires.
\end{itemize}
\end{definition}

\courssub{Reading strategies: skimming, scanning, guessing}

\begin{propriete}[Stratégie 1 : Skimming (Lecture diagonale)]
Objectif : Saisir l'idée générale (\eng{gist}) sans lire chaque mot.
Technique : Lire le titre, l'introduction, la conclusion, la première phrase de chaque paragraphe (\eng{topic sentence}), les mots-clés en gras.
Application : Ce texte a une structure \eng{Problem – Solution} (Paragraphe 1 : avantages, Paragraphe 2 : dangers, Paragraphe 3 : solutions).
\end{propriete}

\begin{propriete}[Stratégie 2 : Scanning (Repérage d'informations précises)]
Objectif : Trouver une information ciblée (chiffre, nom, mot-clé) sans lire tout le texte.
Technique : Définir la question (\eng{How many hours?}, \eng{What do doctors warn about?}). Balayer le texte verticalement pour repérer le mot-clé (\eng{hours}, \eng{doctors}, \eng{warn}). Lire seulement la phrase concernée.
Exercice : Repérez dans le texte : \eng{two hours}, \eng{tired eyes}, \eng{headaches}, \eng{poor concentration}, \eng{violent games}, \eng{offline evenings}.
\end{propriete}

\begin{propriete}[Stratégie 3 : Guessing meaning from context (Déduction du sens)]
Objectif : Comprendre un mot inconnu sans dictionnaire.
Indices : 1. La classe grammaticale (nom, verbe, adjectif). 2. Le contexte immédiat (mots autour). 3. Les connecteurs (\eng{however}, \eng{because}, \eng{for example}). 4. La morphologie (préfixes, suffixes).
Exemple : \eng{“recreational”} (paragraphe 1) $\rightarrow$ contexte \eng{“relax”, “leisure”} $\rightarrow$ adjectif $\rightarrow$ « de loisir / récréatif ». \eng{“scrolling”} (paragraphe 2) $\rightarrow$ verbe \eng{to scroll} + \eng{-ing} $\rightarrow$ action de faire défiler.
\end{propriete}

\begin{exercice}[Entraînement aux stratégies]
\begin{enumerate}
\item \textbf{Skimming :} Choisissez le meilleur résumé du texte. (A) Les smartphones sont dangereux et doivent être interdits. (B) Les TIC offrent des loisirs mais exigent des règles pour éviter l'addiction. (C) Les jeunes Camerounais ne font que jouer en ligne.
\item \textbf{Scanning :} Trouvez dans le texte les équivalents anglais de : « temps d'écran récréatif », « activités de plein air », « jeux de société », « mettre en garde ».
\item \textbf{Guessing :} Déduisez le sens de \eng{companion} (l.1), \eng{tutorials} (l.4), \eng{strangers} (l.10), \eng{smartly} (l.12).
\end{enumerate}
\end{exercice}

\begin{corrige}[Exercice Stratégies]
\begin{enumerate}
\item \textbf{B.} Le texte présente les deux faces (joie/danger) et propose l'équilibre (règles, soirées hors ligne).
\item \textbf{Scanning :} \eng{recreational screen time}, \eng{outdoor activity}, \eng{board games}, \eng{warn}.
\item \textbf{Guessing :} \eng{companion} = compagnon (contexte \eng{favourite}, \eng{chat with friends}) ; \eng{tutorials} = tutoriels (contexte \eng{learn skills}) ; \eng{strangers} = inconnus / étrangers (contexte \eng{contact teenagers online}, danger) ; \eng{smartly} = intelligemment / avec discernement (opposé à \eng{throw away}, contexte \eng{Experts advise}).
\end{enumerate}
\end{corrige}

\courssub{Comprehension practice}

\begin{exercice}[Questions de compréhension détaillée]
\textbf{A. Right or Wrong? Justify by quoting the text.}
\begin{enumerate}
\item ICTs are only dangerous for teenagers.
\item Doctors say screen time causes health problems.
\item The text suggests throwing away smartphones.
\item Youth centres organise activities without screens.
\end{enumerate}

\textbf{B. Answer the following questions in your own words.}
\begin{enumerate}
\item List three benefits of ICTs mentioned in paragraph 1.
\item What are the consequences of excessive screen time according to doctors?
\item What three rules do experts advise families to set?
\item Explain the expression “screen trap” using the text.
\end{enumerate}

\textbf{C. Vocabulary in context.}
Match each word from the text (1-5) with its definition (a-e).
1. \eng{wisely} \quad 2. \eng{scroll} \quad 3. \eng{warn} \quad 4. \eng{outdoor} \quad 5. \eng{recreational}
a) done for enjoyment \quad b) outside a building \quad c) move text on a screen \quad d) tell someone of danger \quad e) showing good judgement
\end{exercice}

\begin{corrige}[Exercice Compréhension]
\textbf{A.}
\begin{enumerate}
\item \textbf{Wrong.} “Used wisely, ICTs are wonderful recreational resources…” (l.3-4).
\item \textbf{Right.} “Doctors warn that too much screen time causes tired eyes, headaches and poor concentration.” (l.9-10).
\item \textbf{Wrong.} “The solution is not to throw away our phones, but to use them smartly.” (l.11-12).
\item \textbf{Right.} “youth centres now offer ‘offline evenings’ with board games, football and storytelling” (l.14-15).
\end{enumerate}

\textbf{B.}
\begin{enumerate}
\item Relax after a long day; discover new music; learn skills through tutorials; stay connected with distant relatives. (Any three).
\item Tired eyes, headaches, poor concentration.
\item No phones during meals; max two hours recreational screen time/day; at least one hour outdoor activity.
\item A situation where digital leisure becomes dangerous/addictive, making one forget homework, sleep, meals (l.7-8).
\end{enumerate}

\textbf{C.} 1-e, 2-c, 3-d, 4-b, 5-a.
\end{corrige}

\courssub{Thematic vocabulary: ICTs and digital leisure}

\begin{propriete}[Word Building : Nouns, Verbs, Adjectives]
\begin{center}
\begin{tabular}{|l|l|l|l|}\hline
\rowcolor{popblueL} \textbf{Verb} & \textbf{Noun (Thing/Person)} & \textbf{Noun (Abstract)} & \textbf{Adjective} \\ \hline
\eng{to connect} & \eng{a connection} & \eng{connectivity} & \eng{connected} / \eng{online} \\ \hline
\eng{to addict} & \eng{an addict} & \eng{addiction} & \eng{addictive} \\ \hline
\eng{to recreate} & \eng{recreation} & — & \eng{recreational} \\ \hline
\eng{to balance} & \eng{a balance} & — & \eng{balanced} \\ \hline
\eng{to advise} & \eng{advice} (U) / \eng{an adviser} & — & \eng{advisable} \\ \hline
\eng{to warn} & \eng{a warning} & — & \eng{warned} / \eng{warning} (adj.) \\ \hline
\end{tabular}
\end{center}
\emph{Note :} \eng{Advice} est indénombrable $\rightarrow$ \eng{a piece of advice}. Verbe : \eng{to advise} (z), Nom : \eng{advice} (s).
\end{propriete}

\begin{exemplebox}[Phrasal Verbs \& Collocations du thème]
\begin{itemize}
\item \eng{log in / log out} (se connecter / se déconnecter) — \eng{sign up for} (s'inscrire à).
\item \eng{scroll through / scroll down} (faire défiler).
\item \eng{switch off / turn off} (éteindre) $\neq$ \eng{switch on / turn on}.
\item \eng{cut down on} (réduire) — \eng{cut down on screen time}.
\item \eng{stay up late} (se coucher tard) — souvent à cause des écrans.
\item \eng{hang out (online)} (traîner / passer du temps).
\end{itemize}
\end{exemplebox}

\begin{exercice}[Consolidation lexicale]
\begin{enumerate}
\item Complétez avec la forme correcte du mot entre parenthèses.
\begin{enumerate}
\item Excessive \eng{\_\_\_\_\_\_\_\_\_\_} (use) of social media can lead to \eng{\_\_\_\_\_\_\_\_\_\_} (addict).
\item Doctors \eng{\_\_\_\_\_\_\_\_\_\_} (warning) against the blue light effect on sleep.
\item We need a \eng{\_\_\_\_\_\_\_\_\_\_} (balance) approach to digital tools.
\item This game is highly \eng{\_\_\_\_\_\_\_\_\_\_} (addict); many teens are \eng{\_\_\_\_\_\_\_\_\_\_} (addict) to it.
\end{enumerate}
\item Remplacez les mots soulignés par un phrasal verb de la liste : \eng{cut down on}, \eng{log out}, \eng{switch off}, \eng{sign up for}.
\begin{enumerate}
\item You should \textbf{reduce} your time on TikTok.
\item Remember to \textbf{disconnect} from your account on public computers.
\item \textbf{Turn off} notifications during class.
\item I want to \textbf{register for} an online coding course.
\end{enumerate}
\end{enumerate}
\end{exercice}

\begin{corrige}[Exercice Vocabulaire]
\begin{enumerate}
\item \eng{use} (noun), \eng{addiction} (noun). \item \eng{warn} (verb). \item \eng{balanced} (adj). \item \eng{addictive} (adj), \eng{addicts} (noun pl) ou \eng{addicted} (adj).
\item \eng{cut down on}, \eng{log out}, \eng{switch off}, \eng{sign up for}.
\end{enumerate}
\end{corrige}

\courssub{Guided production: summary and opinion}

\textbf{Objectifs de l'activité}\par
À la fin de cette activité d'intégration, l'élève sera capable de :
\begin{itemize}
\item relire rapidement un texte sur les loisirs numériques en utilisant la stratégie du \cle{scanning} pour repérer des informations précises (avantages et dangers) ;
\item mobiliser le \cle{lexique des TIC} et des loisirs numériques (\eng{screen time}, \eng{online gaming}, \eng{social media}, \eng{addiction}, \eng{balance}\ldots) dans une production authentique ;
\item exprimer une \cle{opinion} et formuler des \cle{conseils} en anglais à l'aide des structures étudiées (\eng{should}, \eng{it is important to}, impératif) ;
\item rédiger en groupe une \cle{affiche} courte et percutante en anglais correct ;
\item présenter oralement un message en une minute, avec une prononciation claire et un débit adapté ;
\item évaluer les productions des autres groupes et justifier un choix (vote argumenté).
\end{itemize}

\textbf{Situation}\par
Le club d'anglais (\eng{English Club}) de ton lycée à Yaoundé prépare la \eng{Digital Wellness Week}, une semaine de sensibilisation à l'usage équilibré des écrans, ouverte à tous les élèves et à leurs parents. Le principal a remarqué que beaucoup d'élèves passent leurs soirées sur WhatsApp, TikTok ou les jeux en ligne, parfois jusqu'à minuit, et arrivent fatigués en classe. Mais il sait aussi que les TIC peuvent être une formidable source de détente, de créativité et d'apprentissage.

Le club lance donc un concours d'affiches intitulé \eng{“Smart fun or screen trap?”} (« Plaisir intelligent ou piège de l'écran ? »). Chaque groupe doit concevoir une affiche en anglais donnant trois conseils pour un usage équilibré des TIC à la maison. Les meilleures affiches seront affichées dans la salle d'informatique et publiées sur la page Facebook de l'établissement.

Pour préparer ton affiche, relis ce court extrait du texte étudié dans la leçon :

\begin{textebox}[From the lesson text — “Screens at home: joy or danger?”]
In many Cameroonian homes today, the smartphone has become the family's favourite companion. After school, teenagers in Yaoundé and Douala chat with friends on WhatsApp, watch comedy videos, or play online games with players from Lagos to London. Used wisely, ICTs are wonderful recreational resources: they help young people relax after a long day, discover new music, learn skills through tutorials, and stay connected with relatives who live far away.

However, digital leisure can quickly become a screen trap. Some students spend five or six hours a day scrolling through social media and forget their homework, their sleep and even their meals. Doctors warn that too much screen time causes tired eyes, headaches and poor concentration. Parents also worry about violent games and strangers who contact teenagers online.

The solution is not to throw away our phones, but to use them smartly. Experts advise families to set clear rules: no phones during meals, a maximum of two hours of recreational screen time per day, and at least one hour of outdoor activity. In some communities, youth centres now offer “offline evenings” with board games, football and storytelling — proof that real life can be fun too.
\end{textebox}

\textbf{Glossaire :} \emph{wisely} (adv.) : sagement ; \emph{to scroll} (v.) : faire défiler ; \emph{a screen trap} (n.) : un piège de l'écran ; \emph{to warn} (v.) : avertir ; \emph{storytelling} (n.) : art de raconter des histoires.

\textbf{Consignes}\par
Travail en groupes de 4. Répartition des rôles : un \eng{reader} (relit le texte), un \eng{secretary} (écrit), un \eng{timekeeper} (surveille le temps), un \eng{spokesperson} (présente l'affiche).

\begin{enumerate}
\item \textbf{Scanning (5 min).} Relis rapidement le texte ci-dessus et surligne exactement \textbf{deux avantages} (\eng{benefits}) et \textbf{deux dangers} (\eng{dangers}) des loisirs numériques. Recopie-les en expressions courtes : ex. \eng{“helps young people relax”}, \eng{“causes poor concentration”}.

\item \textbf{Brainstorming en groupe (5 min).} À partir de vos quatre éléments, discutez : quelles règles simples votre famille ou votre lycée pourrait-elle appliquer ? Proposez au moins cinq idées de conseils à l'oral, en utilisant \eng{we should\ldots}, \eng{it is important to\ldots}, \eng{why not\ldots?}

\item \textbf{Rédaction de l'affiche (15 min).} Choisissez vos \textbf{trois meilleurs conseils} et rédigez l'affiche. Elle doit contenir obligatoirement :
\begin{itemize}
\item le titre \eng{“Smart fun or screen trap?”} ;
\item un slogan court (ex. \eng{“Enjoy your screen, don't let it enjoy you!”}) ;
\item trois conseils numérotés, à l'\textbf{impératif} ou avec \eng{should}, chacun suivi d'une courte justification introduite par \eng{because} ou \eng{so that} ;
\item au moins \textbf{six mots du lexique des TIC} de la leçon, soulignés sur l'affiche.
\end{itemize}

\item \textbf{Relecture (5 min).} Vérifiez ensemble : ordre des mots, orthographe, présence d'un sujet dans chaque phrase (sauf impératif), ponctuation. Le \eng{secretary} corrige l'affiche.

\item \textbf{Présentation orale (1 min par groupe).} Le \eng{spokesperson} présente l'affiche à la classe : il lit le slogan, explique les trois conseils et dit pourquoi le groupe les a choisis. Les autres élèves prennent des notes.

\item \textbf{Vote argumenté (5 min).} La classe vote pour les trois meilleurs conseils. Chaque électeur justifie son vote en une phrase : \eng{“I vote for tip number two because it is realistic and easy to follow.”} Les trois conseils gagnants seront affichés dans la salle.
\end{enumerate}

\textbf{Ressources à mobiliser}\par

\begin{aretenir}[Boîte à outils linguistique]
\textbf{Lexique des TIC et des loisirs numériques :} \eng{screen time} (temps d'écran), \eng{online gaming} (jeux en ligne), \eng{social media} (réseaux sociaux), \eng{to chat / to scroll / to stream} (discuter / faire défiler / regarder en streaming), \eng{addiction} (dépendance), \eng{to stay connected} (rester connecté), \eng{digital balance} (équilibre numérique), \eng{offline activities} (activités hors ligne).

\textbf{Exprimer un conseil :}
\begin{itemize}
\item Impératif : \eng{Limit your screen time.} « Limite ton temps d'écran. »
\item \eng{You should / shouldn't + base verbale} : \eng{You shouldn't play games before exams.}
\item \eng{It is important to + base verbale} : \eng{It is important to sleep eight hours.}
\item \eng{Why not + base verbale ?} : \eng{Why not read a book instead?}
\end{itemize}

\textbf{Justifier :} \eng{because + phrase} ; \eng{so that + phrase} ; \eng{in order to + base verbale}.

\textbf{Donner une opinion :} \eng{In my opinion\ldots}, \eng{I think / I believe that\ldots}, \eng{For me, the best tip is\ldots}
\end{aretenir}

\begin{attention}[Erreurs fréquentes des francophones]
\begin{itemize}
\item Ne dis pas \eng{*to make sport} ou \eng{*to make games} : dis \eng{to play games}, \eng{to do sport} (ou \eng{exercise}).
\item \eng{Advice} est \textbf{indénombrable} : \eng{a piece of advice}, jamais \eng{*an advice}.
\item Après \eng{should}, toujours la \textbf{base verbale} sans \eng{to} : \eng{you should limit}, pas \eng{*you should to limit}.
\item Attention au faux ami : \eng{actually} signifie « en fait », pas « actuellement » (\eng{currently}).
\item Ordre des mots : l'adjectif \textbf{avant} le nom — \eng{a screen trap}, pas \eng{*a trap screen}.
\item \eng{Switch off} (verbe à particule) : \eng{switch it off} (pronom \textbf{avant} la particule), pas \eng{*switch off it}.
\end{itemize}
\end{attention}

\textbf{Proposition de production et corrigé commenté}\par

\begin{exoresolu}[Affiche modèle du groupe 3 et présentation orale]
Rédige l'affiche du concours \eng{“Smart fun or screen trap?”} avec trois conseils justifiés, puis le script de la présentation orale d'une minute.
\tcblower
\textbf{Affiche (production modèle) :}

\eng{SMART FUN OR SCREEN TRAP?}

\eng{“Your phone is a tool, not a master!”}

\eng{1. Limit your recreational screen time to two hours a day, because your eyes and your brain need rest.}

\eng{2. Switch off your phone during meals and family time, so that you can enjoy real conversations.}

\eng{3. Choose one offline activity every evening — sport, reading or music — in order to keep a healthy digital balance.}

\eng{Enjoy your screen — don't let it enjoy you!}

\textbf{Script de présentation orale :}

\eng{Good morning everyone. Our group believes that ICTs are wonderful recreational resources, but they can become a screen trap if we are not careful. That is why our poster says: “Your phone is a tool, not a master.” Our first tip is to limit screen time to two hours a day, because too much scrolling causes headaches and poor concentration. Our second tip is to switch off phones during meals, so that families can talk to each other again. Finally, we recommend one offline activity every evening, in order to sleep better and stay healthy. In our opinion, smart fun is possible — vote for our tips! Thank you.}

\textbf{Commentaire des choix de langue :}
\begin{itemize}
\item Les trois conseils sont à l'\textbf{impératif} (\eng{limit}, \eng{switch off}, \eng{choose}), forme la plus naturelle pour une affiche de sensibilisation.
\item Chaque conseil est \textbf{justifié} par \eng{because}, \eng{so that} et \eng{in order to}, ce qui montre la maîtrise des trois structures de justification.
\item Le lexique des TIC est bien présent : \eng{screen time}, \eng{scrolling}, \eng{offline activity}, \eng{digital balance}, \eng{switch off}.
\item Le slogan utilise un \textbf{parallélisme} (\eng{Enjoy your screen — don't let it enjoy you!}), procédé typique des campagnes anglophones.
\item La présentation orale suit un plan clair : salutation, opinion (\eng{our group believes that\ldots}), annonce des trois conseils avec des connecteurs (\eng{first}, \eng{second}, \eng{finally}), conclusion avec appel au vote.
\end{itemize}
\end{exoresolu}

\textbf{Bilan de l'activité}\par
Cette activité d'intégration a permis de réunir, dans une tâche de communication authentique et proche de la vie réelle des élèves, l'ensemble des savoir-faire de la leçon : le \cle{scanning} pour extraire rapidement des informations d'un texte, le \cle{lexique des TIC} réemployé dans un nouveau contexte, les structures de \cle{conseil} et d'\cle{opinion}, ainsi que la production écrite (l'affiche) et orale (la présentation). Le vote argumenté ajoute une dimension d'écoute active et d'interaction réelle : chaque élève comprend les affiches des autres et justifie son choix en anglais.

Pour prolonger l'activité à la maison, chaque élève peut rédiger un court paragraphe d'opinion (six à huit lignes) répondant à la question : \eng{“Are ICTs at home more a source of smart fun or a screen trap for you? Give two reasons.”} Ce paragraphe pourra servir de point de départ à la séquence de \eng{writing} du chapitre. Le professeur veillera à valoriser les trois conseils gagnants en les affichant durablement dans la salle, transformant ainsi la production de la classe en véritable outil de sensibilisation de la communauté scolaire.

\newpage

\coursec{Méthodes et exercices résolus}

\coursec{Lesson 55 — Reading: Texts on using ICTs as recreational resources at home or in the community}

\courssub{Warm-up and Pre-reading}

\begin{experience}[Discussion : vos habitudes numériques]
\begin{enumerate}
\item En binôme, discutez en anglais pendant 2 minutes : \eng{How do you usually spend your free time at home? Do you use a smartphone, a computer or a TV?}
\item Relevez au tableau les activités citées : \eng{watch videos, play games, chat on WhatsApp, listen to music, scroll social media, stream films}.
\item Le professeur écrit les mots-clés : \eng{smartphone, screen time, online, connected, entertainment, telecentre}. Prononciation collective : « smaurt-fone », « skriin taim », « on-lain », « kœ-nekt-id », « èn-tèr-tèin-ment », « tè-le-sen-tr ».
\end{enumerate}
\end{experience}

\courssub{Reading Text: “Weekend Screens”}

\begin{textebox}[Support text — “Weekend Screens: How Cameroonian Youth Relax with ICTs”]
On Saturday afternoons in Yaoundé, many teenagers do not go to the stadium any more; they stay connected. For sixteen-year-old Larissa, the smartphone is “my cinema, my radio and my playground”. She watches comedy sketches on YouTube, chats with her cousins in Bamenda and plays online chess against strangers from Kenya or Ghana.

Her father, Mr Etoundi, is not worried. “When I was young, we listened to the radio for hours. My daughter does the same thing, only the tool has changed,” he explains. At the community telecentre in his neighbourhood, young people who have no computer at home can learn graphic design, edit videos of local football matches and even record songs.

However, doctors warn against excessive screen time. “Recreation should refresh the mind, not exhaust the eyes,” says Dr Mbappe, who advises the 20-20-20 rule: every twenty minutes, look at something twenty feet away for twenty seconds. For Larissa, balance is the key word: “After two hours online, I go out and play handball. The screen is a resource, not a prison.”
\end{textebox}

\begin{aretenir}[Glossaire du texte (mots difficiles)]
\begin{itemize}
\item \eng{sketch} (n.) : sketch, courte scène comique.
\item \eng{playground} (n.) : aire de jeux ; ici : espace de divertissement.
\item \eng{telecentre} (n.) : centre communautaire d'accès à l'informatique et à Internet.
\item \eng{graphic design} (n.) : design graphique, création visuelle assistée par ordinateur.
\item \eng{excessive} (adj.) : excessif, trop grand.
\item \eng{refresh} (v.) : rafraîchir, revigorer (l'esprit).
\item \eng{exhaust} (v.) : épuiser, fatiguer complètement.
\item \eng{balance} (n.) : équilibre, juste milieu.
\item \eng{resource} (n.) : ressource, moyen utile.
\item \eng{prison} (n.) : prison (métaphore de l'enfermement).
\end{itemize}
\end{aretenir}

\courssub{Reading Strategies: Skimming, Scanning, Guessing}

\begin{methode}[Réussir la compréhension d'un texte sur les TIC et les loisirs numériques]
\begin{enumerate}
\item \textbf{Avant de lire} : lis le titre, la source et la première phrase. Pose-toi deux questions : \eng{Who?} (qui parle ?) et \eng{What about?} (de quoi s'agit-il ?). Un texte sur \eng{ICTs as recreational resources} parlera de \eng{video games}, \eng{social media}, \eng{streaming}, \eng{online music}, \eng{community forums}.
\item \textbf{Skimming} (survol) : lis rapidement tout le texte en 2 minutes pour dégager l'idée générale. Ne t'arrête pas sur les mots inconnus.
\item \textbf{Scanning} (repérage) : lis les questions, souligne les mots-clés (noms propres, chiffres, verbes d'action), puis retourne dans le texte pour localiser précisément l'information.
\item \textbf{Déduction du sens} : pour un mot inconnu, utilise le contexte (les mots autour), la famille de mots (\eng{entertain} $\rightarrow$ \eng{entertainment}) et les préfixes/suffixes (\eng{useless} = sans utilité).
\item \textbf{Répondre} : réponds toujours par une phrase complète en anglais, en reprenant les termes de la question, jamais par un seul mot. Cite le texte si on te demande de justifier (\eng{Justify your answer with a quotation from the text} « Justifiez votre réponse par une citation du texte »).
\item \textbf{Vérifier} : relis chaque réponse : sujet-verbe, temps, orthographe des mots repris du texte.
\end{enumerate}
\end{methode}

\begin{methode}[Deviner le sens d'un mot sans dictionnaire (Guessing from context)]
\begin{enumerate}
\item \textbf{Identifie la nature du mot} : nom, verbe, adjectif ? La terminaison aide : \emph{-tion/-ment/-ness} = nom ; \emph{-ful/-less/-ive/-able} = adjectif ; \emph{-ise/-ize/-en} = verbe.
\item \textbf{Exploite la famille de mots} : \eng{recreation} vient de \eng{create} ; \eng{excessive} vient de \eng{excess} (excès).
\item \textbf{Lis la phrase entière} : le contexte donne souvent une définition déguisée. Ex. : \eng{Recreation should refresh the mind, not exhaust the eyes} « Les loisirs doivent rafraîchir l'esprit, pas épuiser les yeux » — \eng{exhaust} s'oppose à \eng{refresh}, donc « fatiguer ».
\item \textbf{Cherche les connecteurs} : \eng{however, but} annoncent un contraste ; \eng{because, so} une cause ou conséquence ; \eng{for example} une illustration.
\item \textbf{Vérifie} : remplace le mot par ta traduction dans la phrase ; si la phrase reste logique, c'est bon.
\end{enumerate}
\end{methode}

\courssub{Comprehension Practice}

\begin{exoresolu}[Comprehension questions (type Bac)]
Read the text above and answer in complete sentences. 
1) What does Larissa use her smartphone for? 
2) Why is Mr Etoundi not worried about his daughter? 
3) What activities are offered at the community telecentre? 
4) Explain the 20-20-20 rule in your own words.
\tcblower
\textbf{1)} Mots-clés de la question : \eng{Larissa}, \eng{smartphone}. Scanning $\rightarrow$ paragraphe 1. Réponse : \eng{Larissa uses her smartphone to watch comedy sketches on YouTube, to chat with her cousins in Bamenda and to play online chess.} « Larissa utilise son smartphone pour regarder des sketches comiques sur YouTube, discuter avec ses cousins à Bamenda et jouer aux échecs en ligne. » (On reprend les trois activités citées ; attention à l'infinitif de but \eng{to watch / to chat / to play}.)

\textbf{2)} Mots-clés : \eng{Mr Etoundi}, \eng{not worried}. Paragraphe 2. Réponse : \eng{Mr Etoundi is not worried because he thinks his daughter does the same thing he did with the radio when he was young; only the tool has changed.} « M. Etoundi n'est pas inquiet car il pense que sa fille fait la même chose que lui avec la radio quand il était jeune ; seul l'outil a changé. » (Justification par la comparaison père/fille.)

\textbf{3)} Mots-clés : \eng{telecentre}, \eng{activities}. Réponse : \eng{At the community telecentre, young people can learn graphic design, edit videos of local football matches and record songs.} « Au télécentre communautaire, les jeunes peuvent apprendre le design graphique, monter des vidéos de matchs de football locaux et enregistrer des chansons. »

\textbf{4)} Reformulation (own words) : \eng{The 20-20-20 rule means that every twenty minutes, you should stop looking at your screen and look at something far away for twenty seconds, in order to rest your eyes.} « La règle 20-20-20 signifie que toutes les vingt minutes, vous devez cesser de regarder votre écran et regarder quelque chose de loin pendant vingt secondes, afin de reposer vos yeux. » (On évite de recopier mot à mot ; \eng{twenty feet away} devient \eng{something far away}.)

\textbf{Conclusion méthodologique} : chaque réponse est une phrase complète, ancrée dans le paragraphe repéré par scanning, sans recopie intégrale inutile.
\end{exoresolu}

\courssub{Thematic Vocabulary: ICTs and Digital Leisure}

\begin{methode}[Apprendre et réutiliser le vocabulaire des TIC]
\begin{enumerate}
\item \textbf{Classe le lexique en familles} : le matériel (\eng{smartphone, laptop, screen, headset}), les actions (\eng{to browse, to stream, to download, to upload, to chat, to post, to scroll}), les lieux et services (\eng{telecentre, cybercafé, social network, gaming platform}), les personnes (\eng{user, gamer, follower, subscriber}).
\item \textbf{Apprends les collocations} (mots qui vont ensemble) : \eng{to spend time online}, \eng{to play video games}, \eng{screen time}, \eng{to surf the Internet}, \eng{to stay connected}.
\item \textbf{Distingue les faux amis} : \eng{actually} = en fait (pas « actuellement » = \eng{currently}) ; \eng{a library} = une bibliothèque (pas une librairie = \eng{a bookshop}) ; \eng{sensible} = raisonnable (sensible = \eng{sensitive}).
\item \textbf{Réemploie immédiatement} : fabrique une phrase personnelle avec chaque mot nouveau (\eng{I stream music every evening on my phone.} « Je diffuse de la musique en streaming chaque soir sur mon téléphone. »).
\end{enumerate}
\end{methode}

\begin{center}
\begin{tabularx}{\linewidth}{|X|X|X|}
\hline
\rowcolor{popblueL} English & Français & Exemple \\
\hline
to browse the Internet & naviguer sur Internet & \eng{She browses the Internet for music videos.} « Elle navigue sur Internet pour des clips. » \\
\hline
screen time & temps d'écran & \eng{Doctors advise limiting screen time.} « Les médecins conseillent de limiter le temps d'écran. » \\
\hline
to stream a film & regarder en streaming & \eng{We streamed a football match last night.} « Nous avons regardé un match en streaming hier soir. » \\
\hline
a telecentre & un télécentre & \eng{The telecentre offers free computer lessons.} « Le télécentre propose des cours d'informatique gratuits. » \\
\hline
to stay connected & rester connecté & \eng{He stays connected with his cousins abroad.} « Il reste en contact avec ses cousins à l'étranger. » \\
\hline
an online game & un jeu en ligne & \eng{Online games can be educational.} « Les jeux en ligne peuvent être éducatifs. » \\
\hline
\end{tabularx}
\end{center}

\begin{exoresolu}[Vocabulary gap-fill (type Bac)]
Complete with words from the box: \eng{stream — screen time — telecentre — connected — browse}. 
1) Teenagers should reduce their \dots\ to protect their eyes. 
2) In my neighbourhood, students without computers go to the \dots\ to do research. 
3) My uncle likes to \dots\ the Internet for news every morning. 
4) During the holidays, we \dots\ films on my father's tablet. 
5) WhatsApp helps families in the village stay \dots\ with relatives in town.
\tcblower
\textbf{1)} \eng{screen time} — collocation avec \eng{reduce} et indice \eng{protect their eyes} (santé des yeux).
\textbf{2)} \eng{telecentre} — indice : \eng{students without computers} cherchent un lieu public équipé.
\textbf{3)} \eng{browse} — collocation : \eng{browse the Internet} ; l'indice \eng{for news} confirme la navigation.
\textbf{4)} \eng{stream} — regarder des films en ligne sans les télécharger ; sujet \eng{we} au présent, donc forme de base.
\textbf{5)} \eng{connected} — collocation : \eng{stay connected with somebody} ; sens : garder le contact.

\textbf{Conclusion} : dans un texte à trous, cherche d'abord la collocation (\eng{stay + connected}), puis l'indice de sens dans la phrase.
\end{exoresolu}

\begin{exoresolu}[Vocabulary in context (type Bac)]
Find in the text words or expressions meaning: 
a) se détendre / se rafraîchir l'esprit (paragraph 3) ; 
b) trop de temps passé devant un écran (paragraph 3) ; 
c) un équilibre (paragraph 3). 
Then explain in French what “the screen is a resource, not a prison” means.
\tcblower
\textbf{a)} \eng{refresh the mind} = « rafraîchir l'esprit, se ressourcer ». Indice : le verbe \eng{refresh} contient \eng{fresh} (frais) ; le contexte médical du paragraphe 3 confirme l'idée de repos.

\textbf{b)} \eng{excessive screen time} = « temps d'écran excessif ». Famille de mots : \eng{excess} (excès) + suffixe adjectival \emph{-ive}. Le verbe \eng{warn against} (mettre en garde contre) signale qu'il s'agit d'un danger.

\textbf{c)} \eng{balance} = « l'équilibre ». Indice de contexte : Larissa alterne deux heures en ligne et du sport dehors, donc elle cherche un juste milieu.

\textbf{Explication de la métaphore} : \eng{The screen is a resource, not a prison} « L'écran est une ressource, pas une prison » signifie que l'écran est un outil utile dont on peut profiter librement, mais qu'il ne doit pas nous enfermer ni nous empêcher de vivre d'autres expériences (sport, amis, plein air). La métaphore de la prison évoque la dépendance aux écrans.

\textbf{Conclusion} : nature du mot + famille + contexte + vérification = traduction fiable sans dictionnaire.
\end{exoresolu}

\courssub{Guided Production: Summary and Opinion}

\begin{methode}[Écrire un summary efficace]
\begin{enumerate}
\item \textbf{Repère les 3 ou 4 idées principales} (une par paragraphe en général) : ici, 1) les jeunes utilisent les TIC pour se divertir ; 2) les parents et les télécentres soutiennent cet usage ; 3) les médecins recommandent la modération.
\item \textbf{Reformule} : interdiction de recopier des phrases entières du texte. Utilise des synonymes : \eng{teenagers} $\rightarrow$ \eng{young people} ; \eng{worried} $\rightarrow$ \eng{concerned}.
\item \textbf{Utilise des connecteurs} : \eng{First, Moreover, However, In conclusion}.
\item \textbf{Reste neutre} : pas d'opinion personnelle dans un résumé ; pas de détails (noms propres, exemples précis).
\item \textbf{Respecte la limite} : si 3 sentences sont demandées, écris exactement 3 phrases. Compte les mots si une limite est donnée.
\end{enumerate}
\end{methode}

\begin{exoresolu}[Summary writing (type Bac)]
Summarize the text “Weekend Screens” in three sentences (about 40 words).
\tcblower
\textbf{Étape 1 — Idées principales} : P1 = les TIC servent de loisirs aux jeunes ; P2 = les adultes et les télécentres encadrent positivement cet usage ; P3 = les médecins conseillent de limiter le temps d'écran.

\textbf{Étape 2 — Reformulation et connecteurs.}

\textbf{Réponse rédigée} : \eng{Nowadays, many young Cameroonians use ICTs such as smartphones and the Internet for entertainment. Parents and community telecentres generally support this digital leisure, which can even be educational. However, doctors recommend limiting screen time and keeping a balance with outdoor activities.} « De nos jours, de jeunes Camerounais utilisent les TIC comme smartphones et Internet pour se divertir. Parents et télécentres soutiennent généralement ce loisir numérique, qui peut même être éducatif. Cependant, les médecins recommandent de limiter le temps d'écran et de garder un équilibre avec les activités de plein air. » (38 mots.)

\textbf{Vérification} : 3 phrases exactement ; aucune phrase copiée du texte ; connecteurs \eng{Nowadays / However} ; aucune opinion personnelle ; présent simple partout (vérités générales).

\textbf{Conclusion} : un bon résumé = idées principales + reformulation + connecteurs + limite respectée.
\end{exoresolu}

\begin{methode}[Exprimer une opinion argumentée (opinion paragraph)]
\begin{enumerate}
\item \textbf{Annonce ta position} avec une formule claire : \eng{In my opinion, \dots} / \eng{I believe that \dots} / \eng{As far as I am concerned, \dots}
\item \textbf{Donne une ou deux raisons} introduites par \eng{because, since, as} ou \eng{First of all, Moreover}.
\item \textbf{Illustre avec un exemple personnel ou local} (ton quartier, ton école, le Cameroun) : \eng{For instance, in my school in Douala, \dots}
\item \textbf{Conclus} en une phrase : \eng{That is why I think that \dots} / \eng{To sum up, \dots}
\item \textbf{Soigne la langue} : présent simple pour les opinions générales ; évite \eng{I am agree} (faute classique : \eng{I agree}).
\end{enumerate}
\end{methode}

\begin{exoresolu}[Giving your opinion (type Bac)]
“ICTs are good recreational resources for young people.” Do you agree? Write a short paragraph (5–6 sentences) giving your opinion.
\tcblower
\textbf{Plan} : position (oui, avec nuance) $\rightarrow$ raison 1 (accès à la culture et aux amis) $\rightarrow$ raison 2 (apprentissage) $\rightarrow$ exemple local $\rightarrow$ nuance (modération) $\rightarrow$ conclusion.

\textbf{Réponse rédigée} : \eng{In my opinion, ICTs are excellent recreational resources for young people, provided they are used wisely. First of all, they allow teenagers to relax and stay in touch with friends, even those who live far away. Moreover, digital leisure can be educational: many students learn English by watching videos or playing online games. For instance, in my neighbourhood in Yaoundé, the telecentre helps pupils discover music production after school. However, young people should not spend the whole day in front of a screen, because their eyes and their body need rest. To sum up, ICTs are a wonderful tool for recreation as long as we keep a healthy balance.} « À mon avis, les TIC sont d'excellentes ressources de loisir pour les jeunes, à condition qu'elles soient utilisées sagement. D'abord, elles permettent aux adolescents de se détendre et de garder le contact avec leurs amis, même ceux qui habitent loin. De plus, le loisir numérique peut être éducatif : beaucoup d'élèves apprennent l'anglais en regardant des vidéos ou en jouant en ligne. Par exemple, dans mon quartier à Yaoundé, le télécentre aide les élèves à découvrir la production musicale après l'école. Cependant, les jeunes ne devraient pas passer la journée entière devant un écran, car leurs yeux et leur corps ont besoin de repos. En résumé, les TIC sont un outil merveilleux pour le loisir tant qu'on maintient un équilibre sain. »

\textbf{Justifications grammaticales} : présent simple pour les généralités (\eng{allow, learn, need}) ; \eng{provided (that)} + présent = « à condition que » ; \eng{as long as} = « tant que » ; gérondif après \eng{by} (\eng{by watching, by playing}).

\textbf{Conclusion} : opinion claire + deux raisons + exemple local + nuance + conclusion = paragraphe d'opinion complet et bien noté.
\end{exoresolu}

\courssub{Common Errors to Avoid}

\begin{attention}[Les 5 erreurs qui coûtent des points au Bac]
\begin{enumerate}
\item \textbf{Faux amis} : ne traduis jamais « actuellement » par \eng{actually} (qui signifie « en fait »). Dis \eng{currently} ou \eng{nowadays}. De même, « une librairie » se dit \eng{a bookshop}, pas \eng{a library} (bibliothèque). « Sensible » (qui ressent) = \eng{sensitive} ; \eng{sensible} = raisonnable.
\item \textbf{Calque du français} : « Je suis d'accord » ne se traduit pas par \eng{I am agree} mais par \eng{I agree} (\eng{agree} est un verbe). Et « passer du temps » = \eng{to spend time}, pas \eng{to pass time}.
\item \textbf{Temps verbaux} : pour une opinion générale ou une habitude, utilise le présent simple (\eng{Teenagers play online games}), pas le présent continu. Réserve le continu à l'action en cours : \eng{Larissa is playing chess now.}
\item \textbf{Prépositions} : \eng{to listen \textbf{to} music}, \eng{to depend \textbf{on} the Internet}, \eng{to be interested \textbf{in} video games}, \eng{to chat \textbf{with} friends}, \eng{\textbf{on} the Internet}. Une préposition fausse coûte un point à chaque occurrence.
\item \textbf{Orthographe et ordre des mots} : \eng{entertainment} (un seul \emph{a}, pas \emph{entertainement}), \eng{which} (pas \emph{wich}), \eng{because} (pas \emph{because of} devant une phrase). L'adjectif précède toujours le nom : \eng{a recreational resource}, jamais \eng{a resource recreational}.
\end{enumerate}
\end{attention}

\newpage

\coursec{Exercices}

\courssub{Je vérifie mes connaissances}

\begin{exercice}[QCM : comprendre le texte de la leçon]
Relis le texte de la leçon, puis choisis la bonne réponse pour chaque question.
\begin{enumerate}
\item The main topic of the text is:
\begin{enumerate}
\item the dangers of the internet
\item ICTs as recreational resources at home and in the community
\item how to repair a computer
\item the history of television
\end{enumerate}
\item According to the text, young people in Cameroon mostly use their smartphones to:
\begin{enumerate}
\item do their homework only
\item listen to music, chat and watch videos
\item read newspapers
\item make phone calls to their teachers
\end{enumerate}
\item A community telecentre is a place where:
\begin{enumerate}
\item people buy mobile phones
\item people can access computers and the internet
\item films are produced
\item students sleep
\end{enumerate}
\item The word \eng{leisure} in the text means:
\begin{enumerate}
\item work \item free time \item money \item school
\end{enumerate}
\end{enumerate}
\end{exercice}

\begin{exercice}[Vrai ou faux, avec justification]
Lis les phrases suivantes. Écris \eng{True} ou \eng{False}, puis justifie ta réponse par une courte phrase tirée ou inspirée du texte de la leçon.
\begin{enumerate}
\item ICTs can only be used for serious work, never for fun.
\item Many families in Cameroon watch television together in the evening.
\item Playing video games all night has no effect on a student's health.
\item Cybercafés and telecentres give internet access to people who have no connection at home.
\item The text advises young people to balance screen time with other activities.
\end{enumerate}
\end{exercice}

\begin{exercice}[Vocabulaire : à chaque mot sa définition]
Relie chaque mot de la colonne A à sa définition de la colonne B.
\begin{center}
\begin{tabularx}{\linewidth}{|l|X|}
\hline
\rowcolor{popblueL} \textbf{A — Word} & \textbf{B — Definition} \\
\hline
1. to browse & a. a public place with computers connected to the internet \\
\hline
2. a cybercafé & b. to look at pages on the internet \\
\hline
3. to stream & c. a game played on a computer or a phone \\
\hline
4. a video game & d. to watch or listen to content directly online \\
\hline
5. a password & e. a secret word used to open an account \\
\hline
6. to download & f. to copy a file from the internet onto your device \\
\hline
\end{tabularx}
\end{center}
\end{exercice}

\begin{exercice}[Grammaire : complète avec le bon temps]
Complète les phrases avec le verbe entre parenthèses au present simple, au present continu ou au present perfect.
\begin{enumerate}
\item Every weekend, my cousin \dots{} (watch) football matches on his phone.
\item Listen! The children \dots{} (play) a video game in the living room.
\item I \dots{} never \dots{} (use) a telecentre, but I would like to try one.
\item At the moment, more and more Cameroonians \dots{} (buy) smartphones.
\item She \dots{} (download) three songs since this morning.
\item Students usually \dots{} (chat) with their friends after school.
\end{enumerate}
\end{exercice}

\courssub{Je m'entraîne}

\begin{exercice}[Traduction : du français vers l'anglais]
Traduis les phrases suivantes en anglais correct.
\begin{enumerate}
\item Les jeunes utilisent les TIC pour se divertir à la maison et dans la communauté.
\item Mon frère passe deux heures par jour sur les réseaux sociaux.
\item Au cybercafé, on peut naviguer sur internet et imprimer des documents.
\item Regarder des films en ligne est mon loisir préféré le week-end.
\item Les parents devraient contrôler le temps que leurs enfants passent devant les écrans.
\end{enumerate}
\end{exercice}

\begin{exercice}[Reformulation : dire autrement]
Réécris chaque phrase avec tes propres mots, sans changer le sens. C'est un entraînement au résumé (\eng{paraphrasing}).
\begin{enumerate}
\item ICTs offer many recreational opportunities to young people.
\item Instead of going out, some teenagers prefer chatting online.
\item The telecentre allows the whole community to enjoy digital resources.
\item Too much screen time can be harmful to students' health.
\item Listening to music on the phone helps me relax after school.
\end{enumerate}
\end{exercice}

\begin{exercice}[Texte à trous : un week-end numérique]
Complète le texte avec les mots de la liste suivante (chaque mot ne sert qu'une fois) :
\begin{center}
\eng{streamed — password — downloaded — cybercafé — leisure — chatted — screen — online}
\end{center}
\begin{textebox}[A digital weekend]
Last Saturday, I had a lot of free time, so I decided to enjoy my favourite \dots{} activities. In the morning, I went to the \dots{} near the market because we have no internet at home. I typed my \dots{} to open my email account and I \dots{} with my cousin who lives in Douala. Then I \dots{} two songs by my favourite artist and I \dots{} a football match on my phone. In the evening, my little sister wanted to play \dots{} games, but my mother told her not to spend too much time in front of the \dots{}. We finally watched a film together as a family.
\end{textebox}
\end{exercice}

\begin{exercice}[Classement : Benefits, Risks, Places, Activities]
Classe les douze expressions suivantes dans le tableau, dans la bonne colonne.
\begin{center}
\eng{watching films online — eye problems — a telecentre — chatting with friends — making new friends abroad — addiction to games — a cybercafé — listening to music — learning new skills — at home — cyberbullying — playing video games}
\end{center}
\begin{center}
\begin{tabularx}{\linewidth}{|X|X|X|X|}
\hline
\rowcolor{popblueL} \textbf{Benefits} & \textbf{Risks} & \textbf{Places} & \textbf{Activities} \\
\hline
 & & & \\
\hline
 & & & \\
\hline
 & & & \\
\hline
\end{tabularx}
\end{center}
\end{exercice}

\courssub{Je me prépare au Bac}

\begin{exercice}[Type Bac — Reading comprehension]
Lis le texte suivant, puis réponds aux questions.

\begin{textebox}[Social media and Ghanaian teenagers]
In Accra, the capital of Ghana, social media have become the favourite recreational resource of many teenagers. After school, seventeen-year-old Ama does not go straight home. She stops at a small internet café near her school, where she spends about one hour chatting with her friends on her phone. “It is cheap and the connection is fast,” she explains with a smile. At home, her younger brother Kofi prefers watching comedy videos and playing football games online.

Their mother, Mrs Mensah, is worried. “When I was young, we played outside with our friends. Today, children stay alone with their screens,” she says. However, she admits that social media also have advantages. Ama, for example, belongs to an online reading club where students discuss books in English every Friday evening. Kofi has learned to edit videos, a skill he wants to use in the future.

A teacher at Ama's school, Mr Owusu, believes that the problem is not the technology itself but the way young people use it. “Social media can open the world to our students,” he says, “but they must learn to control their screen time.” His school now organises workshops to teach pupils how to use the internet safely and wisely. The students learn to protect their passwords, to avoid strangers online and to check information before sharing it.

For Ama, the message is clear: “My phone is like a friend, but a friend should not take all your time.”
\end{textebox}

\textbf{A. True or False?} Justify each answer with a short quotation from the text. (5 marks)
\begin{enumerate}
\item Ama goes home immediately after school.
\item Kofi enjoys playing games on the internet.
\item Mrs Mensah thinks social media have only disadvantages.
\item Ama discusses books online every Friday evening.
\item Mr Owusu's school teaches students how to use the internet safely.
\end{enumerate}

\textbf{B. Choose the correct answer.} (3 marks)
\begin{enumerate}
\item Ama stops at the internet café because:
\begin{enumerate}
\item it is expensive \item the connection is fast \item her mother works there
\end{enumerate}
\item Kofi wants to use his video-editing skill:
\begin{enumerate}
\item in the future \item to become a teacher \item to punish his sister
\end{enumerate}
\item According to Mr Owusu, the real problem is:
\begin{enumerate}
\item the technology itself \item the way young people use it \item the price of phones
\end{enumerate}
\end{enumerate}

\textbf{C. Vocabulary.} Find in the text words or expressions that mean: (2 marks)
\begin{enumerate}
\item worried (paragraph 3)
\item abilities (paragraph 3)
\item to keep safe (paragraph 4)
\item in a sensible way (paragraph 4)
\end{enumerate}
\end{exercice}

\begin{exercice}[Type Bac — Scanning chronométré et déduction lexicale]
\textbf{A. Scanning (3 minutes).} Relis le texte \eng{Social media and Ghanaian teenagers} et retrouve le plus vite possible les cinq informations suivantes :
\begin{enumerate}
\item the name of the city where the story takes place
\item the time Ama spends at the internet café
\item the names of the two children
\item one piece of advice given during the school workshops
\end{enumerate}

\textbf{B. Guessing meaning from context.} Devine le sens des mots soulignés dans le paragraphe suivant grâce au contexte, donne un équivalent français, puis vérifie au dictionnaire.
\begin{textebox}[A new habit]
Every evening, after dinner, the Ngo family \textbf{gathers} in the living room to watch a film together. The children are \textbf{glued} to the screen and nobody wants to miss the end of the story. But last week, their father decided to \textbf{limit} screen time to one hour a day. At first, the children were \textbf{upset}, but now they enjoy playing cards and talking together. Their mother says this new habit has \textbf{strengthened} the family \textbf{bond}.
\end{textebox}
Mots à expliquer : \eng{gathers}, \eng{glued}, \eng{limit}, \eng{upset}, \eng{strengthened}, \eng{bond}.
\end{exercice}

\begin{exercice}[Type Bac — Guided composition]
\textbf{Topic:} \eng{Do you think ICTs are good recreational resources for students?}

Write a composition of \textbf{80 to 100 words} giving your opinion. Your composition must contain:
\begin{itemize}
\item an introduction stating your opinion clearly (\eng{In my opinion...} / \eng{I believe that...});
\item \textbf{two arguments} supported by examples from your life or your community (home, school, cybercafé, telecentre);
\item one expression of contrast (\eng{however}, \eng{on the other hand}, \eng{although});
\item a short conclusion.
\end{itemize}

\textbf{Useful language:} \eng{in my opinion — first of all — moreover — for example — however — on the other hand — to conclude — as far as I am concerned}.
\end{exercice}

\newpage

\coursec{Corrigés des exercices}

\begin{corrige}[Exercice 1 — QCM : comprendre le texte de la leçon]
\begin{enumerate}
\item \textbf{Correct answer: b)} \eng{ICTs as recreational resources at home and in the community.} — C'est le thème central du texte : les usages de loisir des technologies, à la maison et dans les lieux publics.
\item \textbf{Correct answer: b)} \eng{listen to music, chat and watch videos.} — Le texte précise que les jeunes Camerounais utilisent surtout leur smartphone pour se divertir (musique, discussions, vidéos), et non uniquement pour travailler.
\item \textbf{Correct answer: b)} \eng{people can access computers and the internet.} — Un \eng{telecentre} est un lieu communautaire équipé d'ordinateurs connectés, ouvert au public.
\item \textbf{Correct answer: b)} \eng{free time.} — \eng{Leisure} signifie « temps libre, loisirs ». Attention au faux ami : ce n'est ni « travail » ni « argent ».
\end{enumerate}
\textbf{Point de méthode :} au QCM, élimine d'abord les réponses absurdes (\eng{students sleep}, \eng{how to repair a computer}), puis vérifie chaque réponse plausible dans le texte avant de choisir.
\end{corrige}

\begin{corrige}[Exercice 2 — Vrai ou faux, avec justification]
\begin{enumerate}
\item \textbf{False.} \eng{ICTs can be used for fun: people listen to music, watch films and play games.} — Le texte montre précisément que les TIC servent aussi aux loisirs.
\item \textbf{True.} \eng{Many families in Cameroon watch television together in the evening.} — Le texte décrit ce moment familial du soir devant la télévision.
\item \textbf{False.} \eng{Playing video games all night is bad for a student's health: it causes tiredness and eye problems.} — Le texte mentionne les risques de l'excès d'écran (fatigue, problèmes de vue, mauvais résultats scolaires).
\item \textbf{True.} \eng{Cybercafés and telecentres give internet access to people who have no connection at home.} — C'est exactement le rôle de ces lieux communautaires décrit dans le texte.
\item \textbf{True.} \eng{The text advises young people to balance screen time with other activities such as sport and reading.} — La conclusion du texte recommande un usage équilibré.
\end{enumerate}
\textbf{Point de méthode :} une justification doit être une phrase courte, tirée ou inspirée du texte, jamais une simple opinion personnelle.
\end{corrige}

\begin{corrige}[Exercice 3 — Vocabulaire : à chaque mot sa définition]
\begin{center}
\begin{tabularx}{\linewidth}{|l|X|l|}
\hline
\rowcolor{popblueL} \textbf{Word} & \textbf{Definition} & \textbf{Traduction} \\
\hline
1. to browse & b. to look at pages on the internet & naviguer (sur internet) \\
\hline
2. a cybercafé & a. a public place with computers connected to the internet & un cybercafé \\
\hline
3. to stream & d. to watch or listen to content directly online & regarder/écouter en streaming \\
\hline
4. a video game & c. a game played on a computer or a phone & un jeu vidéo \\
\hline
5. a password & e. a secret word used to open an account & un mot de passe \\
\hline
6. to download & f. to copy a file from the internet onto your device & télécharger \\
\hline
\end{tabularx}
\end{center}
\textbf{Réponses :} 1-b, 2-a, 3-d, 4-c, 5-e, 6-f.
\textbf{Point de vigilance :} ne confonds pas \eng{to download} (télécharger vers ton appareil) et \eng{to upload} (envoyer un fichier vers internet).
\end{corrige}

\begin{corrige}[Exercice 4 — Grammaire : complète avec le bon temps]
\begin{enumerate}
\item \eng{Every weekend, my cousin \textbf{watches} football matches on his phone.} — Present simple : habitude marquée par \eng{every weekend} ; 3\textsuperscript{e} personne du singulier, d'où \eng{watch\textbf{es}}.
\item \eng{Listen! The children \textbf{are playing} a video game in the living room.} — Present continu : action en cours, signalée par \eng{Listen!} ; \eng{play} ne prend qu'un seul \eng{-ing} : \eng{playing}.
\item \eng{I \textbf{have} never \textbf{used} a telecentre, but I would like to try one.} — Present perfect : expérience de vie avec \eng{never}, sans date précise.
\item \eng{At the moment, more and more Cameroonians \textbf{are buying} smartphones.} — Present continu : tendance actuelle, marquée par \eng{at the moment}.
\item \eng{She \textbf{has downloaded} three songs since this morning.} — Present perfect : action commencée dans le passé avec un résultat présent, marquée par \eng{since}.
\item \eng{Students usually \textbf{chat} with their friends after school.} — Present simple : habitude marquée par \eng{usually}.
\end{enumerate}
\textbf{Points de vigilance :} les marqueurs temporels guident le choix du temps (\eng{every weekend, usually} $\rightarrow$ present simple ; \eng{Listen!, at the moment} $\rightarrow$ present continu ; \eng{never, since, already} $\rightarrow$ present perfect). N'oublie pas le \eng{-s} ou \eng{-es} à la 3\textsuperscript{e} personne du singulier au present simple.
\end{corrige}

\begin{corrige}[Exercice 5 — Traduction : du français vers l'anglais]
\begin{enumerate}
\item \eng{Young people use ICTs to entertain themselves at home and in the community.} — Autre formulation correcte : \eng{Young people use ICTs for entertainment at home and in the community.}
\item \eng{My brother spends two hours a day on social media.} — Attention : « deux heures par jour » se dit \eng{two hours a day}, et non \eng{per day} dans ce contexte courant ; \eng{social media} est généralement au pluriel.
\item \eng{At the cybercafé, you can browse the internet and print documents.} — On peut aussi écrire \eng{one can surf the internet} ; « imprimer » se dit \eng{to print}.
\item \eng{Watching films online is my favourite leisure activity at the weekend.} — Le gérondif \eng{watching} est sujet de la phrase ; « le week-end » se dit \eng{at the weekend} (ou \eng{on weekends} en anglais américain).
\item \eng{Parents should control the time their children spend in front of screens.} — \eng{Should} traduit « devraient » ; attention à la structure \eng{the time (that) their children spend}.
\end{enumerate}
\textbf{Points de vigilance :} « se divertir » ne se traduit pas par \eng{to divert oneself} ; utilise \eng{to entertain oneself}, \eng{to have fun} ou \eng{for entertainment}. « Les réseaux sociaux » = \eng{social media} ou \eng{social networks}, jamais \eng{social medias} avec un \eng{-s}.
\end{corrige}

\begin{corrige}[Exercice 6 — Reformulation : dire autrement]
Voici des reformulations possibles (d'autres réponses correctes sont acceptées si le sens est conservé) :
\begin{enumerate}
\item \eng{ICTs offer many recreational opportunities to young people.} $\rightarrow$ \eng{Thanks to ICTs, young people can enjoy many kinds of entertainment.}
\item \eng{Instead of going out, some teenagers prefer chatting online.} $\rightarrow$ \eng{Some teenagers would rather talk to their friends on the internet than go outside.}
\item \eng{The telecentre allows the whole community to enjoy digital resources.} $\rightarrow$ \eng{Thanks to the telecentre, everyone in the community can use computers and the internet.}
\item \eng{Too much screen time can be harmful to students' health.} $\rightarrow$ \eng{Spending too many hours in front of a screen can be dangerous for students.}
\item \eng{Listening to music on the phone helps me relax after school.} $\rightarrow$ \eng{When I listen to music on my phone after school, I feel more relaxed.}
\end{enumerate}
\textbf{Point de méthode :} pour paraphraser, remplace les mots par des synonymes (\eng{harmful} $\rightarrow$ \eng{dangerous}), change la structure de la phrase (actif/passif, nom/verbe), mais ne change jamais l'idée principale.
\end{corrige}

\begin{corrige}[Exercice 7 — Texte à trous : un week-end numérique]
\begin{textebox}[A digital weekend — corrected version]
Last Saturday, I had a lot of free time, so I decided to enjoy my favourite \textbf{leisure} activities. In the morning, I went to the \textbf{cybercafé} near the market because we have no internet at home. I typed my \textbf{password} to open my email account and I \textbf{chatted} with my cousin who lives in Douala. Then I \textbf{downloaded} two songs by my favourite artist and I \textbf{streamed} a football match on my phone. In the evening, my little sister wanted to play \textbf{online} games, but my mother told her not to spend too much time in front of the \textbf{screen}. We finally watched a film together as a family.
\end{textebox}
\textbf{Justifications :}
\begin{itemize}
\item \eng{leisure activities} : « activités de loisirs » — le contexte du temps libre l'impose.
\item \eng{cybercafé} : lieu public avec internet, justifié par \eng{because we have no internet at home}.
\item \eng{password} : on tape un mot de passe pour ouvrir un compte (\eng{to open my email account}).
\item \eng{chatted} : discuter avec un cousin, au prétérit comme tout le récit.
\item \eng{downloaded} : on télécharge des chansons sur son appareil.
\item \eng{streamed} : on regarde un match en direct en ligne.
\item \eng{online games} : « jeux en ligne » ; \eng{online} est ici adjectif.
\item \eng{screen} : \eng{in front of the screen}, « devant l'écran ».
\end{itemize}
\textbf{Point de vigilance :} tous les verbes du récit sont au prétérit (\eng{chatted}, \eng{downloaded}, \eng{streamed}) car l'action est située \eng{last Saturday}.
\end{corrige}

\begin{corrige}[Exercice 8 — Classement : Benefits, Risks, Places, Activities]
\begin{center}
\begin{tabularx}{\linewidth}{|X|X|X|X|}
\hline
\rowcolor{popblueL} \textbf{Benefits} & \textbf{Risks} & \textbf{Places} & \textbf{Activities} \\
\hline
making new friends abroad & eye problems & a telecentre & watching films online \\
\hline
learning new skills & addiction to games & a cybercafé & chatting with friends \\
\hline
--- & cyberbullying & at home & listening to music / playing video games \\
\hline
\end{tabularx}
\end{center}
\textbf{Détail du classement :}
\begin{itemize}
\item \textbf{Benefits} (avantages) : \eng{making new friends abroad}, \eng{learning new skills}.
\item \textbf{Risks} (risques) : \eng{eye problems}, \eng{addiction to games}, \eng{cyberbullying}.
\item \textbf{Places} (lieux) : \eng{a telecentre}, \eng{a cybercafé}, \eng{at home}.
\item \textbf{Activities} (activités) : \eng{watching films online}, \eng{chatting with friends}, \eng{listening to music}, \eng{playing video games}.
\end{itemize}
\textbf{Point de vigilance :} \eng{chatting with friends} est une \emph{activité} ; son avantage possible serait \eng{making new friends}. Distingue bien l'action de son résultat positif.
\end{corrige}

\begin{corrige}[Exercice 9 — Type Bac : Reading comprehension]
\textbf{A. True or False?} (5 marks — 1 mark par réponse justifiée : 0,5 pour \eng{True/False}, 0,5 pour la citation)
\begin{enumerate}
\item \textbf{False.} \eng{“After school, seventeen-year-old Ama does not go straight home. She stops at a small internet café near her school.”}
\item \textbf{True.} \eng{“Kofi prefers watching comedy videos and playing football games online.”}
\item \textbf{False.} \eng{“She admits that social media also have advantages.”} — Mme Mensah reconnaît des avantages malgré ses inquiétudes.
\item \textbf{True.} \eng{“Ama belongs to an online reading club where students discuss books in English every Friday evening.”}
\item \textbf{True.} \eng{“His school now organises workshops to teach pupils how to use the internet safely and wisely.”}
\end{enumerate}
\textbf{B. Choose the correct answer.} (3 marks — 1 mark par item)
\begin{enumerate}
\item \textbf{b)} \eng{the connection is fast} — citation : \eng{“It is cheap and the connection is fast.”}
\item \textbf{a)} \eng{in the future} — citation : \eng{“a skill he wants to use in the future.”}
\item \textbf{b)} \eng{the way young people use it} — citation : \eng{“the problem is not the technology itself but the way young people use it.”}
\end{enumerate}
\textbf{C. Vocabulary.} (2 marks — 0,5 mark par mot)
\begin{enumerate}
\item \eng{worried} = \eng{worried} (Mrs Mensah is worried) — « inquiet ».
\item \eng{abilities} = \eng{skill} / \eng{skills} — « compétences ».
\item \eng{to keep safe} = \eng{to protect} (\eng{to protect their passwords}) — « protéger ».
\item \eng{in a sensible way} = \eng{wisely} (\eng{safely and wisely}) — « sagement, judicieusement ».
\end{enumerate}
\textbf{Points de vigilance :} recopie les citations exactement, entre guillemets ; une justification sans citation perd la moitié des points. Pour le vocabulaire, le mot doit être \emph{tiré du texte}, pas un synonyme trouvé ailleurs.
\end{corrige}

\begin{corrige}[Exercice 10 — Type Bac : Scanning chronométré et déduction lexicale]
\textbf{A. Scanning.} Réponses attendues :
\begin{enumerate}
\item The city: \eng{Accra} (the capital of Ghana).
\item The time Ama spends at the internet café: \eng{about one hour}.
\item The two children: \eng{Ama} and \eng{Kofi}.
\item One piece of advice from the workshops (une seule suffit) : \eng{protect your passwords} / \eng{avoid strangers online} / \eng{check information before sharing it}.
\end{enumerate}
\textbf{Point de méthode :} en scanning, ne relis pas tout le texte : repère les \emph{mots-repères} (noms propres en majuscules, chiffres, mots de la question) et laisse tes yeux sauter directement dessus.

\textbf{B. Guessing meaning from context.}
\begin{center}
\begin{tabularx}{\linewidth}{|l|X|X|}
\hline
\rowcolor{popblueL} \textbf{Word} & \textbf{Clue in the context} & \textbf{French equivalent} \\
\hline
gathers & \eng{in the living room... together} & se réunit, se rassemble \\
\hline
glued & \eng{nobody wants to miss the end} & collés (à l'écran), captivés \\
\hline
limit & \eng{to one hour a day} & limiter, réduire \\
\hline
upset & \eng{At first... but now they enjoy} & contrariés, mécontents \\
\hline
strengthened & \eng{playing cards and talking together} & a renforcé \\
\hline
bond & \eng{the family...} & lien (familial) \\
\hline
\end{tabularx}
\end{center}
\textbf{Point de méthode :} pour deviner un mot, utilise (1) les mots qui l'entourent, (2) les connecteurs de contraste (\eng{but}, \eng{at first... now}), (3) ta connaissance du monde. Vérifie ensuite au dictionnaire : \eng{glued} vient de \eng{glue} « colle » — image très parlante !
\end{corrige}

\begin{corrige}[Exercice 11 — Type Bac : Guided composition]
\textbf{Barème indicatif (10 marks) :}
\begin{itemize}
\item Respect de la consigne et de la longueur (80--100 mots) : 2 marks.
\item Introduction avec opinion claire : 1 mark.
\item Deux arguments illustrés d'exemples : 4 marks (2 par argument).
\item Une expression de contraste correctement employée : 1 mark.
\item Conclusion : 1 mark.
\item Correction grammaticale et vocabulaire du thème : 1 mark.
\end{itemize}

\textbf{Proposition de rédaction modèle (environ 95 mots) :}
\begin{textebox}[Model composition]
In my opinion, ICTs are good recreational resources for students when they are used wisely. First of all, they help us relax after school. For example, I often listen to music and watch comedy videos on my phone in the evening. Moreover, ICTs can teach us new skills: my cousin learned to edit videos at the telecentre near our quarter, and he now makes short films for his friends. However, spending too much time in front of a screen can be harmful to our health and our studies. To conclude, ICTs are excellent tools for leisure, but we must control our screen time.
\end{textebox}

\textbf{Points de vigilance :}
\begin{itemize}
\item Annonce ton opinion dès la première phrase (\eng{In my opinion...}).
\item Chaque argument doit être suivi d'un exemple concret (\eng{For example...}).
\item Place le connecteur de contraste correctement : \eng{However} en début de phrase suivi d'une virgule ; \eng{although} en tête de subordonnée.
\item Évite les faux amis : « se détendre » = \eng{to relax} (pas \eng{to rest oneself}) ; « nuisible » = \eng{harmful}.
\item Compte tes mots : en dessous de 80 ou au-dessus de 100, tu perds des points.
\end{itemize}
\end{corrige}

\newpage

\coursec{Fiche bilan}

\begin{popfigure}
\begin{tikzpicture}[mindmap, grow cyclic, every node/.style={concept, font=\small}, root concept/.append style={font=\bfseries\small, fill=popblue, text=white, minimum size=3.2cm}, level 1/.append style={level distance=3.4cm, sibling angle=51.4, minimum size=2.5cm}, level 1 concept/.append style={font=\footnotesize}]
\node[root concept]{ICTs as recreational resources}
  child[concept color=poporange, text=white]{node[align=center]{Key vocabulary\\ leisure, streaming, gaming}}
  child[concept color=popgreen, text=white]{node[align=center]{Skimming\\ read fast for the\\ general idea}}
  child[concept color=poppurple, text=white]{node[align=center]{Scanning\\ look for dates,\\ names, numbers}}
  child[concept color=poppink, text=white]{node[align=center]{Guessing meaning\\ from context clues}}
  child[concept color=popgold, text=black]{node[align=center]{Grammar\\ gerunds, present\\ perfect, modals}}
  child[concept color=popblue!70, text=white]{node[align=center]{Comprehension\\ questions and\\ interpretation}}
  child[concept color=popmuted, text=white]{node[align=center]{Production\\ summary and\\ opinion}};
\end{tikzpicture}
\legende{Carte mentale de la leçon : les TIC comme ressources de loisir}
\end{popfigure}

\begin{bilan}
\textbf{1. Lexique clé à connaître par cœur}
\begin{itemize}
  \item \eng{leisure / free time} : les loisirs, le temps libre ; \eng{a recreational resource} : une ressource de loisir
  \item \eng{to stream a film / streaming} : regarder en streaming ; \eng{to download / to upload} : télécharger (recevoir / envoyer)
  \item \eng{online gaming / a gamer} : le jeu en ligne / un joueur ; \eng{social media} : les réseaux sociaux
  \item \eng{to browse the internet / to surf the web} : naviguer sur internet ; \eng{a screen} : un écran ; \eng{screen time} : le temps passé devant un écran
  \item \eng{a smartphone, a tablet, a laptop} : un smartphone, une tablette, un ordinateur portable
  \item \eng{a cybercafé / a community centre} : un cybercafé / un centre communautaire
  \item \eng{to chat with friends} : discuter avec des amis ; \eng{to share photos and videos} : partager des photos et des vidéos
  \item \eng{addictive / an addiction} : addictif / une addiction ; \eng{harmful / beneficial} : nuisible / bénéfique
  \item \eng{to relax, to entertain, to connect people} : se détendre, divertir, relier les gens
\end{itemize}

\textbf{2. Grammaire à revoir pour ce thème}
\begin{center}
\begin{tabularx}{\linewidth}{|X|X|X|}\hline
\rowcolor{popblueL} Règle & Exemple anglais & Traduction \\ \hline
Gérondif après \eng{enjoy, like, spend time, be fond of} & \eng{I enjoy playing video games.} & J'aime jouer aux jeux vidéo. \\ \hline
\eng{Spend + time + V-ing} & \eng{She spends two hours chatting online.} & Elle passe deux heures à discuter en ligne. \\ \hline
Present perfect (expérience, bilan) & \eng{Many youths have discovered e-books.} & Beaucoup de jeunes ont découvert les livres numériques. \\ \hline
Modaux d'avis : \eng{should, must, can} & \eng{Parents should limit screen time.} & Les parents devraient limiter le temps d'écran. \\ \hline
Comparatifs & \eng{Streaming is cheaper than going to the cinema.} & Le streaming coûte moins cher que le cinéma. \\ \hline
\end{tabularx}
\end{center}

\textbf{3. Méthodes express de lecture}
\begin{itemize}
  \item \cle{Skimming} : lire le titre, la première phrase de chaque paragraphe et la conclusion pour trouver l'idée générale — ne pas traduire mot à mot.
  \item \cle{Scanning} : repérer rapidement une information précise (chiffre, date, nom propre, lieu) en laissant glisser les yeux sur le texte.
  \item \cle{Guessing} : deviner le sens d'un mot inconnu grâce au contexte, à la nature du mot, à la famille de mots et aux mots français ressemblants (attention aux faux amis).
  \item \cle{Résumé} : 3 à 5 phrases, avec ses propres mots, sans recopier de longues phrases du texte.
  \item \cle{Opinion} : \eng{In my opinion... / I think that... / I agree because... / However...} + un argument + un exemple.
\end{itemize}

\textbf{4. Pièges fréquents des francophones}
\begin{itemize}
  \item \eng{actually} signifie « en fait » et non « actuellement » (\eng{currently}).
  \item Pas de « s » à \eng{information, advice, news} (indénombrables).
  \item On dit \eng{to listen to music}, \eng{to spend time on something} : la préposition est obligatoire.
  \item \eng{I am agree} est faux : dites \eng{I agree}.
\end{itemize}
\end{bilan}

\begin{aretenir}[Checklist avant le Bac]
\begin{itemize}
  \item[$\square$] Je connais au moins 15 mots du lexique des TIC et des loisirs numériques.
  \item[$\square$] Je sais faire un skimming efficace en moins de deux minutes.
  \item[$\square$] Je sais repérer une date, un chiffre ou un nom par scanning.
  \item[$\square$] Je sais deviner le sens d'un mot inconnu grâce au contexte.
  \item[$\square$] Je maîtrise le gérondif après \eng{enjoy}, \eng{like}, \eng{spend time}.
  \item[$\square$] Je sais employer le present perfect pour parler d'expériences.
  \item[$\square$] Je connais les faux amis : \eng{actually}, \eng{library}, \eng{sensible}.
  \item[$\square$] Je sais répondre à une question de compréhension par une phrase complète.
  \item[$\square$] Je sais résumer un texte en 3 à 5 phrases avec mes propres mots.
  \item[$\square$] Je sais donner mon opinion avec \eng{In my opinion} et la justifier par un exemple.
  \item[$\square$] Je vérifie l'accord sujet-verbe et les prépositions avant de rendre ma copie.
\end{itemize}
\end{aretenir}

\findecours

\end{document}
